\documentclass[11pt,spanish]{article}

% =========================================================
% PLANTILLA BASE PARA INFORMES ACADÉMICOS / TÉCNICOS
% =========================================================

% --- Idioma y codificación ---
\usepackage[spanish,es-tabla]{babel}
\usepackage[T1]{fontenc}
\usepackage{newunicodechar}
\newunicodechar{₡}{\textcolonmonetary}
\newunicodechar{÷}{\(\div\)}
\newunicodechar{×}{\(\times\)}
\newunicodechar{−}{\(-\)}

% --- Márgenes / papel ---
\usepackage[
    left=2.5cm,
    right=2.5cm,
    top=2.5cm,
    bottom=2.5cm,
    headsep=0.5cm,
    footskip=0.8cm
]{geometry}

% --- Matemáticas ---
\usepackage{amsmath}

% Colores para Nivel de Riesgo
\usepackage[table]{xcolor}
\usepackage{colortbl}
\definecolor{C_MuyBajo}{HTML}{47D45A}
\definecolor{C_Bajo}{HTML}{92D050}
\definecolor{C_Medio}{HTML}{FFFF00}
\definecolor{C_Alto}{HTML}{FFC000}
\definecolor{C_MuyAlto}{HTML}{EE0000}
\definecolor{C_Critico}{HTML}{C00000}

% --- Tablas, enlaces y elementos visuales ---
\usepackage{array}
\usepackage{tabularx}
\usepackage{booktabs}
\usepackage{longtable}
\usepackage{graphicx}
\usepackage{float}

\usepackage{tikz}
\usetikzlibrary{positioning,shapes.geometric,arrows.meta,calc}
\usepackage{enumitem}
\usepackage{ragged2e}
\usepackage[hidelinks]{hyperref}
\usepackage{colortbl}
\usepackage{multirow}
\usepackage{hhline}

% Colores principales del documento
\definecolor{tecblue}{HTML}{0B2F5B}
\definecolor{teclightblue}{HTML}{EAF1F8}
\definecolor{tablegray}{HTML}{F4F6F8}
\definecolor{diagramgreen}{HTML}{B2C19C}
\definecolor{diagramlight}{HTML}{E5E9DE}
\definecolor{diagramtext}{HTML}{565A60}
\usepackage{pdfpages}

% Subtítulos menores en bloque, no en línea.
\makeatletter
\renewcommand\paragraph{\@startsection{paragraph}{4}{\z@}%
  {1.25ex \@plus1ex \@minus.2ex}%
  {0.6ex}%
  {\normalfont\normalsize\bfseries}}
\makeatother

% --- Encabezado y pie de página ---
\usepackage{fancyhdr}
\usepackage{lastpage}

% --- Interlineado ---
\linespread{1.2}
\sloppy

% =========================================================
% DATOS DEL DOCUMENTO
% =========================================================
\newcommand{\Institucion}{Instituto Tecnológico de Costa Rica}
\newcommand{\Campus}{Campus Tecnológico Central Cartago}
\newcommand{\DocumentoTitulo}{Plataforma Universitaria Integrada}
\newcommand{\DocumentoSubtitulo}{Modernización y Centralización de Procesos Académicos y Administrativos}
\newcommand{\Curso}{Administración de proyectos}
\newcommand{\Profesor}{Alicia Marcela Salazar Hernández}
\newcommand{\FechaDocumento}{21 de junio de 2026}
\newcommand{\PeriodoAcademico}{I Semestre}
\newcommand{\AnioDocumento}{2026}

\newcommand{\AutoresPortada}{%
    \begin{tabular}{l@{\hspace{1.5cm}}r}
        Mauricio González Prendas & 2024143009\\
        Isaac Jiménez Blanco & 2024202804\\
        Tamara Robles Camacho & 2024099342
    \end{tabular}%
}

% Datos que aparecen en el encabezado.
\newcommand{\DocHeaderLeft}{}
\newcommand{\DocHeaderRight}{\DocumentoTitulo}
\newcommand{\DocFooterRight}{Página \thepage\ de \pageref{LastPage}}

% Metadatos PDF.
\title{\DocumentoTitulo}
\author{\AutoresPortada}
\date{\FechaDocumento}

% =========================================================
% CONFIGURACIÓN DE TABLAS Y UTILIDADES
% =========================================================
\newcolumntype{Y}{>{\RaggedRight\arraybackslash}X}
\newcolumntype{P}[1]{>{\RaggedRight\arraybackslash}p{#1}}
\newcolumntype{L}[1]{>{\RaggedRight\arraybackslash}p{#1}}
\newcolumntype{C}[1]{>{\centering\arraybackslash}p{#1}}
\newcolumntype{R}[1]{>{\raggedleft\arraybackslash}p{#1}}
\renewcommand{\arraystretch}{1.22}
\setlength{\LTpre}{0.5em}
\setlength{\LTpost}{0.5em}

% Imagen segura: si el archivo no existe, muestra un recuadro de marcador.
\newcommand{\SafeImage}[2][0.92\textwidth]{%
    \IfFileExists{#2}{%
        \includegraphics[width=#1]{#2}%
    }{%
        \fbox{%
            \begin{minipage}[c][4cm][c]{#1}
            \centering
            Imagen pendiente\\
            \texttt{#2}
            \end{minipage}%
        }%
    }%
}

% Figura reutilizable.
\newcommand{\EvidenceFigure}[3]{%
    \begin{figure}[H]
    \centering
    \SafeImage{#1}
    \caption{#2}
    \label{#3}
    \end{figure}%
}

% =========================================================
% ENCABEZADO Y PIE
% =========================================================
\pagestyle{fancy}
\fancyhf{}
\fancyhead[R]{\DocumentoTitulo}
\renewcommand{\headrulewidth}{0.4pt}
\fancyfoot[R]{Página \thepage\ de \pageref{LastPage}}
\renewcommand{\footrulewidth}{0.4pt}

% Estilo equivalente para páginas horizontales.
\fancypagestyle{widefancy}{%
    \fancyhf{}%
    \fancyhead[R]{\DocumentoTitulo}%
    \renewcommand{\headrulewidth}{0.4pt}%
    \fancyfoot[R]{Página \thepage\ de \pageref{LastPage}}%
    \renewcommand{\footrulewidth}{0.4pt}%
}

% =========================================================
% PÁGINAS HORIZONTALES PARA TABLAS ANCHAS
% =========================================================
\makeatletter
\newcommand{\SyncPageBuilderDimensions}{%
    \global\vsize=\textheight
    \global\@colroom=\textheight
    \global\@colht=\textheight
}
\makeatother

\newenvironment{widepage}{%
    \clearpage
    \hypersetup{pageanchor=false}%
    \paperwidth=297mm
    \paperheight=210mm
    \pdfpagewidth=297mm
    \pdfpageheight=210mm
    \newgeometry{left=2.5cm,right=2.5cm,top=2.5cm,bottom=2.5cm,headsep=0.5cm,footskip=0.8cm}%
    \setlength{\textwidth}{247mm}%
    \setlength{\textheight}{160mm}%
    \setlength{\linewidth}{\textwidth}%
    \setlength{\columnwidth}{\textwidth}%
    \hsize=\textwidth
    \setlength{\headwidth}{\textwidth}%
    \SyncPageBuilderDimensions
    \pagestyle{widefancy}%
}{%
    \clearpage
    \paperwidth=210mm
    \paperheight=297mm
    \pdfpagewidth=210mm
    \pdfpageheight=297mm
    \restoregeometry
    \SyncPageBuilderDimensions
    \hypersetup{pageanchor=true}%
    \pagestyle{fancy}%
}

\newenvironment{shortwidepage}{%
    \clearpage
    \hypersetup{pageanchor=false}%
    \paperwidth=297mm
    \paperheight=95mm
    \pdfpagewidth=297mm
    \pdfpageheight=95mm
    \newgeometry{left=2.5cm,right=2.5cm,top=2.5cm,bottom=2.5cm,headsep=0.5cm,footskip=0.8cm}%
    \setlength{\textwidth}{247mm}%
    \setlength{\textheight}{45mm}%
    \setlength{\linewidth}{\textwidth}%
    \setlength{\columnwidth}{\textwidth}%
    \hsize=\textwidth
    \setlength{\headwidth}{\textwidth}%
    \SyncPageBuilderDimensions
    \pagestyle{widefancy}%
}{%
    \clearpage
    \paperwidth=210mm
    \paperheight=297mm
    \pdfpagewidth=210mm
    \pdfpageheight=297mm
    \restoregeometry
    \SyncPageBuilderDimensions
    \hypersetup{pageanchor=true}%
    \pagestyle{fancy}%
}

\newenvironment{mediumwidepage}{%
    \clearpage
    \hypersetup{pageanchor=false}%
    \paperwidth=297mm
    \paperheight=145mm
    \pdfpagewidth=297mm
    \pdfpageheight=145mm
    \newgeometry{left=2.5cm,right=2.5cm,top=2.5cm,bottom=2.5cm,headsep=0.5cm,footskip=0.8cm}%
    \setlength{\textwidth}{247mm}%
    \setlength{\textheight}{95mm}%
    \setlength{\linewidth}{\textwidth}%
    \setlength{\columnwidth}{\textwidth}%
    \hsize=\textwidth
    \setlength{\headwidth}{\textwidth}%
    \SyncPageBuilderDimensions
    \pagestyle{widefancy}%
}{%
    \clearpage
    \paperwidth=210mm
    \paperheight=297mm
    \pdfpagewidth=210mm
    \pdfpageheight=297mm
    \restoregeometry
    \SyncPageBuilderDimensions
    \hypersetup{pageanchor=true}%
    \pagestyle{fancy}%
}

% =========================================================
% PORTADA
% =========================================================
\newcommand{\MakeCover}{%
\begin{titlepage}
\thispagestyle{empty}
\centering

{\bfseries \huge \Institucion \par}
\vspace{0.25cm}
{\LARGE \Campus \par}

\vfill

{\huge \DocumentoTitulo \par}
\vspace{0.4cm}
{\Large \DocumentoSubtitulo \par}

\vfill

{\bfseries \LARGE Profesora \par}
{\Large \Profesor \par}

\vfill

{\bfseries \LARGE Estudiantes \par}
{\Large \AutoresPortada \par}

\vfill

{\bfseries\LARGE Fecha \par}
{\Large \FechaDocumento \par}

\vfill

{\LARGE \PeriodoAcademico \par}
{\LARGE \AnioDocumento \par}

\end{titlepage}%
}

% =========================================================
% DOCUMENTO
% =========================================================
\begin{document}
\hypersetup{pageanchor=false}

% =========================
% Portada
% =========================
\MakeCover

% =========================
% Índice
% =========================
\pagenumbering{roman}
\pagestyle{empty}
\addtocontents{toc}{\protect\thispagestyle{empty}}
\tableofcontents
\clearpage
\pagenumbering{arabic}
\hypersetup{pageanchor=true}
\pagestyle{fancy}

% =========================
% Contenido
% =========================

% --- Entregable 01 ---
\section{Caso de negocio}

\subsection{Situación actual}

La Universidad Tecnológica La Mejor es una institución superior que actualmente presenta muchos problemas a causa de cómo gestionan sus procesos académicos, administrativos y financieros, ya que esto se hace mediante una serie de sistemas independientes que no se encuentran integrados entre sí, debido a que esta segmentación tecnológica representa un obstáculo severo para la eficiencia operativa de toda la institución.

Los sistemas actualmente en uso son hojas de cálculo para el registro de calificaciones, junto con correo electrónico para la comunicación constante entre departamentos, y sistemas propietarios sin conexión para la gestión financiera, además de registros físicos o digitales dispersos para el historial académico estudiantil, por lo tanto, la información no está centralizada para los diferentes actores universitarios.

\subsubsection{Diagnóstico operativo}

El análisis de la situación actual muestra los siguientes hallazgos críticos

\begin{itemize}
    \item Duplicidad de datos, debido a que mucha de la información se registra de manera manual en múltiples sistemas, generando así inconsistencias graves y errores operativos constantes.
    \item Falta de visibilidad, ya que las autoridades no disponen de indicadores consolidados en tiempo real para la toma de decisiones estratégicas institucionales.
    \item Procesos lentos, por lo tanto, la matrícula estudiantil, el registro de calificaciones y los tramites financieros dependen de una exhaustiva intervención manual en cada etapa del proceso diario.
    \item Comunicación deficiente, debido a que la información no fluye de forma automática entre el departamento académico, administrativo y financiero.
    \item Insatisfacción estudiantil, esto implica que los estudiantes deben acudir presencialmente, o contactar múltiples oficinas, para realizar sus gestiones básicas universitarias.
\end{itemize}

\subsection{Problema identificado}

El problema central de la Universidad Tecnológica La Mejor se basa en la ausencia de una plataforma integrada que centralice y automatice los procesos académicos, administrativos y financieros de la institución, ya que esta carencia genera ineficiencias operativas medibles, junto con un deterioro notable en la experiencia de todos los usuarios, además de limitaciones significativas para la toma de decisiones basada en datos concretos y reportes financieros exactos.

\subsubsection{Árbol de problemas}

Los efectos identificados como consecuencia directa del problema central son los detallados a continuación

\begin{itemize}
    \item Incremento en los tiempos de atención al estudiante por procesos manuales repetitivos, de esta manera, se reduce fuertemente la eficiencia diaria del personal administrativo.
    \item Alta tasa de errores en registros académicos y financieros, debido a la falta de validación automatizada en todos los sistemas actuales del campus.
    \item Imposibilidad de generar reportes ejecutivos en tiempo real para la toma de decisiones, por lo tanto, los directivos carecen de información oportuna para evaluar la universidad.
    \item Carga administrativa excesiva sobre el personal docente y administrativo, ya que deben duplicar esfuerzos físicos en múltiples plataformas tecnológicas anticuadas.
    \item Pérdida de competitividad institucional frente a universidades con sistemas modernos, esto implica que los estudiantes prefieran explorar otras opciones educativas de la región.
\end{itemize}

\subsubsection{Impacto cuantificable}

Se estima que la situación actual genera los siguientes impactos negativos profundos en la dinámica de la institución

\begin{itemize}
    \item El tiempo promedio de atención por trámite estudiantil se encuentra entre 45 a 90 minutos presenciales, ya que no existen plataformas virtuales adecuadas para los jóvenes.
    \item El porcentaje de errores en registros académicos manuales es aproximadamente del 8 al 12 por ciento, debido a la ausencia total de validaciones automáticas modernas.
    \item Las horas trabajadas de los funcionarios semanales destinadas a tareas administrativas redundantes representan más de 120 horas por departamento, en consecuencia, se desperdicia tiempo valioso del personal capacitado.
    \item El costo de oportunidad estimado por ineficiencias operativas es significativo para la sostenibilidad institucional, por lo tanto, se deben buscar urgentemente alternativas de mejora eficientes.
\end{itemize}

\subsection{Oportunidad de mejora}

La modernización tecnológica de la Universidad Tecnológica La Mejor representa una gran oportunidad para transformar su modelo operativo, debido a que esto permitirá mejorar sustancialmente la experiencia de todos los funcionarios y estudiantes de la institución, además de posicionarla como una universidad de referencia fundamental en gestión educativa dentro de toda la región centroamericana.

\subsubsection{Visión de la solución}

El desarrollo de una plataforma universitaria integrada permitirá consolidar en un único punto de acceso seguro todos los procesos académicos, administrativos y financieros, para lo cual ofrecerá interfaces completamente diferenciadas y adaptadas funcionalmente a estudiantes, docentes, personal administrativo, área financiera y miembros de la alta dirección ejecutiva.

\subsubsection{Oportunidades identificadas}

La automatización de procesos permite reducir de forma considerable el tiempo dedicado a tareas manuales y repetitivas, logrando beneficios de la siguiente forma

\begin{itemize}
    \item Centralización de la información para garantizar una única fuente de verdad aplicable a todos los datos institucionales, ya que esto evita discrepancias molestas y constantes.
    \item Mejora en la experiencia estudiantil al brindar un portal completamente accesible para realizar trámites en línea, por lo tanto, se evitan para siempre las filas presenciales en el campus.
    \item Habilitación de inteligencia institucional para generar tableros ejecutivos con indicadores académicos y financieros en tiempo real, debido a que esto apoya fuertemente la toma de decisiones.
    \item Reducción de errores operativos al implementar validaciones automáticas que eliminen inconsistencias financieras de manera inmediata, de esta manera, se protege la información contable de la entidad.
    \item Escalabilidad al diseñar una arquitectura técnica flexible que permita incorporar nuevos módulos o funcionalidades en el futuro, esto implica que el nuevo sistema no quedará obsoleto a corto plazo.
\end{itemize}

\subsection{Beneficios esperados}

Para este nuevo sistema informático se esperan múltiples beneficios que mejorarán directamente la experiencia de todos los actores de la institución, ya que la novedosa plataforma centralizará la gestión general universitaria.

\subsubsection{Beneficios cuantitativos}

En la siguiente tabla se muestran los beneficios numéricos, comparando la situación actual contra la gran meta esperada tras implementar la solución final

\begin{table}[H]
\centering
\small
\caption{Beneficios esperados}\label{tab:beneficios}
\begin{tabularx}{\textwidth}{p{5.2cm} p{4.1cm} Y}
\toprule
\textbf{Beneficio} & \textbf{Situación actual} & \textbf{Meta esperada} \\
\midrule
Tiempo promedio por trámite estudiantil & 45 a 90 min presenciales & 5 a 10 min en línea \\
Tasa de error en registros académicos & 8\% a 12\% & Menor al 1\% \\
Horas administrativas por semana & Mayor a 120 horas por departamento & Menor a 40 horas por departamento \\
Satisfacción estudiantil según encuestas & No medida o es muy baja & Mayor al 85\% de satisfacción general \\
Disponibilidad de reportes ejecutivos & 72 horas en promedio & Tiempo real inmediato \\
Acceso a trámites fuera del campus & No disponible del todo & 24/7 desde dispositivos inteligentes \\
\bottomrule
\end{tabularx}
\end{table}

\subsubsection{Beneficios cualitativos}

La plataforma también generará grandes ventajas organizacionales de tipo no numérico

\begin{itemize}
    \item Mejora en la imagen institucional y en la percepción de modernidad por parte de estudiantes y docentes, debido a que contarán con un sistema de última generación muy eficiente.
    \item Mayor transparencia en todos los procesos académicos y financieros para todos los actores, por lo tanto, existirá muchísima más confianza institucional hacia la administración.
    \item Fortalecimiento de la cultura organizacional orientada a la eficiencia y la innovación, ya que el personal dejará de realizar trabajos rutinarios excesivos en las oficinas.
    \item Cumplimiento firme de estándares de calidad universitaria y posibilidad real de obtener una acreditación institucional prestigiosa, para lo cual el sistema es una herramienta tecnológica clave.
    \item Retención estudiantil mejorada gracias a una experiencia universitaria mucho más fluida y satisfactoria, esto implica que los estudiantes se mantendrán sumamente motivados en sus respectivas carreras.
    \item Reducción del estrés operativo diario de todo el personal administrativo y del cuerpo docente, en consecuencia, mejorará de forma sostenida el clima laboral general.
\end{itemize}

\subsection{Alternativas evaluadas}

Para atender la problemática identificada, se evaluaron cuidadosamente tres alternativas estratégicas con distintos niveles de inversión, riesgo y beneficio, debido a que la primera alternativa fue mantener el proyecto a como lo están manejando actualmente, mientras que la segunda alternativa consistió en un sistema de software comercial que centraliza la gestión operativa de recursos en una única plataforma genérica, y finalmente, la tercera alternativa correspondió a un desarrollo totalmente a la medida donde se puede adaptar al cien por ciento todo el producto a las necesidades de la institución universitaria. Asimismo, a continuación se presenta el análisis comparativo completo

\begin{widepage}
{\small
\setlength{\tabcolsep}{3pt}

\begin{longtable}{p{4.1cm} p{4.8cm} p{5.8cm} p{7.2cm}}
\caption{Análisis comparativo de alternativas}\label{tab:analisis-alternativas}\\[-0.5em]
\toprule
\textbf{Criterio} & \textbf{Alternativa A Status quo} & \textbf{Alternativa B Sistema comercial} & \textbf{Alternativa C Desarrollo a la medida} \\
\midrule
\endfirsthead
\toprule
\textbf{Criterio} & \textbf{Alternativa A Status quo} & \textbf{Alternativa B Sistema comercial} & \textbf{Alternativa C Desarrollo a la medida} \\
\midrule
\endhead
Costo de implementación & Ninguno & Muy alto de más de 65 millones & Moderado según el alcance definido \\
Tiempo de implementación & Inmediato & 12 a 24 meses de trabajo profundo & 3 meses de trabajo ágil \\
Adaptación a la institución & Alta porque es la configuración actual & Baja por ser un producto genérico y cerrado & Muy alta porque es absolutamente personalizado \\
Escalabilidad técnica & Nula & Media y dependiente de terceros & Alta y autónoma \\
Mantenimiento posterior & Sumamente complejo de ejecutar & Alta dependencia del proveedor externo & Gran autonomía interna institucional \\
Mejora operativa esperada & Ninguna mejora perceptible & Alta capacidad de mejora & Muy alta capacidad de mejora global \\
Riesgo de adopción final & Bajo & Alto & Medio \\
Recomendación ejecutiva & No es viable operativamente a futuro & No es viable por el alto costo financiero asociado & Esta es la excelente opción seleccionada \\
\bottomrule
\end{longtable}
}
\end{widepage}

\subsubsection{Justificación de la alternativa seleccionada}

La alternativa de desarrollo a la medida fue seleccionada indiscutiblemente por ofrecer el mejor balance general entre costo y tiempo de implementación, así como también una gran adaptabilidad a los procesos específicos de la universidad y un fuerte potencial de mejora operativa en todos los departamentos, ya que a diferencia de un software comercial genérico, una solución desarrollada internamente puede ajustarse con gran precisión a todos los flujos de trabajo institucionales, permitiendo evitar licencias de uso sumamente costosas y garantizando de esa forma la autonomía tecnológica a largo plazo para toda la organización educativa.

\subsection{Justificación económica}

La viabilidad económica de este importante proyecto integral se sustenta sólidamente en el análisis profundo de los costos de implementación versus todos los grandes beneficios operativos y estratégicos que generará la plataforma web a mediano y largo plazo, en consecuencia, la prestigiosa universidad recuperará rápidamente su gran inversión económica inicial.

\subsubsection{Estimación de costos del proyecto}

En la siguiente tabla se muestra cada categoría de costo con su respectivo valor financiero estimado y el porcentaje exacto respecto al gran total del costo del desarrollo

\begin{table}[H]
\centering
\small
\caption{Estimación de costos}\label{tab:costos}
\begin{tabularx}{\textwidth}{Y R{4.1cm} R{3.2cm}}
\toprule
\textbf{Categoría de costo} & \textbf{Costo estimado} & \textbf{Porcentaje del total} \\
\midrule
Mano de obra del equipo de proyecto experto & ₡14,000,000 & 51\% \\
Licencias comerciales y herramientas de software & ₡1,700,000 & 6\% \\
Infraestructura tecnológica en la nube y hosting & ₡1,900,000 & 7\% \\
Capacitación general al usuario y gestión del cambio & ₡2,300,500 & 9\% \\
Pruebas operativas exhaustivas y control de calidad & ₡3,000,000 & 11\% \\
Contingencias financieras estimadas del diez por ciento & ₡1,600,000 & 6\% \\
Documentación técnica y gestión completa del proyecto & ₡2,700,000 & 10\% \\
\midrule
\textbf{Total estimado del proyecto informático} & \textbf{₡17,625,200} & \textbf{100\%} \\
\bottomrule
\end{tabularx}
\end{table}

\subsubsection{Retorno sobre la inversión}

Considerando los grandes ahorros operativos estimados que serán directamente derivados de la reducción contundente de horas administrativas y la disminución de errores graves en todos los registros contables, junto con la eliminación definitiva de procesos manuales redundantes, se proyecta un rápido retorno de la alta inversión dentro del mismísimo primer año de operación continua del nuevo sistema informático.

\begin{table}[H]
\centering
\small
\caption{Beneficios financieros}\label{tab:beneficios-financieros}
\begin{tabularx}{\textwidth}{Y R{5cm}}
\toprule
\textbf{Fuente de ahorro o de un gran beneficio financiero} & \textbf{Estimación anual en colones} \\
\midrule
Reducción contundente de horas administrativas redundantes & ₡9,000,000 \\
Eliminación definitiva de errores de cobro y correcciones manuales & ₡4,000,000 \\
Reducción de todos los costos en materiales físicos e impresión de papel & ₡2,500,000 \\
Mejora general en la retención de toda la población estudiantil universitaria & ₡8,000,000 \\
Reducción de los tiempos prolongados de atención presencial al estudiante & ₡3,500,000 \\
\midrule
\textbf{Ahorro global o beneficio anual institucional estimado} & \textbf{₡27,000,000} \\
\bottomrule
\end{tabularx}
\end{table}

\subsubsection{Análisis numérico de costo y beneficio}

En la siguiente tabla se muestra un análisis detallado de la inversión, evaluando la estrecha relación entre el costo inicial y el gran beneficio de implementar este proyecto tecnológico vanguardista en la institución, así como el beneficio anual estimado, junto con el importante periodo temporal de recuperación, además del retorno total de la inversión que ha sido calculado cuidadosamente tanto para el primero como para el tercer año operativo del sistema.

Las fórmulas financieras que fueron estrictamente usadas para sacar dichos resultados matemáticos tan favorables son las presentadas detalladamente a continuación.

La primera fórmula del período de recuperación es la inversión de capital dividida directamente entre el beneficio operativo anual.

La segunda fórmula del retorno sobre la inversión para el primer año es el beneficio financiero del primer año menos la inversión inicial, dividido todo ese resultado entre la misma inversión inicial y finalmente multiplicado por cien.

La tercera fórmula del retorno sobre la inversión para el tercer año es el beneficio acumulado de tres largos años menos la inversión de inicio, dividido todo esto entre la inversión y multiplicado finalmente por cien.

La cuarta fórmula correspondiente a la relación entre costo y beneficio es el beneficio monetario anual dividido entre la inversión del primer año, o también, el beneficio acumulado global dividido entre la misma inversión del tercer periodo anual.

A continuación se presentan de manera sumamente clara todos los resultados correspondientes a las métricas descritas

\begin{table}[H]
\centering
\small
\caption{Resumen de inversión}\label{tab:resumen-inversion}
\begin{tabularx}{\textwidth}{Y p{5.2cm}}
\toprule
\textbf{Inversión total y completa del desarrollo del proyecto web} & ₡17,625,200 \\
\textbf{Beneficio económico anual estimado para la toda la entidad} & ₡27,000,000 \\
\textbf{Período estimado temporal de recuperación de la alta inversión} & Aproximadamente doce meses \\
\textbf{Retorno porcentual sobre la inversión que ha sido proyectado para el año uno} & Un negativo 0.7\% \\
\textbf{Retorno porcentual sobre la inversión que ha sido proyectado para el año tres} & Mayor al 197\% \\
\textbf{Métrica de la relación final obtenida entre el alto costo y el enorme beneficio} & Es de 0.99 para el primer año y \textgreater{}2.97 para el tercer periodo anual \\
\bottomrule
\end{tabularx}
\end{table}

\newpage
% --- Entregable 02 ---
\section{Acta de constitución del proyecto}

\subsection{Definición del problema y propósito}

La Universidad Tecnológica La Mejor presenta problemas importantes en la forma en que gestiona sus procesos académicos, administrativos y financieros, debido a que actualmente la información se encuentra distribuida en varios sistemas independientes, hojas de cálculo, correos electrónicos y registros físicos o digitales completamente separados. Esta situación provoca una constante duplicidad de datos y errores graves en registros, junto con atrasos notables en trámites, por lo tanto, existe muy poca visibilidad para la toma de decisiones institucionales.

El resultado que no se desea mantener es una gestión universitaria fragmentada, donde estudiantes, docentes, personal administrativo y autoridades deben consultar información en diferentes lugares, ya que esta situación hace que los procesos sean mucho más lentos, esto implica que la institución dependa excesivamente de tareas operativas manuales.

Lo que se busca lograr con el proyecto es diseñar y desarrollar un prototipo web navegable de una plataforma universitaria integrada, para lo cual la solución permitirá representar en un solo sistema virtual los principales procesos estudiantiles, docentes, administrativos, financieros y ejecutivos de la universidad. El enfoque del proyecto es puramente la gestión del proyecto y el desarrollo de un entorno visual funcional, debido a que no se contempla el desarrollo de un sistema de procesamiento real ni de una base de datos operativa.

Los procesos que generan actualmente este resultado deficiente son los detallados a continuación

\begin{itemize}
    \item Proceso manual de matrícula de cursos estudiantiles.
    \item Proceso desactualizado de consulta de horarios.
    \item Proceso ineficiente de registro y consulta de calificaciones.
    \item Proceso fragmentado de estado de cuenta y pagos.
    \item Proceso aislado de gestión de estudiantes, docentes y cursos.
    \item Proceso tardío de generación de reportes académicos y financieros.
    \item Proceso de comunicación deficiente entre las áreas académicas, administrativas y financieras.
\end{itemize}

Las medidas que permitirán evidenciar una mejora real son las siguientes

\begin{itemize}
    \item Reducción comprobable del tiempo promedio de atención por trámite estudiantil.
    \item Disminución sistemática de errores en registros académicos y financieros.
    \item Cantidad de módulos representados correctamente en el prototipo entregado.
    \item Cumplimiento estricto de criterios de aceptación definidos para cada módulo.
    \item Nivel de satisfacción esperado por parte de todos los usuarios internos y estudiantes.
    \item Disponibilidad inmediata de indicadores ejecutivos en una pantalla consolidada.
\end{itemize}

El desempeño actual muestra trámites estudiantiles excesivamente lentos, con tiempos estimados entre 45 y 90 minutos de atención presencial, de esta manera, también se identifica una tasa de error aproximada entre el 8 y 12 por ciento en registros manuales, además de más de 120 horas semanales dedicadas a tareas administrativas redundantes por cada departamento universitario.

La meta principal de desempeño es representar una solución tecnológica que reduzca el tiempo de trámite estudiantil a un rango estimado entre 5 y 10 minutos en línea, para lo cual se requiere disminuir los errores esperados a menos del 1 por ciento mediante validaciones estrictas del sistema, y de esa manera, mostrar información ejecutiva de forma sumamente consolidada. Para efectos de este proyecto académico, se espera completar toda la documentación de administración del proyecto y el prototipo visual navegable dentro del periodo planificado de tres meses, en consecuencia, esto implica trabajar con un cronograma controlado, priorizando los módulos mínimos solicitados y evitando agregar funcionalidades que se encuentren fuera del alcance estipulado.

\subsection{Caso de negocio}

El problema actual tiene un impacto directo y perjudicial en los estudiantes, docentes, personal administrativo, área financiera y autoridades académicas de la Universidad Tecnológica La Mejor.

En el caso de los estudiantes, el impacto negativo se refleja en trámites muy lentos y la necesidad obligatoria de acudir a varias oficinas, junto con la dificultad para consultar información académica o financiera desde un único lugar unificado, ya que esto afecta severamente la experiencia estudiantil y genera muchísima insatisfacción general con los servicios institucionales.

En cuanto a los docentes, el impacto operativo se observa en la gran cantidad de procesos manuales requeridos para registrar calificaciones, controlar la asistencia y consultar información de cursos, de esta manera, se aumenta significativamente la carga operativa diaria, por lo tanto, se reduce el valioso tiempo disponible para las actividades propiamente académicas.

Para el personal administrativo y financiero, el problema central genera una constante duplicidad de registros y una continua revisión manual de la información, junto con una enorme dificultad para mantener todos los datos actualizados, asimismo, esto aumenta dramáticamente el riesgo de cometer errores humanos al trabajar con información sumamente dispersa.

En la alta dirección, el impacto se relaciona íntimamente con la muy poca visibilidad que poseen para tomar decisiones estratégicas, debido a que, al no contar con indicadores académicos y financieros centralizados, los reportes finales dependen de consolidaciones manuales muy tediosas, esto implica que pueden tardar muchísimo tiempo en generarse.

Este proyecto tecnológico es prioridad debido a que la universidad necesita modernizar urgentemente sus procesos y centralizar la información para lograr mejorar la eficiencia operativa global, además de que el proyecto permite trabajar una solución moderna que está alineada completamente con proyectos informáticos donde se deben gestionar adecuadamente aspectos como alcance, cronograma, costos, calidad, riesgos, recursos, comunicaciones y cambios.

Los entregables clave esperados para el final del proyecto son los siguientes

\begin{itemize}
    \item Sección fundacional del proyecto.
    \item Definición detallada del alcance.
    \item Estructura de desglose del trabajo del proyecto.
    \item Cronograma estructurado en un gestor de proyectos.
    \item Matriz de asignación de responsabilidades.
    \item Organigrama completo del proyecto.
    \item Plan integral de recursos.
    \item Estimación financiera de costos.
    \item Plan de gestión de calidad.
    \item Plan estructurado de comunicaciones.
    \item Plan de gestión y registro activo de riesgos.
    \item Matriz de evaluación de probabilidad e impacto.
    \item Plan formal de adquisiciones.
    \item Indicadores clave de rendimiento del proyecto.
    \item Prototipo visual navegable.
    \item Proceso estructurado de control de cambios.
    \item Solicitudes formales y evaluación técnica de cambios.
    \item Tablero ejecutivo visual del proyecto.
    \item Informe ejecutivo final.
    \item Acta de cierre formal.
    \item Reflexiones académicas y recomendaciones.
\end{itemize}

Los importantes beneficios indirectos que pueden surgir de la ejecución del proyecto son una mucho mejor imagen institucional y una mayor transparencia en los procesos, junto con una notable reducción de la carga operativa, además de establecer una base inicial tecnológica para futuras etapas de automatización institucional.

\subsubsection{Objetivos del proyecto}

\begin{enumerate}
    \item Desarrollar un prototipo visual web navegable que represente fielmente los módulos principales de la Universidad Tecnológica La Mejor.
    \item Centralizar visualmente, de manera eficiente, los procesos académicos, administrativos y financieros dentro de una misma plataforma unificada.
    \item Diseñar un portal estudiantil con inicio de sesión y funciones de matrícula, horarios, calificaciones, estado de cuenta e historial académico.
    \item Diseñar un portal docente completo con gestión directa de cursos, registro diario de asistencia y registro de calificaciones.
    \item Diseñar un portal administrativo robusto con gestión general de estudiantes, docentes, cursos y períodos académicos.
    \item Diseñar un portal financiero seguro con pagos en línea de matrícula, pagos de cursos y acceso al historial de pagos.
    \item Diseñar un tablero ejecutivo interactivo con indicadores valiosos de estudiantes activos, ingresos por matrícula e indicadores académicos y financieros.
    \item Aplicar rigurosamente los grupos de procesos metodológicos mediante entregables técnicos de inicio, planificación, ejecución, monitoreo y control, junto con el cierre.
\end{enumerate}

\subsubsection{Criterios de éxito del proyecto}

El proyecto será considerado sumamente exitoso si se cumplen los siguientes criterios de calidad

\begin{itemize}
    \item El prototipo web incluye absolutamente todos los módulos mínimos solicitados en el alcance original del proyecto.
    \item La navegación técnica permite ingresar de forma fluida al sistema principal y al portal estudiantil, docente, administrativo, financiero y ejecutivo.
    \item Las secciones formales del documento mantienen una total coherencia entre el alcance, cronograma, riesgos, calidad, costos y cambios.
    \item El cronograma general se entrega de forma completa, incluyendo tareas, subtareas, dependencias, duraciones, responsables, recursos, costos, hitos y la respectiva ruta crítica.
    \item El prototipo diseñado no presenta errores críticos de navegación funcional al momento de realizar la entrega final.
    \item Las secciones entregables se completan y revisan satisfactoriamente dentro del periodo definido en el cronograma aprobado.
    \item El patrocinador principal y todo el equipo de trabajo validan que el proyecto cumple cabalmente con el propósito definido en esta sección fundacional.
\end{itemize}

\subsection{Participantes clave}

\subsubsection{Patrocinador principal}

El doctor Roberto Mora Fallas es el rector actual de la Universidad Tecnológica La Mejor.

Su gran participación es estrictamente necesaria porque representa directamente a la alta dirección institucional, de esta manera, autoriza el inicio formal y presupuestario del proyecto, además de que valida que el trabajo responda a la necesidad imperiosa de modernizar y centralizar los procesos académicos, administrativos y financieros.

Su importante rol es brindar un patrocinio constante, además de aprobar esta sección fundacional y revisar los avances más relevantes, junto con validar todos los entregables principales, para lo cual también debe apoyar decididamente al equipo cuando existan decisiones complejas que afecten el alcance, tiempo, costo o calidad esperada.

\subsubsection{Líder del proyecto}

La estudiante Tamara Robles Camacho es la directora formal del proyecto.

Su participación activa es absolutamente necesaria porque debe coordinar de manera experta el proyecto y mantener siempre alineado el trabajo de todo el equipo técnico con los objetivos estratégicos definidos desde el inicio, asimismo, es la gran responsable de dar un seguimiento estricto a las actividades, riesgos, comunicaciones y solicitudes de cambios.

Su exigente rol es liderar toda la fase de planificación y organizar reuniones operativas de seguimiento, junto con distribuir las responsabilidades diarias y comunicar los avances logrados al patrocinador, para lo cual debe verificar minuciosamente que las secciones se completen dentro del alcance que ha sido previamente establecido.

\subsubsection{Miembros del equipo}

El estudiante Isaac Jiménez Blanco es analista de negocio y documentador primario.

Su participación técnica es muy necesaria porque el proyecto requiere una documentación sumamente clara y coherente que esté alineada con todas las prácticas vistas en el curso universitario, además de que apoya directamente el análisis continuo de alcance, riesgos, calidad, comunicaciones, adquisiciones y los cambios solicitados.

Su rol designado es elaborar y revisar las distintas secciones, de manera que pueda mantener la coherencia absoluta entre ellas, además de apoyar el registro activo de riesgos y documentar exhaustivamente los procesos de cambio, garantizando que toda la información del proyecto sea trazable.

El estudiante Mauricio González Prendas es desarrollador web y asegurador de calidad primario.

Su participación técnica es requerida porque el proyecto necesita urgentemente un prototipo web navegable y completamente funcional, para lo cual también se necesita validar constantemente la calidad final del producto informático mediante continuas revisiones visuales y pruebas manuales estrictas basadas en los criterios de aceptación previamente definidos.

Su rol operativo es desarrollar la interfaz visual de usuario y construir con gran cuidado las pantallas principales, organizando la navegación del prototipo, de esta manera, debe realizar pruebas básicas de control y corregir todos los defectos visuales que sean encontrados durante la etapa de revisión por pares.

El estudiante Carlos Sainz Arce es desarrollador web secundario y asegurador de calidad secundario.

Su participación técnica es necesaria para apoyar decisivamente el rápido desarrollo de un prototipo web funcional y verdaderamente completo, debido a que el trabajo es sumamente extenso.

Su rol informático es implementar programáticamente los módulos correspondientes al portal financiero y ejecutivo, junto con realizar las pruebas de calidad de la interfaz, colaborando estrechamente con el primer desarrollador.

El estudiante Jorge Abarca Quiros es analista de negocio secundario y documentador secundario.

Su participación administrativa es necesaria para distribuir de forma equitativa toda la pesada carga de trabajo enfocada en la redacción analítica y la revisión profunda de las secciones integradas.

Su rol intelectual es elaborar y revisar meticulosamente la mitad de las secciones, apoyando activamente a la gerencia en la difícil gestión general de riesgos, costos operativos y las comunicaciones institucionales.

\subsubsection{Otros participantes importantes}

La licenciada Patricia Vargas Soto es la directora administrativa.

Su participación operativa es importante porque representa íntegramente todas las necesidades específicas del área administrativa de la universidad, para lo cual su rol fundamental es validar minuciosamente los procesos que estén relacionados con estudiantes, docentes, cursos y todos los períodos académicos vigentes.

El ingeniero Carlos Brenes Mora es el director de tecnología.

Su participación técnica es muy importante porque debe revisar obligatoriamente la viabilidad tecnológica del prototipo elaborado, y de esta manera, asegurar que la solución informática sea fácilmente mantenible y coherente con una arquitectura web institucional.

La profesora Andrea Solís Zamora es la representante de los docentes.

Su participación académica es importante porque aporta fielmente la perspectiva cotidiana de todos los docentes, para lo cual su rol es validar que el nuevo portal educativo permita representar la estructura de los cursos y el control de asistencia, junto con el registro de las calificaciones de forma sumamente clara.

Los estudiantes activos son los usuarios finales del portal.

Su participación diaria es sumamente importante porque son uno de los principales y más numerosos grupos beneficiados con el proyecto, para lo cual su rol consiste en validar empíricamente la facilidad de uso del portal estudiantil y comprobar la claridad de los procesos de matrícula, consulta de horarios, visualización de calificaciones y estado financiero.

El departamento financiero es el área contable.

Su participación económica es importante porque este grupo valida directamente todas las necesidades funcionales que están relacionadas con los pagos de la matrícula y los pagos de cursos, junto con el historial de pagos de los alumnos y los indicadores financieros de la institución.

El proveedor de servicios en la nube es un proveedor externo.

Su participación técnica puede ser necesaria en el momento en que el prototipo se despliega en un ambiente web público, para lo cual su rol es brindar disponibilidad ininterrumpida del servicio y ofrecer un soporte básico de la infraestructura utilizada.

\subsection{Alcance}

\subsubsection{Elementos dentro del alcance}

El proyecto cubrirá completamente el diseño visual y el desarrollo técnico de un prototipo web navegable para la Universidad Tecnológica La Mejor, debido a que el alcance fijado incluye los procesos principales del portal estudiantil, docente, administrativo y financiero, junto con un tablero ejecutivo.

El alcance detallado incluye los siguientes aspectos

\begin{itemize}
    \item Inicio de sesión seguro.
    \item Tablero principal informativo.
    \item Portal estudiantil interactivo.
    \item Perfil del estudiante actualizado.
    \item Matrícula de cursos en línea.
    \item Consulta dinámica de horarios.
    \item Consulta de calificaciones finales.
    \item Estado de cuenta detallado.
    \item Historial académico completo.
    \item Portal docente funcional.
    \item Gestión integral de cursos.
    \item Registro diario de asistencia.
    \item Registro seguro de calificaciones.
    \item Portal administrativo centralizado.
    \item Gestión completa de estudiantes.
    \item Gestión central de docentes.
    \item Gestión estructurada de cursos.
    \item Gestión de períodos académicos.
    \item Portal financiero automatizado.
    \item Pagos en línea de matrícula.
    \item Pagos en línea de cursos.
    \item Historial de pagos realizados.
    \item Tablero ejecutivo institucional.
    \item Cantidad total de estudiantes activos.
    \item Ingresos monetarios por matrícula.
    \item Indicadores académicos globales.
    \item Indicadores financieros consolidados.
    \item Documentación de gestión del proyecto estructurada según lo solicitado.
    \item Uso exclusivo de control de versiones web únicamente para el resguardo del código del prototipo.
\end{itemize}

También se incluye obligatoriamente la elaboración detallada de elementos de gestión bajo el enfoque del curso de administración de proyectos, ya que estos deben cubrir metódicamente los grupos de inicio, planificación, ejecución, monitoreo y control, junto con el respectivo cierre operativo.

\subsubsection{Elementos fuera del alcance}

Queda estrictamente fuera del alcance el desarrollo técnico de un sistema informático funcional profundo y la construcción lógica de una base de datos real, en consecuencia, el proyecto tampoco incluye ningún tipo de integración con sistemas reales de la universidad, ni tampoco autenticación institucional verdadera, ni mucho menos una pasarela de pagos real o la migración manual de datos históricos.

También queda completamente fuera del alcance el desarrollo informático de una aplicación móvil nativa y la automatización completa de los procesos institucionales, junto con el despliegue técnico en un ambiente productivo oficial y la operación real del sistema informático después de realizarse el cierre definitivo del proyecto académico.

No se incluye en absoluto la sustitución total de los anticuados sistemas institucionales existentes ni la compra de licencias comerciales de software, debido a que el prototipo funcionará exclusivamente como una representación visual y navegable de la plataforma que ha sido propuesta, usando datos completamente simulados cuando sea estrictamente necesario.

\subsection{Hitos}

Asumiendo que el proyecto inicia formalmente el 8 de junio de 2026 y se desarrolla durante un periodo planificado de aproximadamente tres meses, los hitos principales que deben cumplirse son los mencionados a continuación

\begin{enumerate}
    \item \textbf{Aprobación fundacional y confirmación del alcance general}

    La duración estimada es de un solo día, específicamente el 8 de junio de 2026.

    La fecha estimada de cierre es el 8 de junio de 2026.

    \item \textbf{Finalización de la planificación base}

    La duración estimada abarca las primeras semanas del proyecto.

    La fecha estimada de cierre definitivo se define según el cronograma aprobado en Microsoft Project.

    \item \textbf{Elaboración de alcance, estructura de trabajo, organigrama, recursos y costos}

    La duración estimada comprende la fase de planificación inicial del proyecto.

    La fecha estimada de cierre máximo se define según el avance de la línea base del cronograma.

    \item \textbf{Elaboración del cronograma, plan de calidad, comunicaciones, riesgos y adquisiciones}

    La duración estimada de esta fase corresponde al periodo de planificación detallada.

    La fecha estimada de cierre formal se mantiene alineada con el cronograma aprobado.

    \item \textbf{Desarrollo técnico de los módulos principales del prototipo web}

    La duración estimada para el desarrollo abarca las semanas centrales del proyecto.

    La fecha estimada de cierre técnico corresponde al hito de prototipo de la interfaz completado.

    \item \textbf{Elaboración de indicadores, control de cambios y evaluación de solicitudes}

    La duración estimada para estos elementos forma parte del monitoreo y control del proyecto.

    La fecha estimada de cierre depende de las solicitudes de cambio evaluadas.

    \item \textbf{Tablero ejecutivo, informe de cierre, acta general y reflexiones}

    La duración estimada final corresponde a la fase de cierre del proyecto.

    La fecha estimada de entrega preliminar se establece antes del hito final del cronograma.

    \item \textbf{Revisión exhaustiva final, correcciones, integración y entrega final}

    La duración estimada para revisiones se concentra al final del proyecto.

    La fecha estimada de cierre absoluto corresponde al hito de entrega final definido en Microsoft Project.
\end{enumerate}

\subsection{Facilitadores y mitigación de riesgos}

Para que el proyecto avance correctamente y sin atrasos, se deben tener condiciones operativas mínimas que permitan controlar estrictamente el trabajo diario, para lo cual es necesario reducir los riesgos y mantener perfectamente alineado a todo el equipo de trabajo.

\subsubsection{Gobernanza clara del proyecto}

Debe existir una estructura sumamente clara de roles con un patrocinador visible, una directora de proyecto empoderada, junto con miembros del equipo comprometidos y participantes clave informados, debido a que esto permite saber exactamente quién aprueba, quién ejecuta, quién valida y quién debe ser informado en cada etapa.

\subsubsection{Alcance definido desde el inicio}

El alcance del trabajo debe indicar inequívocamente cuáles módulos se construirán y qué elementos quedan totalmente fuera del proyecto, ya que esto reduce dramáticamente el riesgo de un incremento descontrolado del alcance y evita por completo que se agreguen funciones que no fueron planificadas.

\subsubsection{Uso de un repositorio seguro para el prototipo}

Se debe utilizar un repositorio de versiones en línea únicamente para resguardar el código fuente del desarrollo web, por lo tanto, esto ayuda enormemente a mantener un estricto control de versiones y permite organizar los avances, además de evitar cualquier pérdida catastrófica del trabajo técnico.

\subsubsection{Comunicación constante del equipo de trabajo}

El equipo asignado debe mantener reuniones y revisiones muy frecuentes para verificar diariamente el avance, las dudas operativas y los cambios requeridos, esto implica que se logran evitar inconsistencias y la terrible duplicidad de tareas, así como los atrasos en las entregas críticas.

\subsubsection{Criterios de aceptación por cada módulo}

Cada módulo informático debe tener obligatoriamente criterios básicos de aceptación funcional, por ejemplo, debe existir una navegación muy clara y la información debe ser totalmente visible, de esta manera, las pantallas resultan coherentes y garantizan el cumplimiento de todas las funciones mínimas solicitadas por el cliente.

\subsubsection{Plan riguroso de calidad}

Debe existir una fase estructurada de revisión exhaustiva del producto tecnológico y de todos los elementos administrativos, ya que esto incluye invariablemente la realización de pruebas manuales y una revisión visual profunda mediante una lista de verificación de navegación, además de validar el sistema directamente contra el alcance aprobado.

\subsubsection{Control formal y documentado de cambios}

Cualquier cambio operativo que afecte el alcance, cronograma, costo, calidad o los riesgos debe documentarse formalmente y sin excepciones, por lo tanto, esto permite analizar cuidadosamente todo el impacto negativo antes de que se proceda a aprobar las nuevas solicitudes del cliente.

\subsubsection{Seguimiento estricto del cronograma}

El cronograma del proyecto debe elaborarse detalladamente en un software especializado y debe incluir exhaustivamente todas las tareas, subtareas, dependencias, duraciones, recursos, hitos, ruta crítica, responsables y costos financieros, debido a que esto permite controlar rigurosamente el avance diario y ayuda a detectar cualquier atraso incipiente.

\subsubsection{Gestión proactiva de riesgos}

Se debe mantener de manera obligatoria un registro vivo de riesgos que indique con precisión la causa, probabilidad, impacto, nivel, responsable y la respectiva estrategia de respuesta planificada, para lo cual esto permite al equipo anticipar los problemas graves mucho antes de que lleguen a afectar la fecha de entrega.

\subsubsection{Reserva económica de contingencia}

El presupuesto estimado del proyecto debe mantener celosamente una reserva financiera para mitigar posibles imprevistos, ya que esta reserva permite responder ágilmente a los necesarios ajustes de calidad o retrasos y las múltiples necesidades técnicas que no fueron previstas desde el puro inicio del trabajo.

\subsubsection{Seguimiento frecuente por la duración planificada del proyecto}

Debido a que el proyecto cuenta con una duración planificada de aproximadamente tres meses, todo el equipo debe realizar reuniones y revisiones frecuentes sobre el avance real, de esta manera, esto permite detectar los atrasos de forma rápida y lograr corregir incongruencias, además de priorizar firmemente las entregas obligatorias antes de que expire la fecha final establecida en el cronograma.

\subsection{Barreras y riesgos}

Las principales barreras operativas y riesgos sistémicos del proyecto se relacionan estrechamente con el aumento del alcance y el cumplimiento del cronograma planificado, junto con la coordinación interna del equipo y la calidad del prototipo final, además de la coherencia documental requerida.

\subsubsection{Alcance progresivo y creciente}

Existe un riesgo latente de que se quieran agregar funciones visuales adicionales que definitivamente no estaban contempladas al inicio del trabajo, por lo tanto, esto puede afectar gravemente el cronograma y aumentar desproporcionadamente el esfuerzo requerido, reduciendo así la calidad global.

\subsubsection{Inconsistencias y desactualización}

Como las numerosas secciones están divididas equitativamente entre los diferentes integrantes del grupo, puede existir una preocupante inconsistencia general en los nombres, el presupuesto, los roles operativos, las fechas estipuladas o el mismo alcance, esto implica que puede verse afectada seriamente la calidad general del proyecto final al momento de la entrega.

\subsubsection{Subestimación del esfuerzo técnico requerido}

El prototipo web incluye obligatoriamente la programación de varios módulos robustos y muchísimas pantallas interactivas, debido a que si se subestima groseramente el esfuerzo técnico necesario, el desarrollo puede atrasarse irremediablemente o incluso quedar sumamente incompleto.

\subsubsection{Falta de criterios claros de calidad}

Si el equipo de trabajo no logra definir criterios de aceptación que sean meridianamente claros, el grupo podría considerar equivocadamente que está listo un módulo que en realidad todavía no cumple debidamente con la navegación, diseño gráfico o la funcionalidad mínima requerida por los usuarios.

\subsubsection{Fuerte dependencia de herramientas informáticas}

Los problemas sistémicos con el repositorio de código o el software de gestión de proyectos y las distintas herramientas de diseño o del ambiente de desarrollo, pueden generar enormes retrasos operativos o incluso una irrecuperable pérdida de avance si no se manejan con cuidado extremo.

\subsubsection{Disponibilidad muy limitada de los integrantes}

El equipo técnico está conformado por muy pocos miembros de trabajo por lo que, invariablemente, cada persona tiene bajo su responsabilidad varias secciones y áreas de programación, ya que esto puede generar fácilmente una severa sobrecarga laboral y afectar negativamente la revisión final del trabajo.

\subsubsection{Cambios solicitados en etapas tardías}

Si el cliente se dispone a solicitar cambios sustanciales cuando la fecha de entrega está muy cerca, puede afectarse drásticamente el cumplimiento del cronograma y la calidad del código fuente, por esa razón, se debe usar invariablemente el proceso formal de control de cambios.

\subsubsection{Riesgo constante de baja usabilidad}

El prototipo de software puede llegar a cumplir con todas las pantallas solicitadas contractualmente, pero aun así no ser fácil de usar ni intuitivo, de esta manera, esto afectaría profundamente la experiencia de uso de todos los estudiantes, docentes y el personal administrativo.

\subsubsection{Riesgo de falta de trazabilidad}

Si las secciones escritas no se conectan de manera lógica y secuencial entre sí, puede ser excesivamente difícil poder demostrar la relación metodológica existente entre el alcance, la estructura de trabajo, el cronograma base, los riesgos, la calidad implementada y el cierre definitivo.

\subsubsection{Restricción temporal del cronograma aprobado}

El proyecto tecnológico debe completarse dentro del plazo definido en el cronograma aprobado, por lo que el equipo ejecutor debe controlar de forma estricta la unificación del documento, el prototipo visual, la revisión de calidad y el cierre formal, ya que cualquier desviación puede afectar la profundidad esperada si el grupo no prioriza correctamente el volumen de trabajo.

\subsubsection{Supuestos del proyecto}

\begin{itemize}
    \item El equipo técnico tendrá acceso irrestricto a todas las herramientas que son absolutamente necesarias para desarrollar el prototipo y elaborar ágilmente la gestión escrita.
    \item El prototipo visual puede usar tranquilamente datos que sean simulados para representar la información.
    \item No se requiere programar ningún tipo de servicio profundo ni habilitar una base de datos funcional.
    \item La plataforma de software será netamente web y navegable desde cualquier dispositivo estándar.
    \item Las diversas secciones mantendrán una coherencia milimétrica entre sí a pesar de estar integradas en un solo trabajo final.
    \item El presupuesto base utilizado por la gerencia será exclusivamente el definido en el apartado del caso de negocio original.
    \item El patrocinador ejecutivo y todos los interesados clave estarán disponibles a tiempo para realizar las validaciones generales requeridas.
    \item El cronograma final del proyecto será elaborado íntegramente en la herramienta especializada de gestión de proyectos definida previamente.
    \item El control de versiones en la nube se utilizará sola y exclusivamente para guardar el código fuente del prototipo web.
    \item Los diversos cambios relevantes serán documentados detalladamente y evaluados de forma imparcial antes de poder ser aprobados formalmente.
\end{itemize}

\subsection{Estimaciones de apoyo}

Para lograr desarrollar exitosamente el proyecto se requiere asegurar un apoyo continuo en materia de personal operativo, herramientas tecnológicas, conocimientos técnicos específicos, tiempo hábil de trabajo y un presupuesto holgado, ya que aunque el proyecto no contempla propiamente una solución productiva institucional, sí requiere indefectiblemente una planificación muy adecuada y un excelente prototipo web funcional que represente fielmente a la plataforma universitaria esperada.

\subsubsection{Personal técnico requerido}

El equipo núcleo principal del proyecto estará compuesto exactamente por cinco personas capacitadas, logrando así la distribución lógica de las enormes responsabilidades para garantizar un desarrollo metodológico más robusto y ordenado.

\begin{itemize}
    \item Una gerente de proyecto experimentada y líder del grupo.
    \item Dos analistas de negocio expertos que fungirán como documentadores formales.
    \item Dos desarrolladores de interfaces web que también actuarán como aseguradores de calidad.
    \item Un apoyo gerencial de los distintos interesados externos para lograr la validación correcta de todos los procesos de negocio.
\end{itemize}

\subsubsection{Roles principales definidos}

Tamara Robles Camacho es la directora del proyecto.

La directora es la única responsable oficial de coordinar integralmente el proyecto y de lograr revisar todos los avances operativos, además de poder gestionar asertivamente las delicadas comunicaciones institucionales y asegurar tenazmente el cumplimiento general de todas las metas establecidas.

Isaac Jiménez Blanco es analista de negocio y documentador primario.

El documentador es el máximo responsable de apoyar exhaustivamente toda la extensa documentación técnica, incluyendo el registro del alcance, riesgos, calidad operativa, comunicaciones eficaces, planes de adquisiciones y la gestión metódica de los cambios requeridos por el cliente.

Mauricio González Prendas es desarrollador web y asegurador primario de calidad.

El programador es totalmente responsable de construir desde cero los diferentes módulos funcionales del prototipo web navegable, además de tener que realizar simultáneamente diversas pruebas técnicas orientadas al estricto control de calidad del código.

Carlos Sainz Arce es desarrollador web secundario y asegurador de calidad de apoyo.

El desarrollador de apoyo es directamente responsable de construir todos los módulos complejos correspondientes al portal financiero de la universidad y al importante tablero ejecutivo, para lo cual también debe realizar exigentes pruebas operativas de calidad del entorno visual.

Jorge Abarca Quiros es analista de negocio secundario y documentador auxiliar.

El auxiliar documental es firmemente responsable de poder elaborar de forma consistente los distintos apartados técnicos exigidos, además de apoyar activamente en toda la gestión técnica de evaluación de riesgos, costos del proyecto y las fundamentales comunicaciones del equipo con los interesados.

\subsubsection{Equipamiento y herramientas informáticas requeridas}

\begin{itemize}
    \item Computadoras personales portátiles para todo el equipo de trabajo asignado.
    \item Conexión a internet ininterrumpida y de alta velocidad estable.
    \item Repositorio de código en la nube para resguardar el sistema web.
    \item Herramienta especializada de software para diseñar el cronograma completo.
    \item Entornos y herramientas de desarrollo informático orientadas a web.
    \item Un editor moderno de código fuente para programación avanzada.
    \item Múltiples navegadores web de última generación para pruebas de compatibilidad.
    \item Herramientas ofimáticas estandarizadas de procesamiento de texto y hojas de cálculo numéricas.
    \item Herramientas de comunicación inmediata como correo electrónico corporativo o sistemas de mensajería instantánea grupal.
\end{itemize}

\subsubsection{Especialización técnica requerida}

\begin{itemize}
    \item Experiencia en administración de proyectos estructurada bajo un enfoque metodológico internacional riguroso.
    \item Gran conocimiento en gestión de alcance, tiempo, costo, calidad, mitigación de riesgos y comunicaciones asertivas.
    \item Capacidad de desarrollo y programación de interfaces utilizando directamente tecnologías web modernas e interactivas.
    \item Conocimiento fundamental sobre diseño básico de interfaces de usuario eficientes y técnicas de navegación fluida.
    \item Dominio del control de calidad riguroso del software que va a ser entregado.
    \item Experiencia en gestión ágil y ordenada de cambios al alcance original.
    \item Notoria habilidad para la elaboración de redacción técnica y administrativa que sea clara y profesional.
\end{itemize}

\subsubsection{Presupuesto financiero estimado}

El presupuesto preliminar asignado para todo el proyecto es de 17625200 colones, esto de acuerdo con la estimación de costos actualizada mediante la técnica PERT y la reserva de contingencia definida para el proyecto, debido a que la distribución financiera referencial será la presentada a continuación.

\begin{itemize}
    \item Fase de inicio del proyecto por 1886667 colones.
    \item Fase de planificación del proyecto por 3686667 colones.
    \item Ejecución y desarrollo de la interfaz del prototipo por 7010000 colones.
    \item Monitoreo y control del proyecto por 2271667 colones.
    \item Cierre del proyecto por 588333 colones.
    \item Rubros adicionales de software, hardware y servicios por 666000 colones.
    \item Reserva financiera para contingencias por 1515866 colones.
    \item Total global estimado para la ejecución total ascendiendo a 17625200 colones exactos.
\end{itemize}

Este amplio presupuesto opera como una base referencial de carácter preliminar, ya que cualquier tipo de ajuste financiero que sea importante deberá ser evaluado minuciosamente y revisado siempre mediante un proceso de control de cambios, especialmente si este llega a afectar negativamente el alcance, cronograma, costo, calidad o los riesgos estipulados.

\subsection{Acuerdo del contrato}

\subsubsection{Rol del patrocinador principal}

Yo como patrocinador entiendo perfectamente que mantengo bajo mi responsabilidad formal las siguientes obligaciones con el proyecto.

\begin{itemize}
    \item Brindar completa guía institucional y apoyo directo en la definición de esta sección fundacional trabajando estrechamente junto con la directora de todo el proyecto asignado.
    \item Comprometerse activamente a revisar sin retrasos las etapas principales y críticas del gran proyecto informático encomendado.
    \item Ayudar proactivamente a todo el equipo técnico a lograr superar satisfactoriamente las barreras u obstáculos institucionales que puedan llegar a afectar negativamente el alcance, tiempo esperado, costo operativo o la indispensable calidad final.
    \item Verificar exhaustivamente que los enormes beneficios que han sido definidos desde un principio sean totalmente coherentes e integrales con las complejas necesidades institucionales más urgentes.
    \item Aprobar y ratificar de manera completamente formal tanto el inicio general del importante proyecto tecnológico como sus eventuales cambios operativos principales que requieran análisis y autorización de muy alto nivel administrativo.
\end{itemize}

\subsubsection{Rol de la líder general del proyecto}

Yo como directora de proyecto comprendo con total claridad que tengo bajo mi responsabilidad directa las siguientes tareas primordiales.

\begin{itemize}
    \item Guiar de manera constante y experta a todo el equipo de trabajo durante la complicada planificación estratégica y posterior ejecución meticulosa del proyecto universitario completo.
    \item Organizar las múltiples revisiones internas del propio equipo y brindar así las diversas actualizaciones que sean periódicas sobre el importante avance técnico logrado cada semana operativa.
    \item Coordinar las fundamentales revisiones formales directamente con el patrocinador principal del valioso proyecto en curso para mantener la transparencia administrativa exigida.
    \item Asegurar rigurosamente mediante el control estricto que todas las acciones asignadas al gran proyecto en verdad sean debidamente completadas en tiempo y forma, sin excusas ni excepciones algunas.
    \item Notificar obligatoriamente y sin demoras al patrocinador institucional si de manera imprevista llega a cambiar de forma drástica el alcance, propósito inicial, cronograma general estipulado, gran presupuesto planificado o el mismísimo equipo técnico asignado al trabajo diario.
    \item Gestionar de manera muy eficiente y activa todos los diversos riesgos técnicos identificados, además de las variadas comunicaciones diarias y los inevitables cambios organizacionales surgidos dentro del alcance general de todo este proyecto universitario innovador.
\end{itemize}

\subsubsection{Rol de los miembros asignados del equipo técnico}

Nosotros como miembros oficiales y activos del grupo de desarrollo comprendemos de forma inequívoca que tenemos asignadas las siguientes responsabilidades operativas e ineludibles obligaciones.

\begin{itemize}
    \item Contribuir fuertemente y con alto compromiso técnico o administrativo al gran objetivo general de la implementación exitosa de este enorme proyecto tecnológico académico en desarrollo activo.
    \item Comprometerse de forma sumamente férrea e incansable con lograr consolidar completamente tanto el ambicioso propósito fundacional como cada una de las grandes metas operativas debidamente definidas y documentadas por la gerencia.
    \item Completar en el tiempo estimado y con la máxima excelencia profesional todas las etapas que hayan sido asignadas específicamente a la carga de trabajo personal de cada colaborador técnico.
    \item Aportar un inmenso conocimiento técnico y habilidades sólidas de programación o de documentación, además de la calidad meticulosa requerida que corresponda según la función asignada por el rol formal que ocupa dentro de la propia estructura metodológica estipulada con base en el cronograma integral.
    \item Mantener en todo momento una gran responsabilidad mutua e inquebrantable que genere un entorno solidario y eficiente dentro del equipo de trabajo colaborativo para así impulsar vigorosamente todos los procesos de negocio.
    \item Comunicar obligatoriamente y a tiempo, con mucha claridad y sin ocultamientos, si surgen riesgos inminentes, bloqueos paralizantes graves o las diversas y justas preocupaciones personales que pudiesen amenazar gravemente el avance firme del importante proyecto informático y de esta manera perjudicar las metas críticas planificadas a priori.
    \item Cumplir estricta y fielmente con cada meta temporal impuesta por el riguroso cronograma, además de honrar plenamente con las muchas responsabilidades metodológicas previamente definidas y documentadas en el marco constitutivo fundacional de todo el gran proyecto institucional actual que se encuentra en una franca etapa temprana de desarrollo.
\end{itemize}

\newpage
% --- Entregable 03 ---
\section{Registro de partes interesadas}

Esta sección corresponde al registro de partes interesadas del proyecto Plataforma Universitaria Integrada, desarrollado para la Universidad Tecnológica La Mejor, debido a que esta información forma parte del grupo de procesos de inicio según el marco metodológico del PMI, por lo tanto, su propósito es identificar, clasificar y definir a las personas que formarán parte de este proyecto y que pueden afectar o verse afectadas por el desarrollo y los resultados del mismo.

La identificación temprana de las partes interesadas es fundamental para garantizar una buena comunicación, gestionar expectativas y asegurar el apoyo durante el ciclo de vida del proyecto, en consecuencia, estas partes interesadas se dividirán en dos grupos, los internos (miembros del equipo de desarrollo y administración) y los externos (futuros actores de la plataforma), de esta manera, para cada uno se documenta su cargo, nivel de interés, influencia y poder, así como la estrategia más adecuada para mantener su participación acorde con los objetivos del proyecto.

La escala de referencia para comprender mejor la tabla indica que la influencia y el poder se valoran del 1 al 5 (donde 1 es muy bajo y 5 es muy alto).

\subsection{Partes interesadas internas del proyecto}

En la siguiente tabla se muestran las partes interesadas internas del proyecto.

\begin{widepage}
{\small
\setlength{\tabcolsep}{3pt}

\begin{longtable}{p{3.8cm} p{3.8cm} p{5.8cm} p{1.9cm} p{1.4cm} p{5.8cm}}
\caption{Partes interesadas internas}\label{tab:interesados-internos}\\[-0.5em]
\toprule
\textbf{Nombre} & \textbf{Cargo} & \textbf{Interés} & \textbf{Influencia} & \textbf{Poder} & \textbf{Estrategia de involucramiento} \\
\midrule
\endfirsthead
\toprule
\textbf{Nombre} & \textbf{Cargo} & \textbf{Interés} & \textbf{Influencia} & \textbf{Poder} & \textbf{Estrategia de involucramiento} \\
\midrule
\endhead
Tamara Robles Camacho & Gerente del proyecto & Garantizar que el proyecto se ejecute dentro del alcance, tiempo y costo definidos & 5 & 5 & Liderar reuniones semanales de seguimiento, tomar decisiones sobre cambios y gestionar comunicaciones con el patrocinador \\
Mauricio González Prendas & Desarrollador de la interfaz 1 y calidad & Desarrollar una interfaz funcional, navegable y de calidad que cumpla los criterios de aceptación & 5 & 4 & Participar en la definición del alcance técnico, revisiones de código y pruebas de calidad antes de cada entrega \\
Carlos Sainz Arce & Desarrollador de la interfaz 2 y calidad & Desarrollar una interfaz funcional, navegable y de calidad que cumpla los criterios de aceptación & 5 & 4 & Participar en la definición del alcance técnico, revisiones de código y pruebas de calidad antes de cada entrega junto con el desarrollador 1 \\
Isaac Jiménez Blanco & Analista de negocio y documentador 1 & Asegurar que todos los entregables documentales estén completos, coherentes y alineados con los estándares del PMI & 5 & 4 & Responsable de redactar y revisar entregables, mantener actualizado el registro de riesgos y la documentación del proyecto \\
Jorge Abarca Quiros & Analista de negocio y documentador 2 & Asegurar, junto con el analista y documentador 1, que todos los entregables documentales estén completos, coherentes y alineados con los estándares del PMI & 5 & 4 & Redacta y analiza los entregables, manteniendo actualizada la documentación \\
\bottomrule
\end{longtable}
}
\end{widepage}

\subsection{Partes interesadas externas del proyecto}

En la siguiente tabla se mencionan las partes externas interesadas para este proyecto.

\begin{widepage}
{\small
\setlength{\tabcolsep}{3pt}

\begin{longtable}{p{3.8cm} p{4.3cm} p{5.8cm} p{1.9cm} p{1.4cm} p{5.4cm}}
\caption{Partes interesadas externas}\label{tab:interesados-externos}\\[-0.5em]
\toprule
\textbf{Nombre} & \textbf{Cargo} & \textbf{Interés} & \textbf{Influencia} & \textbf{Poder} & \textbf{Estrategia de involucramiento} \\
\midrule
\endfirsthead
\toprule
\textbf{Nombre} & \textbf{Cargo} & \textbf{Interés} & \textbf{Influencia} & \textbf{Poder} & \textbf{Estrategia de involucramiento} \\
\midrule
\endhead
Dr. Roberto Mora Fallas & Rector y patrocinador del proyecto & Modernizar la institución, mejorar la imagen universitaria y optimizar los procesos internos & 4 & 5 & Presentarle la sección de constitución del proyecto para aprobación, informes ejecutivos mensuales y acta de cierre \\
Lic. Patricia Vargas Soto & Directora administrativa & Reducir la carga operativa del personal y eliminar duplicidad de datos entre departamentos & 4 & 4 & Reuniones de validación de requerimientos administrativos y revisión del módulo de gestión \\
Ing. Carlos Brenes Mora & Director de tecnologías de la información & Garantizar que la solución sea técnicamente viable, segura y mantenible a futuro & 3 & 4 & Consultas técnicas en la fase de planificación y revisión de la arquitectura de la interfaz \\
Prof. Andrea Solís Zamora & Representante docente & Contar con herramientas digitales que simplifiquen el registro de calificaciones y asistencia & 4 & 3 & Participar en sesiones de validación del portal docente y pruebas de usuario \\
Estudiantes (colectivo) & Usuarios finales del portal estudiantil & Acceder a matrícula, calificaciones, horarios y pagos de forma rápida y en línea & 5 & 2 & Encuestas de satisfacción y pruebas de usabilidad del portal estudiantil \\
Departamento financiero & Área financiera & Digitalizar los procesos de cobro de matrícula y generación de reportes financieros & 4 & 3 & Validación de requerimientos del módulo financiero y revisión de casos de uso así como las pruebas de usuario para ese módulo \\
Proveedor de alojamiento web & Proveedor externo & Ofrecer infraestructura confiable para el despliegue del sistema & 2 & 2 & Contrato de servicio con nivel de servicio definido, revisión periódica de disponibilidad \\
\bottomrule
\end{longtable}
}
\end{widepage}

\newpage
% --- Entregable 04 ---
\section{Definición del alcance}

Esta sección corresponde a la definición del alcance del proyecto Plataforma Universitaria Integrada desarrollado para la Universidad Tecnológica La Mejor debido a que este entregable forma parte de la fase de planificación y tiene como finalidad establecer de manera clara qué trabajo será incluido y cuáles condiciones se asumen como verdaderas además de las limitaciones que deberán respetarse durante el desarrollo. La definición del alcance es importante porque permite evitar confusiones entre el equipo de trabajo y los interesados del proyecto ya que ayuda a controlar el crecimiento no planificado conocido como corrupción del alcance al dejar por escrito los límites desde etapas tempranas. Este proyecto consiste en diseñar y desarrollar un prototipo de interfaz web navegable para representar los principales procesos académicos y administrativos de la institución por lo tanto la solución propuesta no contempla una programación del lado del servidor funcional ni una base de datos real debido a que el enfoque principal es la gestión de proyectos y la construcción de una interfaz visual. La definición presentada se basa en el acta de constitución y el caso de negocio además del registro de interesados realizados previamente para lo cual se mantiene coherencia entre los objetivos y las necesidades identificadas.

\subsection{Propósito de la definición del alcance}

El propósito de esta sección es establecer los límites formales del proyecto Plataforma Universitaria Integrada para lo cual se describen los elementos que serán desarrollados y los que no serán considerados durante esta etapa además de los supuestos bajo los cuales se planifica el trabajo y las restricciones existentes. La definición del alcance permite responder de forma clara a varias preguntas fundamentales.

\begin{itemize}
    \item Qué se va a desarrollar
    \item Qué módulos tendrá el prototipo
    \item Qué documentos forman parte del proyecto
    \item Qué actividades no serán realizadas
    \item Qué condiciones se asumen para poder planificar
    \item Qué límites afectan el tiempo y los recursos
    \item Qué entregables deben completarse antes del cierre
\end{itemize}

Esta delimitación será utilizada como referencia para la estructura de desglose del trabajo y el cronograma además de la estimación de costos y el plan de calidad de esta manera cualquier solicitud que modifique el alcance aquí definido deberá analizarse mediante el procedimiento formal de control de cambios.

\subsection{Alcance incluido}

El alcance incluido corresponde a todo el trabajo que será desarrollado dentro del proyecto debido a que incluye tanto los entregables de administración y gestión como el desarrollo de un prototipo web navegable.

\subsubsection{Alcance general del proyecto}

El proyecto contempla el diseño y desarrollo de una Plataforma Universitaria Integrada en formato de prototipo visual de esta manera la plataforma deberá representar los procesos principales de la universidad mediante pantallas navegables y módulos funcionales. El sistema debe permitir visualizar una estructura centralizada para los siguientes portales.

\begin{itemize}
    \item Portal estudiantil
    \item Portal docente
    \item Portal administrativo
    \item Portal financiero
    \item Panel ejecutivo
\end{itemize}

El proyecto también incluye la elaboración de los documentos requeridos aplicando los grupos de procesos de inicio y planificación además de la ejecución y el cierre en consecuencia esto garantiza el cumplimiento metodológico.

\subsubsection{Módulos incluidos en el prototipo}

En la siguiente tabla se muestra el alcance funcional incluido para el prototipo de la interfaz de usuario.

\begin{table}[H]
\centering
\caption{Alcance funcional incluido para el prototipo visual}
\label{tab:modulos-incluidos}
{\small
\begin{tabularx}{\textwidth}{L{4cm} Y}
\toprule
\textbf{Módulo} & \textbf{Funcionalidades incluidas} \\
\midrule
Inicio de sesión & Pantalla de inicio de sesión para simular el acceso a la plataforma \\
Panel principal & Menú inicial con acceso a los portales principales del sistema \\
Portal estudiantil & Perfil del estudiante, matrícula de cursos, horarios, calificaciones, estado de cuenta e historial académico \\
Portal docente & Gestión de cursos, registro de asistencia y registro de calificaciones \\
Portal administrativo & Gestión de estudiantes, gestión de docentes, gestión de cursos y períodos académicos \\
Portal financiero & Pagos de matrícula, pagos de cursos e historial de pagos \\
Panel ejecutivo & Cantidad de estudiantes activos, ingresos por matrícula, indicadores académicos e indicadores financieros \\
\bottomrule
\end{tabularx}
}
\end{table}

\subsubsection{Alcance del portal estudiantil}

El portal estudiantil será diseñado para representar las funciones principales que necesita un estudiante dentro de la plataforma ya que esto debe permitir que el usuario visualice su información académica y financiera desde un único lugar. Se incluyen las siguientes características.

\begin{itemize}
    \item Inicio de sesión simulado
    \item Visualización del perfil del estudiante
    \item Matrícula de cursos disponibles
    \item Consulta de horarios
    \item Consulta de calificaciones
    \item Consulta del estado de cuenta
    \item Historial académico del estudiante
\end{itemize}

Este módulo busca representar una experiencia más centralizada para el estudiante debido a que esto reduce la dependencia de varios sistemas o de los diferentes trámites presenciales.

\subsubsection{Alcance del portal docente}

El portal docente será diseñado para representar las funciones principales que necesita un profesor dentro de la universidad debido a que su objetivo es facilitar la gestión académica de los cursos asignados. Se contemplan las siguientes funciones.

\begin{itemize}
    \item Visualización de cursos asignados
    \item Gestión básica de cursos
    \item Registro de asistencia
    \item Registro de calificaciones
    \item Consulta de estudiantes por curso
\end{itemize}

Este portal no realizará cálculos reales de promedios ni guardará información en una base de datos estructurada por lo tanto la información mostrada podrá ser simulada para fines de prototipo.

\subsubsection{Alcance del portal administrativo}

El portal administrativo será diseñado para representar la gestión interna de la institución ya que este módulo se relaciona con los procesos que actualmente se encuentran dispersos o se realizan de forma manual. Entre los elementos que abarca se encuentran los siguientes.

\begin{itemize}
    \item Gestión de estudiantes
    \item Gestión de docentes
    \item Gestión de cursos
    \item Gestión de períodos académicos
    \item Visualización de información administrativa básica
\end{itemize}

Este portal funcionará como una representación visual de los procesos principales asimismo esto permite operar sin automatizar completamente los flujos reales de la universidad.

\subsubsection{Alcance del portal financiero}

El portal financiero será diseñado para representar los procesos relacionados con pagos y consultas de los estudiantes debido a que su finalidad es mostrar cómo la universidad podría centralizar la información dentro de la plataforma. Las funciones consideradas son las siguientes.

\begin{itemize}
    \item Consulta de pagos de matrícula
    \item Consulta de pagos de cursos
    \item Historial de pagos
    \item Visualización de estado financiero del estudiante
\end{itemize}

El módulo no procesará transacciones reales ni se conectará con entidades bancarias por lo tanto la información será netamente representativa y podrá ser simulada.

\subsubsection{Alcance del panel ejecutivo}

El panel ejecutivo será diseñado para apoyar la toma de decisiones de la alta dirección en consecuencia este módulo permitirá visualizar indicadores académicos y financieros de manera resumida. Se detallan los siguientes aspectos.

\begin{itemize}
    \item Cantidad de estudiantes activos
    \item Ingresos por matrícula
    \item Indicadores académicos básicos
    \item Indicadores financieros básicos
    \item Representación visual de información consolidada
\end{itemize}

Este panel utilizará datos simulados debido a que no existirá conexión con bases de información reales ya que su finalidad es demostrar cómo podría presentarse la información general.

\subsubsection{Alcance de la documentación del proyecto}

El proyecto también incluye los respectivos documentos de administración que permiten demostrar que el equipo aplicó los procesos fundamentales de planificación y control además de las etapas de cierre. Se incluyen los siguientes elementos.

\begin{itemize}
    \item Caso de negocio
    \item Acta de constitución del proyecto
    \item Registro de interesados
    \item Definición del alcance
    \item Estructura de desglose del trabajo
    \item Cronograma del proyecto
    \item Matriz de asignación de responsabilidades
    \item Organigrama del proyecto
    \item Plan de recursos
    \item Estimación de costos
    \item Plan de calidad
    \item Plan de comunicaciones
    \item Plan de gestión de riesgos
    \item Registro de riesgos
    \item Matriz de probabilidad e impacto
    \item Plan de adquisiciones
    \item Indicadores clave de rendimiento
    \item Desarrollo de la interfaz de usuario
    \item Proceso de control de cambios
    \item Solicitud de cambio primera
    \item Solicitud de cambio segunda
    \item Evaluación de los cambios
    \item Panel ejecutivo del proyecto
    \item Informe ejecutivo
    \item Acta de cierre
    \item Reflexiones y recomendaciones
\end{itemize}

\subsection{Alcance excluido}

El alcance excluido corresponde a todo aquello que no será desarrollado en esta etapa del proyecto por lo tanto esta sección es necesaria para evitar expectativas incorrectas sobre el producto final. En la siguiente tabla se presentan los elementos que no forman parte del trabajo.

\begin{longtable}{L{5.2cm} L{9.2cm}}
\caption{Elementos excluidos del alcance}\label{tab:alcance-excluido}\\[-0.5em]
\toprule
\textbf{Elemento excluido} & \textbf{Justificación} \\
\midrule
\endfirsthead
\toprule
\textbf{Elemento excluido} & \textbf{Justificación} \\
\midrule
\endhead
Programación de servidor & El proyecto solicita enfocarse en la interfaz navegable \\
Base de datos real & Los datos podrán ser simulados dentro del prototipo \\
Integración con sistemas reales & No se cuenta con acceso a plataformas institucionales verdaderas \\
Autenticación real & El inicio de sesión será representativo y no validará usuarios \\
Pasarela de pagos real & El módulo financiero no procesará transacciones financieras \\
Migración de datos históricos & No se trabajará con información real de estudiantes o docentes \\
Aplicación móvil nativa & El alcance solicitado corresponde a una plataforma web \\
Despliegue en producción & El prototipo será académico y no operará de forma oficial \\
Automatización de procesos & Se representarán flujos principales en lugar de procesos completos \\
Compra de un sistema comercial & El caso de negocio seleccionó una solución desarrollada a la medida \\
Soporte operativo posterior & El proyecto finaliza con la entrega académica y el cierre \\
Capacitación formal a usuarios & Se puede proponer capacitación pero no una ejecución real \\
\bottomrule
\end{longtable}

También queda fuera del alcance cualquier funcionalidad adicional que no esté incluida en los módulos mínimos definidos ya que no se contemplan notificaciones automáticas o analítica avanzada. Si alguno de estos elementos se desea incluir posteriormente se deberá gestionar como una solicitud de cambio y evaluarse su impacto correspondiente en consecuencia.

\subsection{Supuestos}

Los supuestos son condiciones que se consideran verdaderas para poder planificar el proyecto esto implica que si alguno de ellos cambia se puede afectar la ejecución o requerir ajustes en el plan general.

\begin{longtable}{L{2cm} L{12.4cm}}
\caption{Supuestos del proyecto}\label{tab:supuestos}\\[-0.5em]
\toprule
\textbf{Código} & \textbf{Supuesto} \\
\midrule
\endfirsthead
\toprule
\textbf{Código} & \textbf{Supuesto} \\
\midrule
\endhead
S1 & El proyecto será desarrollado como un prototipo académico web navegable \\
S2 & No será necesario desarrollar programas de servidor ni bases de datos funcionales \\
S3 & La información mostrada en el prototipo podrá ser simulada \\
S4 & El equipo contará con acceso a computadoras e internet además de herramientas de desarrollo \\
S5 & El repositorio se utilizará para administrar el código del prototipo, documentos finales y evidencias del proyecto \\
S6 & Los documentos serán elaborados de forma separada por los integrantes del equipo \\
S7 & Los documentos deberán mantener coherencia en presupuestos y alcances generales \\
S8 & El presupuesto base será el mismo definido en el caso de negocio \\
S9 & El proyecto deberá completarse dentro del periodo planificado de tres meses según el cronograma aprobado \\
S10 & El cronograma será elaborado utilizando una herramienta especializada de gestión \\
S11 & Los interesados estarán disponibles para validaciones generales del trabajo \\
S12 & Los módulos definidos serán suficientes para cumplir con la consigna final \\
S13 & Las revisiones de calidad se realizarán mediante pruebas manuales del prototipo \\
S14 & Los cambios importantes serán evaluados mediante el proceso formal de control respectivo \\
\bottomrule
\end{longtable}

\subsection{Restricciones}

Las restricciones son límites que condicionan la ejecución del proyecto debido a que estas normativas deben considerarse obligatoriamente durante toda la planificación y el desarrollo respectivo.

\begin{longtable}{L{2cm} L{12.4cm}}
\caption{Restricciones del proyecto}\label{tab:restricciones}\\[-0.5em]
\toprule
\textbf{Código} & \textbf{Restricción} \\
\midrule
\endfirsthead
\toprule
\textbf{Código} & \textbf{Restricción} \\
\midrule
\endhead
R1 & El proyecto debe ejecutarse y cerrarse según la fecha final definida en el cronograma aprobado de Microsoft Project \\
R2 & El prototipo debe ser web y completamente navegable por los usuarios \\
R3 & No se permite presentar únicamente un prototipo visual estático en herramientas de diseño gráfico \\
R4 & El sistema debe desarrollarse usando tecnologías web permitidas por las normativas académicas \\
R5 & No se desarrollará un componente de servidor funcional bajo ninguna circunstancia \\
R6 & No se implementará una base de datos real con información verificable \\
R7 & El equipo está compuesto por cinco integrantes lo que requiere una clara división de labores \\
R8 & Los documentos deben cumplir todos los elementos solicitados por la profesora encargada \\
R9 & El cronograma debe elaborarse utilizando un software de planificación avanzado \\
R10 & El presupuesto preliminar debe mantenerse alineado con lo establecido previamente \\
R11 & Las decisiones relevantes deben ser comunicadas al equipo para evitar incongruencias \\
R12 & Cualquier cambio fundamental debe pasar por el respectivo control de cambios estipulado \\
\bottomrule
\end{longtable}

Estas restricciones hacen que el equipo deba priorizar los elementos obligatorios y evitar agregar funcionalidades no solicitadas debido a que el tiempo disponible debe gestionarse según el cronograma aprobado y se debe controlar estrictamente el avance del trabajo.

\subsection{Entregables}

Los entregables del proyecto se dividen en dos grupos principales ya que el primero corresponde a los documentos de gestión y el segundo incluye el producto técnico desarrollado para la validación funcional.

\subsubsection{Entregables de gestión del proyecto}

\begin{longtable}{L{1.5cm} L{5.8cm} L{7.1cm}}
\caption{Entregables de gestión del proyecto}\label{tab:entregables-gestion}\\[-0.5em]
\toprule
\textbf{Código} & \textbf{Entregable} & \textbf{Descripción} \\
\midrule
\endfirsthead
\toprule
\textbf{Código} & \textbf{Entregable} & \textbf{Descripción} \\
\midrule
\endhead
E1 & Caso de negocio & Documento que justifica la necesidad del proyecto y analiza alternativas además de definir beneficios \\
E2 & Acta de constitución del proyecto & Documento que autoriza formalmente el esfuerzo y define sus elementos generales \\
E3 & Registro de interesados & Documento que identifica participantes y su respectivo nivel de interés e influencia \\
E4 & Definición del alcance & Documento que define los límites y supuestos además de las restricciones correspondientes \\
E5 & Estructura de desglose del trabajo & Desglose detallado de las actividades hasta llegar al nivel de paquetes \\
E6 & Cronograma del proyecto & Archivo con tareas y dependencias además de incluir los recursos e hitos \\
E7 & Matriz de asignación de responsabilidades & Herramienta que define responsables y aprobadores para cada actividad importante \\
E8 & Organigrama del proyecto & Representación gráfica y jerárquica de los roles principales establecidos \\
E9 & Plan de recursos & Archivo que define los componentes humanos y tecnológicos necesarios \\
E10 & Estimación de costos & Documento que calcula los gastos de mano de obra y recursos informáticos \\
E11 & Plan de calidad & Estrategia que define objetivos y criterios de aceptación para el cumplimiento \\
E12 & Plan de comunicaciones & Guía que define la frecuencia y los medios de las comunicaciones internas \\
E13 & Plan de gestión de riesgos & Metodología detallada para identificar y gestionar adecuadamente las amenazas \\
E14 & Registro de riesgos & Inventario de posibles amenazas junto a sus respectivas respuestas y responsables \\
E15 & Matriz de probabilidad e impacto & Representación visual del nivel crítico de los riesgos identificados previamente \\
E16 & Plan de adquisiciones & Estrategia sobre las compras necesarias y los desarrollos que se harán internamente \\
E17 & Indicadores clave de rendimiento & Métricas de seguimiento para evaluar el tiempo y los costos del proyecto \\
E19 & Proceso de control de cambios & Procedimiento formal para gestionar de forma ordenada las modificaciones solicitadas \\
E20 & Solicitud de cambio primera & Simulación documentada de una solicitud de modificación real para la interfaz \\
E21 & Solicitud de cambio segunda & Segunda simulación documentada sobre alteraciones en el alcance del proyecto \\
E22 & Evaluación de los cambios & Análisis del impacto de las modificaciones sobre alcance, tiempo, costo, calidad y riesgos \\
E23 & Panel ejecutivo del proyecto & Vista consolidada de indicadores de desempeño para el monitoreo adecuado \\
E24 & Informe ejecutivo & Resumen descriptivo del proyecto dirigido fundamentalmente a la alta gerencia \\
E25 & Acta de cierre & Archivo formal para dar por finalizado el proyecto de manera oficial \\
E26 & Reflexiones y recomendaciones & Análisis crítico sobre la gestión del equipo y los aprendizajes logrados \\
\bottomrule
\end{longtable}

\subsubsection{Entregable técnico del proyecto}
\vspace{-1.5em}
\begin{table}[H]
\centering
\caption{Entregable técnico del proyecto}
\label{tab:entregable-tecnico}
{\small
\begin{tabularx}{\textwidth}{L{1.5cm} L{4.8cm} Y}
\toprule
\textbf{Código} & \textbf{Entregable} & \textbf{Descripción} \\
\midrule
E18 & Desarrollo de la interfaz & Prototipo web navegable con diferentes portales y paneles funcionales en la plataforma \\
\bottomrule
\end{tabularx}
}
\end{table}

El entregable técnico debe permitir la navegación entre las pantallas principales y mostrar visualmente los procesos definidos debido a que no se espera una solución productiva completa sino un prototipo operativo a nivel de vistas.

\subsection{Criterios generales de aceptación}

Aunque los criterios específicos se detallarán en el plan de calidad y en la estructura de desglose del trabajo se definen los parámetros generales para validar que el alcance fue cumplido cabalmente.

\begin{longtable}{L{4.6cm} L{9.8cm}}
\caption{Criterios generales de aceptación}\label{tab:criterios-aceptacion}\\[-0.5em]
\toprule
\textbf{Área} & \textbf{Criterio de aceptación} \\
\midrule
\endfirsthead
\toprule
\textbf{Área} & \textbf{Criterio de aceptación} \\
\midrule
\endhead
Alcance del producto & El prototipo incluye los módulos mínimos solicitados en la consigna inicial \\
Navegación & El usuario puede desplazarse entre los módulos principales sin problemas críticos \\
Portal estudiantil & Se visualizan el perfil y las matrículas además de las calificaciones de forma clara \\
Portal docente & Se exponen los cursos y la respectiva asistencia además del registro de notas \\
Portal administrativo & Se presentan los listados de estudiantes y docentes junto con los períodos vigentes \\
Portal financiero & Se despliegan los pagos de matrícula y los historiales de cancelación \\
Panel ejecutivo & Se muestran los indicadores académicos y financieros fundamentales \\
Documentación & Los apartados mantienen coherencia con las versiones anteriores aprobadas \\
Calidad & No existen errores graves de navegación al momento de realizar la entrega \\
Tiempo & Las secciones se completan antes de la fecha máxima estipulada \\
\bottomrule
\end{longtable}

Estos criterios permiten validar el alcance sin agregar trabajo innecesario de esta manera si un elemento no está incluido en esta definición no deberá desarrollarse salvo que sea aprobado formalmente.

\newpage
% --- Entregable 05 ---
\section{Estructura de desglose del trabajo}

Esta sección contiene la Estructura de desglose del trabajo del proyecto Plataforma Universitaria Integrada. Esta estructura es una descomposición jerárquica orientada a los entregables que el equipo de proyecto ejecutará para cumplir los objetivos y crear los productos requeridos.

La Estructura de desglose del trabajo permite definir el alcance total del proyecto dividiendo el trabajo en paquetes manejables y fácilmente asignables. Al tratarse de un prototipo de interfaz navegable acompañado de un conjunto de documentos de gestión académica, esta estructura se organiza en grandes fases que coinciden con los grupos de procesos y el desarrollo técnico, finalizando con la integración y el cierre.

\subsection{Diagrama de la estructura}

A continuación, se presenta de manera gráfica la Estructura de desglose del trabajo. Debido al nivel de detalle y la extensión del diagrama, este se expone en formato de página horizontal en la página siguiente, utilizando una exportación vectorial desde la herramienta original de modelado para garantizar que todos los textos sean legibles.

\begin{shortwidepage}
\begin{figure}[H]
\centering
\SafeImage[1\textwidth]{wbs_diagrama.pdf}
\caption{Estructura de desglose del trabajo - Plataforma Universitaria Integrada}
\label{fig:wbs}
\end{figure}
\end{shortwidepage}

\subsection{Diccionario de la estructura (paquetes principales)}

Para complementar el diagrama anterior y satisfacer el requerimiento de documentar criterios de aceptación, a continuación se describen los paquetes de trabajo más importantes del proyecto. Estos criterios garantizan que cada porción del trabajo alcance el estándar de calidad definido.

\begin{table}[H]
\centering
\caption{Diccionario de la Estructura de desglose del trabajo y criterios de aceptación}
\label{tab:wbs-diccionario}
{\small
\begin{tabularx}{\textwidth}{p{3.5cm} Y}
\toprule
\textbf{Paquete de trabajo} & \textbf{Criterios de aceptación} \\
\midrule
Documentación de planificación & Los documentos de esta fase deben estar aprobados, contener la información requerida y utilizar un formato profesional uniforme. \\
\midrule
Portal estudiantil & El módulo debe permitir navegación simulada de matrícula, consulta de notas y estados de cuenta, cumpliendo los principios de diseño UX aprobados. \\
\midrule
Portal docente & El módulo debe mostrar pantallas funcionales visualmente para registrar asistencia y calificaciones, aunque sin persistencia de base de datos. \\
\midrule
Portal financiero & El módulo debe evidenciar el flujo visual del pago de matrícula y cursos con simulaciones interactivas. \\
\midrule
Panel ejecutivo & El módulo debe estar integrado al prototipo y también ser entregado de manera automatizada como un archivo Excel. \\
\midrule
Integración y pruebas & Todos los portales deben estar unidos bajo un menú de navegación claro. Debe evaluarse con usuarios simulados sin generar errores visuales graves. \\
\midrule
Documentación de cierre & Los documentos de esta fase deben estar alineados con el cronograma y presupuesto establecido. \\
\bottomrule
\end{tabularx}
}
\end{table}

\newpage
% --- Entregable 06 ---
\section{Cronograma del proyecto}

[Espacio reservado para la incorporación del cronograma correspondiente a esta sección, el cual es responsabilidad de otros miembros del equipo].

\newpage
% --- Entregable 07 ---
\section{Matriz RACI}

La Matriz RACI (Responsable, Aprobador, Consultado e Informado) es una herramienta muy importante en la gestión de proyectos debido a que permite definir con claridad las responsabilidades de cada miembro del equipo y de los interesados para cada actividad principal, por lo tanto, su correcta aplicación elimina ambigüedades y evita duplicidades de esfuerzo ya que asegura que cada tarea tenga un único responsable de su ejecución y un único responsable final de su resultado.

Asimismo, esta sección define la Matriz RACI del proyecto Plataforma Universitaria Integrada, abarcando todas las actividades principales desde el inicio hasta el cierre de este, por lo tanto, se actúa en concordancia con los grupos de procesos del PMI.

\subsection{Equipo del proyecto}

En la siguiente tabla se muestran los roles del equipo junto con su responsabilidad principal y los entregables asignados, de esta manera se clarifica la distribución del trabajo.

\begin{table}[H]
\centering
\caption{Equipo del proyecto}
\label{tab:equipo-proyecto}
{\small
\setlength{\tabcolsep}{3.5pt}
\begin{tabularx}{\textwidth}{p{3.5cm} p{3.2cm} Y p{4.2cm}}
\toprule
\textbf{Rol} & \textbf{Nombre} & \textbf{Responsabilidad principal} & \textbf{Entregables asignados} \\
\midrule
Director de proyecto & Tamara Robles Camacho & Dirección general, gestión documental y coordinación del equipo & E2, E20, E21, E23, E25, E26 \\
Desarrollador de la interfaz y pruebas de calidad 1 & Mauricio González Prendas & Desarrollo de la interfaz de usuario navegable y pruebas de calidad & E18 (Portal estudiantil, portal docente y portal administrativo) \\
Desarrollador de la interfaz y pruebas de calidad 2 & Carlos Sainz Arce & Desarrollo de la interfaz de usuario navegable y pruebas de calidad & E18 (Portal financiero y portal ejecutivo) \\
Analista de negocio y documentador 1 & Isaac Jiménez Blanco & Elaboración de entregables documentales del PMI & E1, E3 al E9, E19 \\
Analista de negocio y documentador 2 & Jorge Abarca Quiros & Elaboración de entregables documentales del PMI & E10 al E17, E22, E24 \\
Rector o patrocinador & Dr. Roberto Mora Fallas & Aprobación y toma de decisiones estratégicas & Aprobación de entregables clave \\
Directora administrativa & Lic. Patricia Vargas Soto & Validación de los requerimientos administrativos & Revisión de módulos administrativos \\
Director de TI & Ing. Carlos Brenes Mora & Revisión técnica y viabilidad de la solución & Revisión técnica de la interfaz de usuario \\
\bottomrule
\end{tabularx}
}
\end{table}

\subsection{Convención RACI}

La siguiente convención aplica a todas las actividades documentadas en la matriz ya que es indispensable para su lectura.

\begin{table}[H]
\centering
\caption{Convención RACI}
\label{tab:convencion-raci}
{\small
\setlength{\tabcolsep}{4pt}
\begin{tabularx}{\textwidth}{Y Y Y Y}
\toprule
\textbf{R (Responsable)} & \textbf{A (Aprobador)} & \textbf{C (Consultado)} & \textbf{I (Informado)} \\
\midrule
Ejecuta la tarea directamente & Responsable final y toma decisiones & Asesora y valida el resultado & Se le informa del progreso \\
\bottomrule
\end{tabularx}
}
\end{table}

\subsubsection{Notas de aplicación}

\begin{itemize}
    \item Cada actividad debe tener exactamente un responsable final, por lo tanto, esto garantiza el éxito de la tarea.
    \item Puede haber múltiples ejecutores o responsables en una misma actividad ya que algunas tareas requieren esfuerzo conjunto.
    \item Los roles de consultado e informado no ejecutan la tarea, pero participan en el proceso de validación o comunicación, de esta manera se mantienen integrados.
    \item El director de proyecto mantiene el rol de aprobador en todas las actividades de gestión y planificación debido a que es el líder responsable.
\end{itemize}

\subsection{Matriz de responsabilidades}

A continuación se muestra la matriz completa de responsabilidades organizada por fase del proyecto y actividad principal con la asignación RACI para cada miembro del equipo e interesados clave debido a que es el núcleo del proyecto.

\begin{widepage}
{\scriptsize
\setlength{\tabcolsep}{2pt}
\begin{longtable}{p{1.8cm} p{5.0cm} C{2.0cm} C{2.0cm} C{2.0cm} C{2.0cm} C{2.0cm} C{2.0cm} C{2.0cm} C{2.0cm}}
\caption{Matriz de responsabilidades}\label{tab:matriz-responsabilidades}\\[-0.5em]
\toprule
\textbf{Fase} & \textbf{Actividad principal} & \textbf{Dir. de proyecto (Tamara)} & \textbf{Desarrollador 1 (Mauricio)} & \textbf{Desarrollador 2 (Carlos)} & \textbf{Analista 1 (Isaac)} & \textbf{Analista 2 (Jorge)} & \textbf{Patrocinador (Dr. Roberto)} & \textbf{Dir. admin. (Lic. Patricia)} & \textbf{Dir. TI (Ing. Carlos)} \\
\midrule
\endfirsthead
\toprule
\textbf{Fase} & \textbf{Actividad principal} & \textbf{Dir. de proyecto (Tamara)} & \textbf{Desarrollador 1 (Mauricio)} & \textbf{Desarrollador 2 (Carlos)} & \textbf{Analista 1 (Isaac)} & \textbf{Analista 2 (Jorge)} & \textbf{Patrocinador (Dr. Roberto)} & \textbf{Dir. admin. (Lic. Patricia)} & \textbf{Dir. TI (Ing. Carlos)} \\
\midrule
\endhead
Inicio & Elaborar caso de negocio & A & I & I & R & C & C & C & I \\
Inicio & Desarrollar acta de constitución del proyecto & A & I & I & R & C & C & I & I \\
Inicio & Identificar y registrar a los interesados & A & I & I & R & C & C & C & C \\
Planificación & Definir el alcance del proyecto & A & C & C & R & C & I & C & C \\
Planificación & Elaborar la EDT con criterios de aceptación & A & C & C & R & C & I & I & I \\
Planificación & Desarrollar cronograma (MS Project) & A & C & C & R & C & I & I & I \\
Planificación & Definir matriz RACI y organigrama & A & C & C & R & C & I & I & I \\
Planificación & Elaborar plan de recursos & A & C & C & R & C & I & C & I \\
Planificación & Estimar costos del proyecto & A & C & C & R & R & C & C & I \\
Planificación & Elaborar plan de calidad & A & C & C & R & R & I & I & C \\
Planificación & Elaborar plan de comunicaciones & A & C & C & R & R & C & C & I \\
Planificación & Elaborar plan de gestión de riesgos & A & C & C & R & R & I & I & C \\
Planificación & Elaborar registro de riesgos (10 o más) & A & C & C & R & R & I & I & C \\
Planificación & Elaborar matriz probabilidad e impacto & A & C & C & R & R & I & I & C \\
Planificación & Elaborar plan de adquisiciones & A & C & C & R & R & C & C & C \\
Ejecución & Desarrollar inicio de sesión y panel principal & C & C & R & I & I & I & I & A \\
Ejecución & Desarrollar portal estudiantil & C & R & C & I & I & I & C & A \\
Ejecución & Desarrollar portal docente & C & R & C & I & I & I & C & A \\
Ejecución & Desarrollar portal administrativo & C & R & C & I & I & I & A & C \\
Ejecución & Desarrollar portal financiero & C & C & R & I & I & I & C & A \\
Ejecución & Desarrollar panel de control ejecutivo & C & C & R & I & I & A & C & C \\
Ejecución & Realizar pruebas de usabilidad y calidad & A & C & R & C & C & I & C & C \\
Monitoreo & Definir y monitorear los indicadores de rendimiento & A & I & I & C & R & C & I & I \\
Monitoreo & Gestionar control de cambios & A & C & C & R & C & C & I & I \\
Monitoreo & Elaborar solicitudes de cambio & A & C & C & C & R & C & I & I \\
Monitoreo & Evaluar impacto de los cambios & A & C & C & C & R & C & I & C \\
Monitoreo & Dar seguimiento al cronograma & A & C & C & R & R & I & I & I \\
Monitoreo & Dar seguimiento a riesgos activos & A & C & C & R & C & I & I & C \\
Monitoreo & Mantener el panel de control ejecutivo & A & C & C & C & R & C & I & I \\
Cierre & Elaborar informe ejecutivo final & A & I & I & C & R & C & C & I \\
Cierre & Formalizar acta de cierre del proyecto & A & I & I & R & C & A & C & I \\
Cierre & Documentar reflexiones y recomendaciones & A & R & R & R & R & I & I & I \\
\bottomrule
\end{longtable}
}
\end{widepage}

\subsection{Análisis de carga por rol}

En la siguiente tabla se observan cuántas veces aparece cada persona como responsable, aprobador y consultado, de esta manera se evalúa la distribución del trabajo.

\begin{table}[H]
\centering
\caption{Análisis de carga por rol}
\label{tab:analisis-carga-rol}
{\small
\setlength{\tabcolsep}{3.5pt}
\begin{tabularx}{\textwidth}{p{4.2cm} C{1.5cm} C{1.5cm} C{1.5cm} Y}
\toprule
\textbf{Rol} & \textbf{\# como R} & \textbf{\# como A} & \textbf{\# como C} & \textbf{Observación} \\
\midrule
Director de proyecto (Tamara) & 0 & 26 & 6 & Lidera toda la gestión del proyecto \\
Desarrollador 1 (Mauricio) & 4 & 0 & 22 & Foco en la ejecución técnica de la interfaz de usuario \\
Desarrollador 2 (Carlos) & 4 & 0 & 21 & Foco en la ejecución técnica de la interfaz de usuario \\
Analista 1 (Isaac) & 18 & 0 & 6 & Principal ejecutor de la documentación del PMI \\
Analista 2 (Jorge) & 14 & 0 & 11 & Principal ejecutor de la documentación del PMI \\
Rector o patrocinador & 0 & 2 & 9 & Toma decisiones estratégicas y aprueba procesos \\
Directora administrativa & 0 & 1 & 14 & Valida los requerimientos administrativos \\
Director de TI & 0 & 4 & 12 & Revisa la viabilidad técnica \\
\bottomrule
\end{tabularx}
}
\end{table}

\subsection{Observaciones finales}

La presente Matriz RACI fue elaborada en concordancia con los roles definidos en el organigrama del proyecto y es coherente con los paquetes de trabajo identificados en la EDT, por lo tanto, cualquier modificación en el equipo o en el alcance deberá reflejarse en una actualización de este apartado mediante el proceso formal de control de cambios.

Se recomienda revisar esta matriz al inicio de cada fase del proyecto para verificar que las asignaciones siguen siendo vigentes, ya que esto asegura que la carga de trabajo por rol se mantenga balanceada.

\newpage
% --- Entregable 08 ---
\section{Organigrama del proyecto}

Esta sección detalla el organigrama del proyecto Plataforma Universitaria Integrada desarrollado para la Universidad Tecnológica La Mejor. El propósito de este esquema organizativo es definir la estructura básica del equipo, los roles principales y la relación jerárquica entre las personas que participan en la gestión y desarrollo de la iniciativa, permitiendo visualizar con claridad quién patrocina el proyecto, quién lo dirige y quiénes ejecutan las actividades principales. De esta manera, la jerarquía documentada ayuda a evitar confusiones sobre las responsabilidades y facilita la comunicación interna del grupo al contemplar roles específicos y bien definidos como el patrocinador, la directora del proyecto, los analistas de negocio y documentadores, además de los desarrolladores de la interfaz y control de calidad.

\subsection{Esquema general del organigrama}

\begin{figure}[H]
\centering
\resizebox{\textwidth}{!}{%
\begin{tikzpicture}[
    font=\scriptsize,
    text=diagramtext,
    box/.style={rounded corners=12pt, draw=none, fill=diagramlight, align=center, text width=3.6cm, minimum height=1.45cm, inner sep=6pt},
    mainbox/.style={box, fill=diagramgreen, text width=4.8cm, minimum height=1.85cm},
    midbox/.style={box, fill=diagramgreen, text width=5.0cm, minimum height=0.95cm},
    team/.style={box, fill=diagramgreen, text width=2.9cm, minimum height=0.85cm},
    connector/.style={dashed, line width=0.8pt, draw=black}
]

\node[mainbox] (sponsor) at (0,0) {\textbf{Dr. Roberto Mora\\Fallas}\\[5pt] Patrocinador\\[5pt] Rector de la Universidad Tecnológica La\\Mejor};

\node[midbox] (pm) at (0,-2.45) {\textbf{Tamara Robles Camacho}\\[5pt] Directora del proyecto};

\node[team] (team) at (0,-4.15) {\textbf{Equipo de trabajo}};

\node[box, text width=3.6cm, minimum height=1.75cm] (analyst) at (-6.25,-6.45) {\textbf{Isaac Jiménez Blanco}\\[4pt] -Analista de Neg. y Doc. 1\\-Documentación y control\\del proyecto};

\node[box, text width=3.35cm, minimum height=1.85cm] (ux) at (-2.1,-6.45) {\textbf{Jorge Abarca Quiros}\\[4pt] -Analista de Neg. y Doc. 2\\-Documentación y riesgos};

\node[box, text width=3.35cm, minimum height=1.85cm] (dev) at (2.1,-6.45) {\textbf{Mauricio González Prendas}\\[4pt] -Dev Interfaz 1 y QA\\-Portales básicos y QA};

\node[box, text width=3.6cm, minimum height=1.75cm] (qa) at (6.25,-6.45) {\textbf{Carlos Sainz Arce}\\[4pt] -Dev Interfaz 2 y QA\\-Portales fin/ejec. y QA};

\draw[connector] (sponsor.south) -- (pm.north);
\draw[connector] (pm.south) -- (team.north);
\draw[connector] (team.south) -- ++(0,-0.45) coordinate (branch);
\draw[connector] (branch) -| (analyst.north);
\draw[connector] (branch) -| (ux.north);
\draw[connector] (branch) -| (dev.north);
\draw[connector] (branch) -| (qa.north);

\end{tikzpicture}%
}
\caption{Organigrama principal del proyecto Plataforma Universitaria Integrada.}
\end{figure}

\clearpage
\begin{mediumwidepage}

\subsection{Interesados de apoyo y validación}

\begin{figure}[H]
\centering
\resizebox{0.96\linewidth}{!}{%
\begin{tikzpicture}[
    font=\scriptsize,
    text=diagramtext,
    box/.style={rounded corners=12pt, draw=none, fill=diagramlight, align=center, text width=3.2cm, minimum height=1.45cm, inner sep=5pt},
    mainbox/.style={box, fill=diagramgreen, text width=5.0cm, minimum height=1.1cm},
    connector/.style={dashed, line width=0.8pt, draw=black, line cap=round}
]

\node[mainbox] (top) at (0,0) {\textbf{Stakeholders externos\\de apoyo y validación}};

\coordinate (barcenter) at (0,-2.0);
\coordinate (leftbar) at (-10.0,-2.0);
\coordinate (rightbar) at (10.0,-2.0);

\node[box, text width=3.2cm, minimum height=1.55cm] (patricia) at (-10.0,-3.65) {\textbf{Lic. Patricia\\Vargas Soto}\\[7pt] Directora\\Administrativa};

\node[box, text width=3.2cm, minimum height=1.55cm] (carlos) at (-6.0,-5.35) {\textbf{Ing. Carlos Brenes\\Mora}\\[7pt] Director de TI};

\node[box, text width=3.2cm, minimum height=1.75cm] (andrea) at (-2.0,-3.95) {\textbf{Prof. Andrea Solís\\Zamora}\\[7pt] Representante\\docente};

\node[box, text width=3.2cm, minimum height=1.75cm] (estudiantes) at (2.0,-3.95) {\textbf{Estudiantes}\\[10pt] Usuarios finales del\\Portal Estudiantil};

\node[box, text width=3.25cm, minimum height=1.55cm] (financiero) at (6.0,-5.35) {\textbf{Departamento\\Financiero}\\[7pt] Área financiera de la\\institución};

\node[box, text width=3.25cm, minimum height=1.65cm] (hosting) at (10.0,-3.65) {\textbf{Proveedor\\de Hosting}\\[7pt] Soporte externo si se\\requiere alojamiento\\de evidencia};

\draw[connector] (top.south) -- (barcenter);
\draw[connector] (leftbar) -- (rightbar);
\draw[connector] (-10.0,-2.0) -- (patricia.north);
\draw[connector] (-6.0,-2.0) -- (carlos.north);
\draw[connector] (-2.0,-2.0) -- (andrea.north);
\draw[connector] (2.0,-2.0) -- (estudiantes.north);
\draw[connector] (6.0,-2.0) -- (financiero.north);
\draw[connector] (10.0,-2.0) -- (hosting.north);

\end{tikzpicture}%
}
\caption{Stakeholders externos de apoyo y validación del proyecto.}
\end{figure}

\end{mediumwidepage}
\clearpage

\subsection{Descripción breve de los roles}

En la siguiente tabla se describen los roles principales definidos para el proyecto, para lo cual se detalla la responsabilidad principal de cada integrante.

\begin{table}[H]
\centering
\small
\caption{Descripción de roles}\label{tab:roles}
\setlength{\tabcolsep}{5pt}
\renewcommand{\arraystretch}{1.25}
\begin{tabularx}{\textwidth}{p{4.0cm} p{4.4cm} Y}
\toprule
\textbf{Rol} & \textbf{Nombre} & \textbf{Responsabilidad principal} \\
\midrule
Patrocinador & Dr. Roberto Mora Fallas & Autorizar el proyecto, brindar apoyo institucional y aprobar decisiones importantes \\
Directora del proyecto & Tamara Robles Camacho & Coordinar el proyecto, dar seguimiento al avance y comunicar decisiones relevantes \\
Analista de Negocio y Doc. 1 & Isaac Jiménez Blanco & Apoyar la definición del alcance, documentación y control del proyecto \\
Analista de Negocio y Doc. 2 & Jorge Abarca Quiros & Elaborar entregables documentales, apoyar gestión de riesgos y costos \\
Desarrollador de la interfaz 1 y QA & Mauricio González Prendas & Construir las pantallas estudiantil, docente y admin, y realizar pruebas de calidad \\
Desarrollador de la interfaz 2 y QA & Carlos Sainz Arce & Construir el portal financiero y ejecutivo, y realizar pruebas de calidad \\
\bottomrule
\end{tabularx}
\end{table}

El equipo está conformado por cinco integrantes, por lo tanto esto permite una mayor especialización en las tareas de documentación y desarrollo técnico.

\subsection{Explicación del organigrama}

El organigrama se organiza en tres niveles principales, ya que en el primer nivel se encuentra el patrocinador, quien representa la autoridad institucional y tiene la capacidad de aprobar el proyecto, validar entregables importantes y apoyar cuando se presenten decisiones que afecten el alcance, el tiempo, el costo o la calidad.

En el segundo nivel se ubica la directora del proyecto, debido a que este rol se encarga de coordinar el trabajo del equipo, revisar avances, mantener comunicación con los interesados y asegurar que los entregables se completen dentro del plazo definido. Asimismo, en este proyecto la entrega final debe realizarse según el cronograma aprobado en Microsoft Project, esto implica que la coordinación del equipo es un punto clave.

En el tercer nivel se encuentra el equipo ejecutor, en consecuencia los analistas de negocio apoyan la comprensión del problema, la definición del alcance y elaboran la documentación del proyecto, mientras que los desarrolladores de la interfaz y QA se dividen la construcción de los módulos del prototipo web y revisan que el prototipo cumpla los criterios mínimos de calidad y navegación.

Los stakeholders externos aparecen como apoyo de validación, por lo tanto estos actores no ejecutan directamente el proyecto, pero pueden aportar información importante sobre procesos administrativos, técnicos, docentes, financieros y estudiantiles, ya que su participación ayuda a que el prototipo represente mejor las necesidades de la Universidad Tecnológica La Mejor.

\subsection{Relación entre roles y entregables}

La siguiente tabla muestra cómo se relacionan los roles del organigrama con los entregables principales del proyecto, de esta manera se define la responsabilidad de cada parte.

\begin{table}[H]
\centering
\small
\caption{Relación entre roles y entregables}\label{tab:roles-entregables}
\setlength{\tabcolsep}{5pt}
\renewcommand{\arraystretch}{1.25}
\begin{tabularx}{\textwidth}{p{4.3cm} Y}
\toprule
\textbf{Rol} & \textbf{Entregables o actividades relacionadas} \\
\midrule
Patrocinador & Aprobación del acta de constitución del proyecto, revisión de avances importantes y validación de cierre \\
Directora del proyecto & Seguimiento del cronograma, coordinación del equipo, control de cambios y comunicación \\
Analistas de Negocio y Doc. & Definición del alcance, documentación, riesgos, adquisiciones y evaluación de cambios \\
Desarrolladores de la interfaz y QA & Login, panel principal, módulos web, panel ejecutivo y pruebas de calidad \\
\bottomrule
\end{tabularx}
\end{table}

Esta distribución permite que cada rol tenga una función clara dentro del proyecto, además de que facilita la creación posterior del cronograma y la asignación de responsables en Microsoft Project.

\newpage
% --- Entregable 09 ---
\section{Plan de recursos}

\subsection{Recursos humanos}

El equipo del proyecto está conformado por cinco integrantes principales con dedicación completa durante la duración del proyecto, complementados por tres stakeholders clave con participación parcial en fases específicas.

\subsubsection{Equipo del proyecto}

En la siguiente tabla se muestran los roles de cada uno de los integrantes junto con su dedicación en el proyecto y las horas estimadas:

\begin{widepage}
{\small
\setlength{\tabcolsep}{3pt}
\begin{longtable}{L{4.3cm} L{4.6cm} L{2.2cm} L{3.1cm} L{5.1cm}}
\caption{Equipo del proyecto}\label{tab:equipo-proyecto}\\[-0.5em]
\toprule
\textbf{Rol} & \textbf{Nombre / perfil} & \textbf{Dedicación} & \textbf{Horas estimadas} & \textbf{Fase(s) asignada(s)} \\
\midrule
\endfirsthead
\toprule
\textbf{Rol} & \textbf{Nombre / perfil} & \textbf{Dedicación} & \textbf{Horas estimadas} & \textbf{Fase(s) asignada(s)} \\
\midrule
\endhead
Directora del proyecto & Tamara Robles Camacho & 100\% & 320 hrs & Todas las fases \\
Desarrollador de la interfaz \& QA 1 & Mauricio González Prendas & 100\% & 240 hrs & Ejecución / Pruebas \\
Desarrollador de la interfaz \& QA 2 & Carlos Sainz Arce & 100\% & 240h & Ejecución / Pruebas \\
Analista de Negocio \& Doc.1 & Isaac Jiménez Blanco & 100\% & 190 hrs & Inicio / Planificación / Cierre \\
Analista de Negocio \& Doc.2 & Jorge Abarca Quiros & 100\% & 190h & Inicio / Planificación / Cierre \\
Rector / Patrocinador & Dr. Roberto Mora Fallas & 5\% & 20 hrs & Inicio / Cierre \\
Directora Administrativa & Lic. Patricia Vargas Soto & 10\% & 30 hrs & Planificación / Ejecución \\
Director de TI & Ing. Carlos Brenes Mora & 10\% & 30 hrs & Planificación / Ejecución \\
\bottomrule
\end{longtable}
}
\end{widepage}

\subsubsection{Roles y responsabilidades clave}

Explicar los roles y sus responsabilidades es fundamental para llevar un orden en la elaboración del proyecto y evitar duplicación o mal información sobre responsabilidades, a continuación, se explican los roles y responsabilidades del equipo interno del proyecto:

\paragraph{Directora del proyecto (Tamara Robles Camacho)}

\begin{itemize}
    \item Responsable de la planificación, seguimiento y control del proyecto.
    \item Gestiona el registro de riesgos, cambios y comunicaciones con stakeholders.
    \item Elabora los entregables de gestión y coordina al equipo.
\end{itemize}

\paragraph{Desarrollador de la interfaz \& QA 1 (Mauricio González Prendas)}

Mauricio es el responsable de liderar la mitad del desarrollo técnico del prototipo y las pruebas de usabilidad del primer bloque de pantallas.

\begin{itemize}
    \item Realiza pruebas funcionales, de usabilidad y control de calidad de la interfaz.
    \item Responsable de la primera entrega del producto principal (interfaz navegable).
\end{itemize}

\paragraph{Desarrollador de la interfaz \& QA 2 (Carlos Sainz Arce)}

Carlos complementa el trabajo de Mauricio y lidera la segunda mitad del desarrollo y pruebas técnicas.

\begin{itemize}
    \item Realiza, junto con el desarrollador 1, pruebas funcionales, de usabilidad y control de calidad de la interfaz.
    \item Responsable de la segunda parte de la entrega del producto principal (interfaz navegable).
\end{itemize}

\paragraph{Analista de negocio \& documentador 1 (Isaac Jiménez Blanco)}

\begin{itemize}
    \item Elabora la mitad de los entregables documentales de planificación y cierre.
    \item Gestiona la documentación de cronograma y comunicaciones.
    \item Apoya, junto con el analista de negocio 2, al PM en la coordinación de actividades de inicio y planificación.
\end{itemize}

\paragraph{Analista de negocio \& documentador 2 (Jorge Abarca Quiros)}

\begin{itemize}
    \item Elabora la mitad de los entregables documentales de planificación y cierre.
    \item Gestiona la documentación de riesgos, costos y comunicaciones.
    \item Apoya, junto con él analista de negocio 1, al PM en la coordinación de actividades de inicio y planificación.
\end{itemize}

\subsubsection{Plan de capacitación}

Dado que el equipo es conformado por estudiantes universitarios con conocimiento técnico previo, se identifican las siguientes necesidades de nivelación:

\begin{itemize}
    \item MS Project: capacitación en línea (4 horas) para el responsable del cronograma.
    \item Metodología PMI/PMBOK: revisión guiada de verificación del PMI/PMBOK y plantillas (3 horas).
    \item React / Tailwind CSS: refuerzo técnico del desarrollador (autogestionado, 8 horas).
\end{itemize}

\subsection{Recursos tecnológicos}

A continuación, se detallan todas las herramientas de software, servicios en la nube y equipos de hardware requeridos para el desarrollo del proyecto.

\begin{widepage}
{\small
\setlength{\tabcolsep}{3pt}
\begin{longtable}{L{4cm} L{7cm} L{3cm} L{4cm} L{4.2cm}}
\caption{Recursos tecnológicos}\label{tab:recursos-tecnologicos}\\[-0.5em]
\toprule
\textbf{Recurso} & \textbf{Descripción / uso} & \textbf{Tipo} & \textbf{Costo} & \textbf{Responsable} \\
\midrule
\endfirsthead
\toprule
\textbf{Recurso} & \textbf{Descripción / uso} & \textbf{Tipo} & \textbf{Costo} & \textbf{Responsable} \\
\midrule
\endhead
VS Code & Editor de código principal para desarrollo frontend & Software & Gratuito & Mauricio \\
React + Vite & Framework y bundler para el frontend del sistema & Software & Gratuito & Mauricio \\
Tailwind CSS & Framework de estilos para la interfaz de usuario & Software & Gratuito & Mauricio \\
GitHub & Control de versiones y colaboración en código & Software & Gratuito & Mauricio \\
Figma & Diseño de wireframes y prototipos de interfaz & Software & Gratuito (plan free) & Mauricio \\
Microsoft Project & Elaboración del cronograma del proyecto & Software & Licencia comercial (\textasciitilde{}₡13,000/mes) & Isaac \\
Microsoft Word / Excel & Elaboración de entregables documentales & Software & Licencia Office 365 (\textasciitilde{}₡7,000/mes) & Todo el equipo \\
Cloudflare Workers & Hosting gratuito para despliegue del frontend demo & Servicio cloud & Gratuito & Mauricio \\
Correo y Teams & Comunicación formal del equipo en tiempo real, teams ofrece la posibilidad de reuniones virtuales del equipo y con stakeholders & Servicio & Gratuito & Las reuniones virtuales son responsables de Tamara y la comunicación de todo el equipo, \\
Computadoras personales (x5) & Equipos de trabajo personales de cada integrante & Hardware & Ya disponibles & Cada integrante \\
Conexión a Internet & Necesaria para reuniones virtuales y despliegue & Servicio & \textasciitilde{}₡25,000/mes c/u & Cada integrante \\
\bottomrule
\end{longtable}
}
\end{widepage}

\subsubsection{Justificación de herramientas seleccionadas}

En esta sección se justifican el porqué de la escogencia de dichas tecnologías y herramientas anteriormente mencionadas:

\begin{itemize}
    \item Se priorizan herramientas gratuitas o con licencia comercial simple para minimizar costos.
    \item React fue seleccionado por ser una de las tecnologías permitidas y por la experiencia previa del equipo.
    \item El uso del repositorio garantiza control de versiones y trazabilidad del código fuente.
    \item Cloudflare Workers permiten demostrar la interfaz funcional sin costo de infraestructura.
\end{itemize}

\subsection{Recursos materiales}

Los recursos materiales necesarios para este proyecto son mínimos dado su carácter digital. Se identifican los siguientes materiales físicos de apoyo:

\begin{table}[H]
\centering
\caption{Recursos materiales}\label{tab:recursos-materiales}
{\small
\setlength{\tabcolsep}{3pt}
\begin{tabularx}{\textwidth}{L{4cm} Y L{2.5cm} R{2.8cm}}
\toprule
\textbf{Material / recurso} & \textbf{Uso / descripción} & \textbf{Cantidad} & \textbf{Costo estimado} \\
\midrule
Impresiones de documentos & Entregables formales impresos para revisión presencial & \textgreater{}80 páginas & ₡9,000 \\
Cuadernos / libretas & Apuntes de reuniones y planificación manual & 5 unidades & ₡9,500 \\
Bolígrafos y marcadores & Uso en sesiones presenciales y pizarras & Set básico & ₡5,000 \\
Pizarra (uso en las oficinas) & Sesiones de planificación y diseño colaborativo & 1 & ₡0 \\
USB / almacenamiento externo & Respaldo de archivos del proyecto & 1 por integrante & ₡10,000 \\
Carpetas y archivadores & Organización física de entregables impresos & 5 unidades & ₡6,500 \\
\midrule
\multicolumn{3}{l}{\textbf{Total de recursos materiales}} & \textbf{₡40,000} \\
\bottomrule
\end{tabularx}
}
\end{table}

\subsection{Resumen de recursos del proyecto}

La siguiente tabla integra todos los recursos identificados por categoría:

\begin{table}[H]
\centering
\caption{Resumen de recursos del proyecto}\label{tab:resumen-recursos}
{\small
\setlength{\tabcolsep}{3pt}
\begin{tabularx}{\textwidth}{L{5.5cm} Y R{4.3cm}}
\toprule
\textbf{Categoría} & \textbf{Descripción} & \textbf{Costo estimado} \\
\midrule
Recursos Humanos & 5 integrantes + 3 stakeholders parciales & ₡14,000,000 \\
Recursos tecnológicos: software & 10 herramientas identificadas, 3 con pago & ₡45,000 \\
Recursos tecnológicos: hardware & 5 computadoras personales & ₡0 \\
Recursos tecnológicos: servicios & Internet, hosting, reuniones virtuales & ₡75,000 \\
Recursos Materiales & Impresiones, papelería, USBs & ₡30,000 \\
\midrule
\multicolumn{2}{l}{\textbf{Inversión total estimada del proyecto}} & \textbf{₡16,100,000} \\
\bottomrule
\end{tabularx}
}
\end{table}

\newpage
% --- Entregable 10 ---
\section{Estimación de costos}

\subsection{Metodología}

La estimación de costos es el proceso de predecir los recursos financieros necesarios para completar cada componente del proyecto. Para este proyecto se aplica una técnica (PERT) aprobada por el PMBOK:

\begin{table}[H]
\centering
\caption{Técnica aplicada}
\label{tab:tecnica-aplicada}
{\small
\begin{tabularx}{\textwidth}{P{3.8cm} Y Y}
\toprule
\textbf{Técnica} & \textbf{Descripción} & \textbf{Aplicación en este proyecto} \\
\midrule
PERT (Tres Puntos) & Considera incertidumbre con tres escenarios: Optimista (O), Más Probable (M) y Pesimista (P). Fórmula: (O + 4M + P) / 6 & Se aplica a los requisitos por semana, con sus tareas, además de los rubros adicionales: software, hardware, servicios y contingencias. \\
\bottomrule
\end{tabularx}
}
\end{table}

\subsection{Estimación PERT}

En la siguiente tabla se muestra la estimación de costos detallada para cada tarea utilizando la técnica PERT:

\begin{widepage}
{\footnotesize
\setlength{\tabcolsep}{3pt}
\begin{longtable}{P{7.9cm} R{3.0cm} R{3.2cm} R{3.0cm} R{4.8cm}}
\caption{Estimación PERT del proyecto}\label{tab:estimacion-pert}\\[-0.5em]
\toprule
\rowcolor{tecblue}
\multicolumn{5}{c}{\textcolor{white}{\textbf{Estimación PERT de la Plataforma Universitaria Integrada}}} \\
\rowcolor{teclightblue}
\multicolumn{5}{c}{Universidad Tecnológica La Mejor  |  I Semestre 2026  |  Plazo: 3 meses  |  Equipo: 5 personas  |  Moneda: Colones (₡)} \\
\midrule
\rowcolor{tablegray}
\textbf{Tareas} & \textbf{Optimista} & \textbf{Más probable} & \textbf{Pesimista} & \textbf{Costo estimado} \par \textbf{(O + 4M + P) / 6} \\
\midrule
\endfirsthead
\toprule
\rowcolor{tablegray}
\textbf{Tareas} & \textbf{Optimista} & \textbf{Más probable} & \textbf{Pesimista} & \textbf{Costo estimado} \par \textbf{(O + 4M + P) / 6} \\
\midrule
\endhead
\rowcolor{diagramlight}
\multicolumn{5}{l}{\textbf{Fase 1: Inicio (Semanas 1 a 2)}} \\
\multicolumn{5}{l}{\textbf{\textit{Gestión de Inicio (Business Case, Charter, Stakeholders)}}} \\
\hspace{0.5cm}Elaborar Business Case completo & ₡480.000 & ₡700.000 & ₡900.000 & ₡696.667 \\
\hspace{0.5cm}Desarrollar Project Charter & ₡300.000 & ₡450.000 & ₡600.000 & ₡450.000 \\
\hspace{0.5cm}Registrar y analizar stakeholders & ₡200.000 & ₡280.000 & ₡420.000 & ₡290.000 \\
\hspace{0.5cm}Coordinación y gestión del PM (Inicio) & ₡300.000 & ₡450.000 & ₡600.000 & ₡450.000 \\
\midrule
\textbf{Subtotal} &  &  &  & \textbf{₡1.886.667} \\
\addlinespace[0.35em]
\rowcolor{diagramlight}
\multicolumn{5}{l}{\textbf{Fase 2: Planificación (Semanas 2 a 5)}} \\
\multicolumn{5}{l}{\textbf{\textit{Documentación PMI: alcance, WBS, cronograma}}} \\
\hspace{0.5cm}Definir alcance incluido/excluido & ₡200.000 & ₡280.000 & ₡420.000 & ₡290.000 \\
\hspace{0.5cm}Elaborar WBS con criterios de aceptación & ₡280.000 & ₡360.000 & ₡540.000 & ₡376.667 \\
\hspace{0.5cm}Cronograma en MS Project (tareas, hitos) & ₡350.000 & ₡480.000 & ₡720.000 & ₡498.333 \\
\multicolumn{5}{l}{\textbf{\textit{Documentación PMI: recursos, costos, calidad}}} \\
\hspace{0.5cm}RACI, organigrama y Plan de Recursos & ₡200.000 & ₡280.000 & ₡420.000 & ₡290.000 \\
\hspace{0.5cm}Estimación de costos (PERT) & ₡180.000 & ₡240.000 & ₡360.000 & ₡250.000 \\
\hspace{0.5cm}Plan de Recursos & ₡110.000 & ₡220.000 & ₡330.000 & ₡220.000 \\
\hspace{0.5cm}Plan de Calidad con métricas & ₡160.000 & ₡220.000 & ₡330.000 & ₡228.333 \\
\multicolumn{5}{l}{\textbf{\textit{Documentación PMI: comunicaciones, riesgos, adquisiciones}}} \\
\hspace{0.5cm}Plan de Comunicaciones & ₡140.000 & ₡200.000 & ₡300.000 & ₡206.667 \\
\hspace{0.5cm}Plan Gestión de Riesgos + Registro & ₡280.000 & ₡380.000 & ₡570.000 & ₡395.000 \\
\hspace{0.5cm}Matriz de Probabilidad e Impacto & ₡120.000 & ₡180.000 & ₡270.000 & ₡185.000 \\
\hspace{0.5cm}Plan de Adquisiciones & ₡120.000 & ₡180.000 & ₡270.000 & ₡185.000 \\
\hspace{0.5cm}Gestión y coordinación del PM (Planif.) & ₡400.000 & ₡540.000 & ₡810.000 & ₡561.667 \\
\midrule
\textbf{Subtotal} &  &  &  & \textbf{₡3.686.667} \\
\addlinespace[0.35em]
\rowcolor{diagramlight}
\multicolumn{5}{l}{\textbf{Fase 3: Ejecución y desarrollo de la interfaz (Semanas 4 a 10)}} \\
\multicolumn{5}{l}{\textbf{\textit{Módulos del sistema: Dev Interfaz 1 (Mauricio)}}} \\
\hspace{0.5cm}Login y panel principal & ₡580.000 & ₡640.000 & ₡980.000 & ₡686.667 \\
\hspace{0.5cm}Portal estudiantil: matrícula y horarios & ₡670.000 & ₡890.000 & ₡1.280.000 & ₡918.333 \\
\hspace{0.5cm}Portal estudiantil: calificaciones e historial & ₡490.000 & ₡660.000 & ₡980.000 & ₡685.000 \\
\hspace{0.5cm}Portal docente: gestión cursos y asistencia & ₡500.000 & ₡700.000 & ₡960.000 & ₡710.000 \\
\hspace{0.5cm}Portal docente: registro de calificaciones & ₡360.000 & ₡490.000 & ₡720.000 & ₡506.667 \\
\hspace{0.5cm}Portal administrativo: CRUD completo & ₡690.000 & ₡890.000 & ₡1.270.000 & ₡920.000 \\
\multicolumn{5}{l}{\textbf{\textit{Módulos del sistema: Dev Interfaz 2 (Carlos)}}} \\
\hspace{0.5cm}Portal financiero: pagos e historial & ₡480.000 & ₡700.000 & ₡960.000 & ₡706.667 \\
\hspace{0.5cm}Dashboard Ejecutivo del sistema & ₡480.000 & ₡640.000 & ₡960.000 & ₡666.667 \\
\multicolumn{5}{l}{\textbf{\textit{Módulos compartidos: ambos devs}}} \\
\hspace{0.5cm}Integración entre módulos y navegación & ₡320.000 & ₡500.000 & ₡660.000 & ₡496.667 \\
\hspace{0.5cm}Pruebas de usabilidad y QA (ambos devs) & ₡480.000 & ₡710.000 & ₡960.000 & ₡713.333 \\
\midrule
\textbf{Subtotal} &  &  &  & \textbf{₡7.010.000} \\
\addlinespace[0.35em]
\rowcolor{diagramlight}
\multicolumn{5}{l}{\textbf{Fase 4: Monitoreo y control (Semanas 1 a 12, continuo)}} \\
\multicolumn{5}{l}{\textbf{\textit{Seguimiento y Control}}} \\
\hspace{0.5cm}Definir y monitorear KPIs del proyecto & ₡200.000 & ₡350.000 & ₡420.000 & ₡336.667 \\
\hspace{0.5cm}Control de cambios + Solicitudes \#1 y \#2 & ₡240.000 & ₡320.000 & ₡480.000 & ₡333.333 \\
\hspace{0.5cm}Evaluación de impacto de cambios & ₡160.000 & ₡270.000 & ₡330.000 & ₡261.667 \\
\hspace{0.5cm}Seguimiento cronograma y riesgos activos & ₡200.000 & ₡280.000 & ₡420.000 & ₡290.000 \\
\hspace{0.5cm}Dashboard Ejecutivo del proyecto & ₡200.000 & ₡350.000 & ₡420.000 & ₡336.667 \\
\hspace{0.5cm}Gestión PM (reuniones, reportes, coordinación) & ₡480.000 & ₡710.000 & ₡960.000 & ₡713.333 \\
\midrule
\textbf{Subtotal} &  &  &  & \textbf{₡2.271.667} \\
\addlinespace[0.35em]
\rowcolor{diagramlight}
\multicolumn{5}{l}{\textbf{Fase 5: Cierre (Semana 12)}} \\
\multicolumn{5}{l}{\textbf{\textit{Entregables de Cierre}}} \\
\hspace{0.5cm}Informe Ejecutivo final & ₡120.000 & ₡220.000 & ₡270.000 & ₡211.667 \\
\hspace{0.5cm}Acta de Cierre formal & ₡100.000 & ₡190.000 & ₡210.000 & ₡178.333 \\
\hspace{0.5cm}Reflexiones y Recomendaciones del equipo & ₡120.000 & ₡200.000 & ₡270.000 & ₡198.333 \\
\midrule
\textbf{Subtotal} &  &  &  & \textbf{₡588.333} \\
\addlinespace[0.35em]
\rowcolor{diagramlight}
\multicolumn{5}{l}{\textbf{Rubros adicionales (software, hardware, servicios)}} \\
\multicolumn{5}{l}{\textbf{\textit{Software}}} \\
\hspace{0.5cm}MS Project licencia comercial (3 meses) & ₡30.000 & ₡70.000 & ₡90.000 & ₡66.667 \\
\hspace{0.5cm}Office 365 para 5 usuarios (3 meses) & ₡90.000 & ₡130.000 & ₡180.000 & ₡131.667 \\
\hspace{0.5cm}Herramientas premium opcionales & ₡0 & ₡20.000 & ₡60.000 & ₡23.333 \\
\multicolumn{5}{l}{\textbf{\textit{Hardware}}} \\
\hspace{0.5cm}USB / almacenamiento externo (5 unidades) & ₡10.000 & ₡20.000 & ₡25.000 & ₡19.167 \\
\hspace{0.5cm}Periféricos complementarios & ₡5.000 & ₡15.000 & ₡20.000 & ₡14.167 \\
\multicolumn{5}{l}{\textbf{\textit{Servicios}}} \\
\hspace{0.5cm}Internet (5 personas x 3 meses x ₡25,000) & ₡225.000 & ₡370.000 & ₡450.000 & ₡359.167 \\
\hspace{0.5cm}Impresiones y material físico & ₡6.000 & ₡10.000 & ₡17.000 & ₡10.500 \\
\hspace{0.5cm}Transporte a reuniones presenciales & ₡20.000 & ₡42.000 & ₡60.000 & ₡41.333 \\*
\midrule
\textbf{Subtotal} &  &  &  & \textbf{₡666.000} \\*
\midrule
Subtotal antes de contingencias &  &  &  & \textbf{₡16.109.333} \\*
Reserva para contingencias (10 \%) &  &  &  & \textbf{₡1.515.867} \\*
\midrule
\rowcolor{teclightblue}
\multicolumn{4}{l}{\textbf{Presupuesto total del proyecto (PERT + contingencias)}} & \textbf{₡17.625.200} \\*
\addlinespace[0.35em]
\rowcolor{tablegray}
\multicolumn{5}{l}{\textbf{Notas:}} \\
\multicolumn{5}{p{22.74cm}}{\scriptsize \textbullet{} Fórmula PERT: Costo Estimado = (Optimista + 4 $\times$ Más Probable + Pesimista) / 6} \\
\multicolumn{5}{p{22.74cm}}{\scriptsize \textbullet{} El escenario Optimista asume todo sale bien, el Pesimista incluye retrasos, correcciones y sobrecostos.} \\
\multicolumn{5}{p{22.74cm}}{\scriptsize \textbullet{} Equipo: Tamara Robles (Dir.), Mauricio González (Dev Interfaz 1), Carlos Sainz Arce (Dev Interfaz 2), Isaac Jiménez (Analista 1), Jorge Abarca Quiros (Analista 2).} \\
\multicolumn{5}{p{22.74cm}}{\scriptsize \textbullet{} Plazo total del proyecto: 3 meses (12 semanas). Moneda: Colones Costarricenses (₡).} \\
\multicolumn{5}{p{22.74cm}}{\scriptsize \textbullet{} Los costos de mano de obra se basan en tarifas de mercado TI Costa Rica para perfiles junior-medio (de ₡8,000 a ₡10,000/hr).} \\
\bottomrule
\end{longtable}
}
\end{widepage}

\subsection{Reserva para contingencias}

Según el estándar PMI, toda estimación debe incluir una reserva para contingencias que cubra riesgos identificados e imprevistos menores. Para este proyecto se aplica el 10\% sobre el subtotal de todos los rubros, que es la práctica estándar en proyectos de TI de mediana escala.

\begin{table}[H]
\centering
\caption{Reserva para contingencias}
\label{tab:reserva-contingencias}
\begin{tabularx}{\textwidth}{Y R{4cm}}
\toprule
\textbf{Concepto} & \textbf{Monto} \\
\midrule
Subtotal de la Fase 1 (Inicio) & ₡1 886 667 \\
Subtotal de la Fase 2 (Planificación) & ₡3 686 667 \\
Subtotal de la Fase 3 (Desarrollo de interfaz) & ₡7 010 000 \\
Subtotal de la Fase 4 (Monitoreo y control) & ₡2 271 667 \\
Subtotal de la Fase 5 (Cierre) & ₡588 333 \\
Subtotal de rubros adicionales (software, hardware y servicios) & ₡666 000 \\
\midrule
Subtotal antes de contingencias & ₡16 109 333 \\
Reserva para contingencias (10 \%) & ₡1 515 867 \\
\midrule
\textbf{Presupuesto total del proyecto} & \textbf{₡17 625 200} \\
\bottomrule
\end{tabularx}
\end{table}

\subsection{Resumen ejecutivo del presupuesto}

La siguiente tabla consolida todos los rubros identificados y presenta la distribución porcentual del presupuesto total del proyecto:

\begin{table}[H]
\centering
\caption{Resumen ejecutivo del presupuesto}
\label{tab:resumen-presupuesto}
\begin{tabularx}{\textwidth}{Y R{3.2cm} R{2.8cm} P{3cm}}
\toprule
\textbf{Rubro} & \textbf{Monto estimado} & \textbf{\% del total} & \textbf{Técnica usada} \\
\midrule
Fase 1 (Inicio) & ₡1 886 667 & 10.70\% & PERT \\
Fase 2 (Planificación) & ₡3 686 667 & 20.92\% & PERT \\
Fase 3 (Desarrollo de interfaz) & ₡7 010 000 & 39.77\% & PERT \\
Fase 4 (Monitoreo y Control) & ₡2 271 667 & 12.89\% & PERT \\
Fase 5 (Cierre) & ₡588 333 & 3.34\% & PERT \\
Subtotal de rubros adicionales (software, hardware y servicios) & ₡666 000 & 3.34\% & PERT \\
Contingencias (10\%) & ₡1 515 867 & 8.60\% & \% sobre subtotal \\
\midrule
\textbf{Presupuesto total del proyecto} & \textbf{₡17 625 200} & \textbf{100\%} & \\
\bottomrule
\end{tabularx}
\end{table}

\subsubsection{Observaciones clave}

\begin{itemize}
    \item La mano de obra representa el 87.62\% del presupuesto total, lo que es típico en proyectos de software donde el capital humano es el principal activo.
    \item El uso de herramientas gratuitas (GitHub, React, VS Code) redujo significativamente el rubro de software.
    \item Las contingencias cubren principalmente retrasos en entregables documentales y ajustes de alcance no previstos.
\end{itemize}

\newpage
% --- Entregable 11 ---
\section{Plan de calidad}

Esta sección detalla el Plan de Calidad del proyecto Plataforma Universitaria Integrada desarrollado para la Universidad Tecnológica La Mejor con el objetivo principal de definir cómo se asegurará que los entregables cumplan con los requisitos establecidos desde las fases iniciales de planificación. El proyecto surge como respuesta a la necesidad de modernizar procesos académicos y administrativos que actualmente operan de manera fragmentada generando duplicidad e ineficiencia, sin embargo, debido a que el alcance de esta entrega se limita a un prototipo de interfaz web navegable en lugar de un sistema completo en producción, los criterios de calidad se adaptan específicamente a evaluar la representación visual de los módulos junto con la consistencia de navegación y la alineación con la documentación administrativa.

Para mantener la coherencia a lo largo de toda la iniciativa técnica se toman como referencia fundamental los documentos organizacionales elaborados previamente que establecen elementos esenciales como el propósito, las restricciones y los criterios generales de aceptación del producto final. En consecuencia y considerando la naturaleza académica del trabajo junto con su tiempo limitado de ejecución, este plan prioriza las revisiones manuales, las validaciones de usabilidad básica y el control de defectos visibles en el diseño gráfico, excluyendo deliberadamente aspectos de mayor complejidad técnica como la lógica de servidor, la integración con bases de datos reales o las pruebas de carga productivas.

\subsection{Objetivos de calidad}

Los objetivos de calidad definen claramente lo que se espera lograr en la iniciativa para considerar que tanto el producto como la gestión cumplen con lo planificado de modo que puedan servir como referencia principal durante el desarrollo y la validación final del sistema.

\begin{itemize}
    \item Garantizar que el prototipo de interfaz web navegable represente correctamente los módulos principales definidos en el alcance permitiendo una navegación clara junto con el uso de datos simulados y la visualización de la información mínima establecida para cada portal correspondiente.
    \item Asegurar que todos los entregables documentales mantengan una coherencia transversal con la documentación estratégica previamente aprobada para mantener los mismos nombres de roles o fechas y criterios generales de aceptación en toda la documentación elaborada.
    \item Reducir drásticamente la presencia de errores críticos de navegación en el prototipo antes de la entrega final considerando como error crítico cualquier problema que impida ingresar a un módulo principal o continuar un flujo básico de navegación.
    \item Verificar constantemente que el proyecto se mantenga dentro del alcance aprobado y que cualquier cambio relevante sea analizado mediante el proceso formal de control de cambios para evitar agregar funcionalidades innecesarias que puedan afectar el tiempo o los costos.
\end{itemize}

\subsection{Normas que se aplicarán}

Las normas establecen reglas internas fundamentales que el equipo deberá seguir rigurosamente para asegurar un desarrollo ordenado y coherente por lo cual en este proyecto se aplicarán las directrices preventivas que se detallan a continuación.

\begin{itemize}
    \item \textbf{Protección de datos simulados:} El proyecto no utilizará datos reales de estudiantes o personal administrativo sino que todos los nombres y calificaciones mostradas en el prototipo serán datos simulados para asegurar que no se incluya información sensible que pueda confundirse con datos reales de la institución.
    \item \textbf{Control de acceso representativo:} El inicio de sesión será una simulación visual que no validará usuarios reales dado que la interfaz deberá diferenciar visualmente los portales principales según el tipo de usuario sin prometer autenticación institucional real al estar fuera del alcance del proyecto.
    \item \textbf{Seguridad básica de interfaz:} Las pantallas deberán presentar una navegación clara y controlada donde el prototipo no solicitará información personal innecesaria ni mostrará mensajes que expongan información sensible del sistema.
    \item \textbf{Consistencia documental:} Todos los documentos del proyecto deberán mantener de manera transversal el mismo nombre del proyecto y los roles principales alineados con el registro de interesados junto con el presupuesto y las fechas respetando la planificación definida originalmente en Microsoft Project.
    \item \textbf{Control de cambios:} Cualquier cambio que afecte el alcance o los costos del proyecto deberá ser tratado mediante el proceso formal de control de cambios evaluando rigurosamente las solicitudes antes de modificar documentos o el prototipo sin agregar funcionalidades no contempladas.
    \item \textbf{Validación de información institucional:} Los procesos académicos y administrativos representados en el prototipo deberán mantenerse a un nivel general académico dado que la plataforma no debe presentarse como un sistema institucional real sino utilizar indicadores simulados coherentes con el propósito de apoyar la toma de decisiones.
\end{itemize}

\subsection{Estándares que se aplicarán}

Los estándares permiten evaluar el proyecto con criterios técnicos y de gestión reconocidos a nivel internacional por lo que para la Plataforma Universitaria Integrada se adoptan las metodologías y buenas prácticas descritas en los siguientes apartados.

\begin{itemize}
    \item \textbf{ISO/IEC 25010 como guía de calidad del producto:} Esta norma internacional se utilizará como referencia principal para evaluar exhaustivamente las características de calidad del prototipo priorizando en este proyecto factores críticos como la adecuación funcional junto con la usabilidad y la seguridad básica además de la eficiencia de rendimiento y la compatibilidad general.
    \item \textbf{Buenas prácticas de gestión de proyectos según PMI:} El proyecto se gestionará integralmente mediante entregables de inicio y planificación seguidos de la ejecución junto con el monitoreo y cierre lo cual permite mantener una estructura sumamente ordenada y totalmente coherente con los documentos formales solicitados. Adicionalmente para el desarrollo visual de la interfaz se mantendrá una organización clara de archivos aplicando nombres comprensibles en todos los módulos y asegurando una navegación consistente entre páginas mediante el uso de un repositorio para controlar los cambios de código antes de la entrega final.
    \item \textbf{Accesibilidad web básica y seguridad de la interfaz:} Aunque el proyecto no busca una certificación formal de accesibilidad se tomarán como referencia buenas prácticas básicas garantizando textos legibles y contraste adecuado junto con menús identificables y diseño responsivo cuando sea posible. Asimismo debido a que no existe un backend ni autenticación real no se aplicarán pruebas de seguridad profundas pero se evitará solicitar datos sensibles innecesarios manteniendo flujos de navegación controlados sin mostrar información confidencial ni formularios confusos.
    \item \textbf{Estándar interno de documentación:} Los documentos deberán estar redactados rigurosamente con una estructura clara y títulos consistentes asegurando en todo momento evitar cualquier contradicción fundamental entre el alcance del proyecto junto con su cronograma y costos además de los roles y criterios de aceptación.
\end{itemize}

\subsection{Métricas de calidad}

Las métricas establecidas permiten verificar cuantitativamente si el proyecto cumple con los objetivos definidos por lo cual se utilizarán indicadores simples debido a que el alcance corresponde a un prototipo académico de interfaz navegable y no a un sistema real en producción.

\begin{table}[H]
\centering
\caption{Métricas de calidad}
\label{tab:metricas-calidad}
{\small
\setlength{\tabcolsep}{3.5pt}
\begin{tabularx}{\textwidth}{p{3.8cm} Y Y}
\toprule
\textbf{Área evaluada} & \textbf{Métrica} & \textbf{Meta de calidad} \\
\midrule
Navegación del prototipo & Módulos principales accesibles desde el dashboard & 100\% de módulos definidos en el alcance \\
Funcionalidad visual & Pantallas mínimas representadas por módulo & 100\% de pantallas principales \\
Errores críticos & Errores que impiden navegar o visualizar módulos & 0 errores críticos al momento de entrega \\
Usabilidad & Claridad de menús, botones y flujo de navegación & Navegación entendible sin instrucciones extensas \\
Documentación & Coherencia entre entregables & Sin contradicciones importantes en roles, fechas, costos o alcance \\
Cronograma & Cumplimiento de fecha final & Entrega según el cronograma aprobado \\
Costos & Alineación con presupuesto base & Mantener total de ₡17,625,200 \\
Cambios & Cambios relevantes gestionados formalmente & 100\% de cambios importantes evaluados \\
Calidad visual & Consistencia de colores, distribución y legibilidad & Interfaz uniforme en módulos principales \\
Evidencia & Capturas, archivos y documentos organizados & Repositorio final completo y accesible \\
\bottomrule
\end{tabularx}
}
\end{table}

\subsection{Actividades de aseguramiento y control}

Las actividades conjuntas de aseguramiento y control permiten garantizar que la calidad se revise progresivamente durante el desarrollo del proyecto y no únicamente al final dado que el aseguramiento preventivo busca evitar defectos mientras que el control técnico se enfoca en detectar y corregir errores inmediatamente antes de la entrega final.

\subsubsection{Aseguramiento de calidad}

\paragraph*{Revisión temprana del alcance}

\begin{itemize}
    \item Confirmar que el prototipo solo incluya los módulos definidos.
    \item Verificar que no se agreguen funciones fuera del alcance.
    \item Revisar que el alcance excluido se respete durante el desarrollo.
\end{itemize}

\paragraph*{Revisión de coherencia documental}

\begin{itemize}
    \item Comparar toda la documentación técnica y estratégica aprobada previamente para evitar posibles contradicciones.
    \item Verificar que los roles del equipo sean consistentes en todos los documentos.
    \item Revisar que el presupuesto total sea el mismo en los entregables donde aparezca.
    \item Confirmar que la fecha final del proyecto coincida con el cronograma aprobado en Microsoft Project.
\end{itemize}

\paragraph*{Revisión del cronograma}

\begin{itemize}
    \item Verificar que las tareas tengan duración, fechas, dependencias, recursos, responsables, costos e hitos.
    \item Confirmar que la ruta crítica esté identificada en Microsoft Project.
    \item Revisar que el cronograma sea coherente con el alcance y los recursos disponibles.
\end{itemize}

\paragraph*{Revisión de diseño y navegación}

\begin{itemize}
    \item Revisar que los módulos principales tengan una estructura visual clara.
    \item Verificar que el Dashboard principal permita acceder a los portales definidos.
    \item Revisar que los nombres de secciones sean entendibles para estudiantes, docentes, administración y área financiera.
\end{itemize}

\paragraph*{Revisión de roles y responsabilidades}

\begin{itemize}
    \item Validar que Tamara Robles Camacho mantenga el rol de Directora del proyecto.
    \item Validar que Isaac Jiménez Blanco mantenga el rol de Analista de Negocio y Documentador 1.
    \item Validar que Jorge Abarca Quiros mantenga el rol de Analista de Negocio y Documentador 2.
    \item Validar que Mauricio González Prendas mantenga el rol de Desarrollador de interfaz 1 y control de calidad.
    \item Validar que Carlos Sainz Arce mantenga el rol de Desarrollador de interfaz 2 y control de calidad.
    \item Confirmar que los interesados externos sean considerados como apoyo de validación, no como ejecutores directos.
\end{itemize}

\subsubsection{Control de calidad}

\paragraph*{Pruebas manuales de navegación}

\begin{itemize}
    \item Ingresar al prototipo desde el login simulado.
    \item Acceder al Dashboard principal.
    \item Entrar al Portal Estudiantil.
    \item Entrar al Portal Docente.
    \item Entrar al Portal Administrativo.
    \item Entrar al Portal Financiero.
    \item Entrar al Dashboard Ejecutivo.
    \item Verificar que cada enlace principal funcione.
\end{itemize}

\paragraph*{Pruebas de contenido visual}

\begin{itemize}
    \item Revisar que cada módulo muestre la información mínima definida.
    \item Confirmar que los datos utilizados sean simulados.
    \item Verificar que los títulos, botones y menús sean consistentes.
\end{itemize}

\paragraph*{Pruebas de usabilidad básica}

\begin{itemize}
    \item Revisar si un usuario puede identificar fácilmente dónde entrar.
    \item Revisar si el flujo de navegación es claro.
    \item Verificar que no existan pantallas confusas o sin salida.
\end{itemize}

\paragraph*{Pruebas de compatibilidad básica}

\begin{itemize}
    \item Probar el prototipo en al menos un navegador principal.
    \item Si el tiempo lo permite, revisar en dos navegadores.
    \item Verificar que el diseño no se rompa en una pantalla de escritorio común.
\end{itemize}

\paragraph*{Control de defectos}

\begin{itemize}
    \item Registrar errores encontrados durante las pruebas.
    \item Clasificar errores como críticos, medios o bajos.
    \item Corregir errores críticos antes de la entrega.
    \item Revalidar los módulos corregidos.
\end{itemize}

\paragraph*{Revisión final de entregables}

\begin{itemize}
    \item Verificar nombres de archivos.
    \item Revisar ortografía básica.
    \item Confirmar que los documentos estén completos.
    \item Validar que el prototipo y los documentos estén organizados en la carpeta de entrega.
\end{itemize}

\subsection{Herramientas}

Las herramientas de apoyo permiten organizar el trabajo, controlar versiones, revisar calidad y dejar evidencia del avance del proyecto. Para este Plan de Calidad se utilizarán las siguientes:

\paragraph*{Gestión del cronograma}

\begin{itemize}
    \item Microsoft Project Profesional 2024 para tareas, subtareas, dependencias, duración, recursos, responsables, costos, hitos, ruta crítica y línea base.
\end{itemize}

\paragraph*{Gestión documental}

\begin{itemize}
    \item Microsoft Word para la elaboración de entregables.
    \item Repositorio centralizado para almacenar y compartir los entregables finales del proyecto.
    \item PDF para entregar evidencias del cronograma, ruta crítica, recursos y estadísticas del proyecto.
\end{itemize}

\paragraph*{Control de versiones}

\begin{itemize}
    \item Repositorio de control de versiones para administrar el código del prototipo de interfaz.
    \item Entregas de código y ramas cuando sea necesario para registrar avances del desarrollo.
\end{itemize}

\paragraph*{Diseño y apoyo visual}

\begin{itemize}
    \item Figma, Canva u otra herramienta visual únicamente como apoyo de diseño.
    \item Estas herramientas no sustituyen el prototipo web navegable, ya que el alcance exige desarrollo web.
\end{itemize}

\paragraph*{Desarrollo de interfaz}

\begin{itemize}
    \item Tecnologías web permitidas como HTML, CSS, JavaScript, React, Angular o Vue.
    \item Editor de código seleccionado por el equipo.
\end{itemize}

\paragraph*{Pruebas y control de calidad}

\begin{itemize}
    \item Checklist manual de pruebas.
    \item Navegadores web para validar navegación.
    \item Capturas de pantalla como evidencia.
    \item Revisión entre compañeros antes de la entrega.
\end{itemize}

\paragraph*{Comunicación del equipo}

\begin{itemize}
    \item Reuniones cortas de seguimiento.
    \item Mensajes grupales para resolver dudas y coordinar cambios.
\end{itemize}

\subsection{Responsables}

La definición de responsables permite que cada actividad de calidad tenga un dueño claro. En este proyecto, algunos roles son combinados debido al tamaño del equipo y al alcance académico del trabajo.

\subsubsection*{Patrocinador}
Dr. Roberto Mora Fallas

\begin{itemize}
    \item Aprueba el Project Charter.
    \item Revisa avances importantes.
    \item Valida el cierre del proyecto.
    \item Representa el interés institucional de modernizar los procesos universitarios.
\end{itemize}

\subsubsection*{Directora del proyecto}
Tamara Robles Camacho

\begin{itemize}
    \item Coordina el trabajo del equipo.
    \item Da seguimiento al cronograma, costos y entregables.
    \item Verifica que el proyecto se mantenga dentro del alcance.
    \item Coordina decisiones relacionadas con cambios, riesgos y comunicación.
    \item Revisa que el equipo cumpla con las fechas definidas.
\end{itemize}

\subsubsection*{Analista de Negocio y Documentador 1}
Isaac Jiménez Blanco

\begin{itemize}
    \item Elabora y revisa documentos de planificación y control (mitad de entregables).
    \item Mantiene coherencia entre alcance, cronograma, calidad, riesgos, adquisiciones y cambios.
    \item Define criterios de aceptación y puntos de revisión.
    \item Apoya la validación de que los módulos representen los procesos principales de la universidad.
\end{itemize}

\subsubsection*{Analista de Negocio y Documentador 2}
Jorge Abarca Quiros

\begin{itemize}
    \item Elabora y revisa documentos de planificación y control (mitad de entregables).
    \item Apoya la gestión de riesgos, costos y comunicaciones del proyecto.
    \item Mantiene actualizada la documentación del proyecto.
\end{itemize}

\subsubsection*{Desarrollador de interfaz 1 y control de calidad}
Mauricio González Prendas

\begin{itemize}
    \item Desarrolla el prototipo web navegable (módulos principales).
    \item Implementa pantallas, módulos y navegación principal.
    \item Ejecuta pruebas manuales de navegación y funcionalidad visual.
    \item Registra errores encontrados en el prototipo.
    \item Aplica correcciones antes de la revisión final.
\end{itemize}

\subsubsection*{Desarrollador de interfaz 2 y control de calidad}
Carlos Sainz Arce

\begin{itemize}
    \item Implementa los módulos del portal financiero y ejecutivo.
    \item Ejecuta pruebas manuales de navegación y funcionalidad visual.
    \item Realiza pruebas de calidad del front-end junto con el desarrollador 1.
\end{itemize}

\subsubsection*{Interesados externos de validación}
Lic. Patricia Vargas Soto

\begin{itemize}
    \item Apoya la revisión del enfoque administrativo.
\end{itemize}

\paragraph*{Ing. Carlos Brenes Mora}

\begin{itemize}
    \item Apoya la revisión técnica general y la viabilidad visual del prototipo.
\end{itemize}

\paragraph*{Prof. Andrea Solís Zamora}

\begin{itemize}
    \item Apoya la validación del Portal Docente.
\end{itemize}

\paragraph*{Estudiantes}

\begin{itemize}
    \item Representan usuarios finales del Portal Estudiantil.
    \item Pueden apoyar en pruebas de usabilidad simuladas.
\end{itemize}

\paragraph*{Departamento Financiero}

\begin{itemize}
    \item Apoya la validación del Portal Financiero y los indicadores relacionados con pagos.
\end{itemize}

\paragraph*{Proveedor de Hosting}

\begin{itemize}
    \item Se considera como referencia externa para infraestructura, aunque no se realizará despliegue productivo.
\end{itemize}

\subsection{Criterios de aceptación}

El proyecto Plataforma Universitaria Integrada será aceptado cuando cumpla con los siguientes criterios:

\paragraph*{Alcance del prototipo}

\begin{itemize}
    \item El prototipo incluye Login, Dashboard principal, Portal Estudiantil, Portal Docente, Portal Administrativo, Portal Financiero y Dashboard Ejecutivo.
    \item Los módulos corresponden al alcance definido previamente.
\end{itemize}

\paragraph*{Navegación}

\begin{itemize}
    \item El usuario puede desplazarse entre el login, dashboard principal y módulos principales.
    \item Los botones o enlaces principales funcionan correctamente.
    \item No existen errores críticos que impidan navegar.
\end{itemize}

\paragraph*{Portal Estudiantil}

\begin{itemize}
    \item Se visualiza perfil del estudiante.
    \item Se representa matrícula de cursos.
    \item Se muestran horarios.
    \item Se muestran calificaciones.
    \item Se visualiza estado de cuenta.
    \item Se visualiza historial académico.
\end{itemize}

\paragraph*{Portal Docente}

\begin{itemize}
    \item Se visualizan cursos asignados.
    \item Se representa registro de asistencia.
    \item Se representa registro de calificaciones.
    \item Se muestra información básica de estudiantes por curso.
\end{itemize}

\paragraph*{Portal Administrativo}

\begin{itemize}
    \item Se visualiza gestión de estudiantes.
    \item Se visualiza gestión de docentes.
    \item Se visualiza gestión de cursos.
    \item Se visualizan períodos académicos.
\end{itemize}

\paragraph*{Portal Financiero}

\begin{itemize}
    \item Se visualizan pagos de matrícula.
    \item Se visualizan pagos de cursos.
    \item Se muestra historial de pagos.
    \item No se procesan pagos reales.
\end{itemize}

\paragraph*{Dashboard Ejecutivo}

\begin{itemize}
    \item Se visualizan indicadores académicos básicos.
    \item Se visualizan indicadores financieros básicos.
    \item Se representa información consolidada para la toma de decisiones.
    \item Los datos mostrados son simulados.
\end{itemize}

\paragraph*{Calidad visual y usabilidad}

\begin{itemize}
    \item La interfaz mantiene consistencia visual entre módulos.
    \item Los textos son legibles.
    \item Los menús son claros.
    \item La distribución de pantallas permite comprender el flujo general del sistema.
\end{itemize}

\paragraph*{Documentación del proyecto}

\begin{itemize}
    \item Los documentos mantienen una estricta coherencia con todos los entregables estratégicos de planificación y control previamente aprobados.
    \item Los entregables tienen nombres correctos.
    \item Las fechas, roles, presupuesto y alcance no se contradicen.
    \item El presupuesto total se mantiene en ₡17,625,200.
\end{itemize}

\paragraph*{Cronograma y evidencia}

\begin{itemize}
    \item El cronograma fue elaborado en Microsoft Project.
    \item El cronograma incluye tareas, subtareas, dependencias, duraciones, recursos, hitos, responsables, costos y ruta crítica.
    \item Se conservan evidencias en PDF del cronograma, recursos, ruta crítica y estadísticas del proyecto.
    \item La entrega se realiza según la fecha final definida en el cronograma aprobado.
\end{itemize}

\subsection{Puntos de verificación}

Durante el desarrollo del proyecto se realizarán puntos de verificación para revisar calidad antes de la entrega final.

\begin{itemize}
    \item \textbf{Revisión de inicio:} Se verifica que los documentos iniciales de planificación se encuentren perfectamente alineados. En esta revisión se comprueba que el problema, el propósito, el patrocinador, los roles, los interesados y los beneficios esperados guarden total coherencia entre sí.
    \item \textbf{Revisión de alcance:} Se valida que la Definición del Alcance indique claramente qué se incluye, qué se excluye, cuáles son los supuestos, restricciones, entregables y criterios generales de aceptación.
    \item \textbf{Revisión del cronograma:} Se revisa que el cronograma en Microsoft Project tenga tareas, dependencias, duración, recursos, hitos, ruta crítica, responsables y costos. También se verifica que la fecha final coincida con el cronograma aprobado y que el presupuesto total sea de ₡17,625,200.
    \item \textbf{Revisión de roles y organización:} Se verifica que el organigrama del proyecto mantenga una estructura clara con patrocinador, directora del proyecto, analistas de negocio y documentadores, y desarrolladores de interfaz y control de calidad. También se revisa que los interesados externos aparezcan exclusivamente como apoyo de validación.
    \item \textbf{Revisión del diseño del prototipo:} Se revisa que el prototipo tenga una estructura clara, pantallas comprensibles y navegación entre módulos. Esta revisión evita retrabajo antes de finalizar el desarrollo.
    \item \textbf{Revisión de la interfaz:} Se verifica que los módulos principales estén implementados adecuadamente a nivel visual y navegable. También se revisa que los datos mostrados sean simulados y que no se prometa ninguna funcionalidad que se encuentre fuera del alcance estipulado.
    \item \textbf{Pruebas manuales:} Se ejecuta un checklist de navegación y funcionalidad visual. Se registran errores, se corrigen los problemas críticos y se revalida el prototipo.
    \item \textbf{Revisión final de documentos:} Se revisa la coherencia entre todos los entregables, especialmente nombres, fechas, roles, presupuesto, alcance, criterios de aceptación y evidencia.
    \item \textbf{Validación de cierre:} Se confirma que el prototipo y los documentos estén listos para la entrega final. También se prepara la repositorio compartido en GitHub con documentos, evidencias y archivos correspondientes.
\end{itemize}

\subsection{Tareas y estimación de tiempo}

La estimación de tiempo del Plan de Calidad se basa en una técnica ascendente, donde se desglosan las actividades necesarias para asegurar y controlar la calidad del proyecto. Estas tareas están incluidas dentro del cronograma general del proyecto y no representan trabajo adicional fuera del alcance.

\paragraph*{Fase 1: Planificación de calidad}

\begin{table}[H]
\centering
\caption{Fase 1: planificación de calidad}
\label{tab:fase-1-planificacion-calidad}
{\small
\setlength{\tabcolsep}{3.5pt}
\begin{tabularx}{\textwidth}{Y p{3.2cm}}
\toprule
\textbf{Tarea} & \textbf{Tiempo estimado} \\
\midrule
Revisión del alcance y criterios generales de aceptación & 2 horas \\
Definición de objetivos de calidad & 2 horas \\
Definición de métricas de calidad & 2 horas \\
Selección de estándares aplicables & 2 horas \\
Documentación del Plan de Calidad & 4 horas \\
Total & 12 horas \\
\bottomrule
\end{tabularx}
}
\end{table}

\paragraph*{Fase 2: Aseguramiento documental}

\begin{table}[H]
\centering
\caption{Fase 2: aseguramiento documental}
\label{tab:fase-2-aseguramiento-documental}
{\small
\setlength{\tabcolsep}{3.5pt}
\begin{tabularx}{\textwidth}{Y p{3.2cm}}
\toprule
\textbf{Tarea} & \textbf{Tiempo estimado} \\
\midrule
Revisión de coherencia documental y de alcance & 2 horas \\
Revisión de coherencia entre interesados, organigrama y responsabilidades asignadas & 2 horas \\
Revisión de coherencia entre Alcance, Cronograma y Costos & 3 horas \\
Revisión de nombres, fechas y entregables & 2 horas \\
Ajustes menores de formato y redacción & 3 horas \\
Total & 12 horas \\
\bottomrule
\end{tabularx}
}
\end{table}

\paragraph*{Fase 3: aseguramiento del prototipo de interfaz}

\begin{table}[H]
\centering
\caption{Fase 3: aseguramiento del prototipo de interfaz}
\label{tab:fase-3-aseguramiento-prototipo}
{\small
\setlength{\tabcolsep}{3.5pt}
\begin{tabularx}{\textwidth}{Y p{3.2cm}}
\toprule
\textbf{Tarea} & \textbf{Tiempo estimado} \\
\midrule
Revisión de estructura del Dashboard principal & 2 horas \\
Revisión del Portal Estudiantil & 2 horas \\
Revisión del Portal Docente & 2 horas \\
Revisión del Portal Administrativo & 2 horas \\
Revisión del Portal Financiero & 2 horas \\
Revisión del Dashboard Ejecutivo & 2 horas \\
Total & 12 horas \\
\bottomrule
\end{tabularx}
}
\end{table}

\paragraph*{Fase 4: Control de calidad y pruebas}

\begin{table}[H]
\centering
\caption{Fase 4: control de calidad y pruebas}
\label{tab:fase-4-control-calidad-pruebas}
{\small
\setlength{\tabcolsep}{3.5pt}
\begin{tabularx}{\textwidth}{Y p{3.2cm}}
\toprule
\textbf{Tarea} & \textbf{Tiempo estimado} \\
\midrule
Ejecución de pruebas manuales de navegación & 4 horas \\
Pruebas de usabilidad básica & 3 horas \\
Pruebas de compatibilidad en navegador & 2 horas \\
Registro de defectos encontrados & 2 horas \\
Corrección de defectos críticos & 4 horas \\
Revalidación de módulos corregidos & 3 horas \\
Total & 18 horas \\
\bottomrule
\end{tabularx}
}
\end{table}

\paragraph*{Fase 5: Validación final y evidencia}

\begin{table}[H]
\centering
\caption{Fase 5: validación final y evidencia}
\label{tab:fase-5-validacion-final}
{\small
\setlength{\tabcolsep}{3.5pt}
\begin{tabularx}{\textwidth}{Y p{3.2cm}}
\toprule
\textbf{Tarea} & \textbf{Tiempo estimado} \\
\midrule
Revisión final del prototipo contra criterios de aceptación & 3 horas \\
Revisión final de documentos & 3 horas \\
Preparación de evidencias en PDF & 2 horas \\
Organización de carpeta final & 2 horas \\
Validación final del equipo & 2 horas \\
Total & 12 horas \\
\bottomrule
\end{tabularx}
}
\end{table}

\paragraph*{Resumen de tiempo de calidad}

\begin{table}[H]
\centering
\caption{Resumen de tiempo de calidad}
\label{tab:resumen-tiempo-calidad}
{\small
\setlength{\tabcolsep}{3.5pt}
\begin{tabularx}{\textwidth}{Y p{3.2cm}}
\toprule
\textbf{Tarea} & \textbf{Tiempo estimado} \\
\midrule
Planificación de calidad & 12 horas \\
Aseguramiento documental & 12 horas \\
Aseguramiento del prototipo de interfaz & 12 horas \\
Control de calidad y pruebas & 18 horas \\
Validación final y evidencia & 12 horas \\
Total & 66 horas \\
\bottomrule
\end{tabularx}
}
\end{table}



\newpage
% --- Entregable 12 ---
\section{Plan de comunicaciones}

Esta sección detalla el plan de comunicaciones del proyecto Plataforma Universitaria Integrada desarrollado para la Universidad Tecnológica La Mejor con la finalidad principal de definir estructuradamente cómo fluirá la información entre el equipo técnico junto con el patrocinador y todos los interesados relacionados. La comunicación representa un pilar fundamental dentro de esta iniciativa debido a que permite mantener a todas las partes debidamente alineadas reduciendo confusiones operativas y facilitando la coordinación de tareas para garantizar que la entrega final se realice puntualmente según el cronograma aprobado considerando además que el equipo reducido exige una alta eficiencia organizativa.

Para asegurar el éxito de este enfoque se ha determinado que el plan de comunicaciones guarde una relación transversal con los demás entregables estratégicos elaborados previamente asegurando de esta manera que las decisiones relevantes se documenten de forma congruente en todas las áreas de alcance o calidad y costos. En consecuencia el equipo utilizará canales oficiales complementarios donde un repositorio de control de versiones funcionará como almacenamiento principal del prototipo y las evidencias documentales mientras que un sistema de mensajería rápida servirá para la coordinación diaria dejando el correo electrónico exclusivamente para formalizar decisiones importantes con los interesados externos logrando así que cada persona reciba la información correcta en el momento oportuno sin pérdida de trazabilidad.

\subsection{Propósito del plan de comunicaciones}

El propósito integral de este plan consiste en establecer formalmente las reglas y canales junto con los responsables o frecuencias necesarias para compartir la información del proyecto respondiendo de antemano interrogantes clave para guiar al equipo.

\begin{itemize}
    \item Qué información específica debe comunicarse y a quién debe enviarse.
    \item Quién es el responsable directo de gestionar dicha comunicación.
    \item Con qué frecuencia y a través de qué canal oficial debe compartirse.
    \item Qué formato debe tener el mensaje y si debe quedar registrado como evidencia.
    \item Cómo se abordarán los riesgos o cambios y problemas detectados.
\end{itemize}

El plan establece un marco estructurado para garantizar que la comunicación general del proyecto se mantenga oportuna y trazable adaptando el nivel de formalidad según la relevancia de la información compartida dado que las coordinaciones diarias internas podrán efectuarse ágilmente por mensajería rápida mientras que cualquier decisión importante sobre el alcance o las validaciones finales deberá documentarse obligatoriamente mediante el repositorio oficial o bien por correo electrónico y actas de reunión.

\subsection{Objetivos del plan de comunicaciones}

Los objetivos principales de este plan orientados a garantizar el flujo correcto de la información se detallan en la siguiente lista

\begin{itemize}
    \item Garantizar que el equipo de trabajo mantenga una comunicación constante sobre los avances o dudas y riesgos junto con los entregables pendientes.
    \item Definir canales oficiales estructurados para evitar que la información importante se pierda en conversaciones informales o quede sin el debido respaldo documental.
    \item Mantener informados al patrocinador y a los interesados externos sobre los progresos relevantes del proyecto especialmente cuando se requiera su validación oportuna.
    \item Asegurar que todos los documentos se mantengan organizados en un espacio compartido utilizando un repositorio de control de versiones como almacenamiento central.
    \item Reducir el riesgo de inconsistencias entre documentos abordando de manera coherente temas como el alcance o los costos y fechas del proyecto.
    \item Establecer un mecanismo formal y estandarizado para comunicar cambios que puedan afectar directamente el tiempo o la calidad y los riesgos identificados.
\end{itemize}

\subsection{Interesados considerados}

Este plan considera los interesados internos y externos definidos en los entregables anteriores.

\subsubsection{Interesados internos}

\begin{table}[H]
\centering
\caption{Interesados internos}
\label{tab:stakeholders-internos}
{\small
\begin{tabularx}{\textwidth}{p{3.7cm} p{4.1cm} Y}
\toprule
\textbf{Nombre} & \textbf{Rol} & \textbf{Necesidad de comunicación} \\
\midrule
Tamara Robles Camacho & Directora del proyecto & Coordinar avances, revisar entregables, comunicar decisiones y dar seguimiento al cronograma \\
Isaac Jiménez Blanco & Analista de Negocio y Documentador 1 & Mantener trazabilidad documental, comunicar riesgos, cambios, calidad y entregables \\
Jorge Abarca Quiros & Analista de Negocio y Documentador 2 & Elaborar y revisar entregables documentales, apoyar la gestión de riesgos, costos y comunicaciones \\
Mauricio González Prendas & Desarrollador de interfaz 1 y control de calidad & Comunicar avance técnico de portales básicos, problemas del prototipo, pruebas y correcciones \\
Carlos Sainz Arce & Desarrollador de interfaz 2 y control de calidad & Comunicar avance técnico de portales financiero y ejecutivo, y pruebas de calidad \\
\bottomrule
\end{tabularx}
}
\end{table}

\subsubsection{Interesados externos}

\begin{table}[H]
\centering
\caption{Interesados externos}
\label{tab:stakeholders-externos}
{\small
\begin{tabularx}{\textwidth}{p{3.7cm} p{3.9cm} Y}
\toprule
\textbf{Nombre} & \textbf{Rol} & \textbf{Necesidad de comunicación} \\
\midrule
Dr. Roberto Mora Fallas & Patrocinador & Recibir avances generales, validar entregables principales y aprobar decisiones relevantes \\
Lic. Patricia Vargas Soto & Directora Administrativa & Validar necesidades administrativas representadas en el prototipo \\
Ing. Carlos Brenes Mora & Director de TI & Revisar aspectos técnicos generales y viabilidad visual del prototipo \\
Prof. Andrea Solís Zamora & Representante Docente & Validar claridad del Portal Docente \\
Estudiantes & Usuarios finales & Apoyar validaciones básicas del Portal Estudiantil si se requiere \\
Departamento Financiero & Área financiera & Validar representación del Portal Financiero \\
Proveedor de Hosting & Proveedor externo & Brindar referencia técnica si se requiere publicar o alojar evidencia del prototipo \\
\bottomrule
\end{tabularx}
}
\end{table}

\subsection{Canales oficiales de comunicación}

Para mantener orden en la comunicación del proyecto se definen los siguientes canales oficiales

\begin{itemize}
    \item \textbf{Repositorio de control de versiones:} Será el canal principal para centralizar el repositorio del prototipo de interfaz junto con documentos finales o evidencias y versiones controladas y también puede utilizarse para registrar solicitudes técnicas de cambios relacionados con el desarrollo.
    \item \textbf{Mensajería instantánea:} Será utilizado para comunicación rápida entre los integrantes del equipo coordinando horarios o resolviendo dudas rápidas y avisando sobre avances inmediatos sin ser considerado como un canal formal para aprobar cambios de alcance.
    \item \textbf{Plataforma de videollamadas:} Será utilizado para reuniones de seguimiento y revisión de entregables o validación de avances documentando obligatoriamente mediante una minuta o mensaje formal cuando una reunión produzca decisiones relevantes.
    \item \textbf{Correo electrónico:} Será utilizado para comunicaciones formales con el patrocinador o los interesados externos y también se utilizará cuando se requiera dejar evidencia de una aprobación o entrega oficial y solicitudes de revisión.
    \item \textbf{Microsoft Project:} Será utilizado para comunicar y respaldar la planificación del cronograma incluyendo tareas o fechas y responsables junto con hitos o dependencias de la ruta crítica comunicando siempre las actualizaciones relevantes al equipo.
    \item \textbf{Microsoft Word y PDF:} Serán utilizados para preparar y entregar documentos formales del proyecto utilizando los formatos PDF como evidencia final inalterable de los entregables.
\end{itemize}

\subsection{Reglas generales de comunicación}

Para evitar confusión y mantener trazabilidad, el equipo aplicará las siguientes reglas

\begin{itemize}
    \item Toda decisión que afecte alcance, tiempo, costo, calidad o riesgos deberá documentarse formalmente.
    \item La mensajería instantánea podrá usarse para coordinación rápida, pero no para aprobar cambios importantes.
    \item Los entregables finales deberán guardarse en el repositorio de control de versiones o en una carpeta compartida asociada.
    \item Cada documento deberá tener un nombre claro y relacionado con el entregable correspondiente.
    \item Las reuniones importantes deberán dejar una minuta breve o resumen de acuerdos.
    \item Los cambios al prototipo deberán registrarse en el repositorio de control de versiones mediante entregas de código o solicitudes.
    \item Los problemas técnicos o defectos relevantes deberán comunicarse al equipo y documentarse.
    \item La directora del proyecto será responsable de dar seguimiento general a la comunicación de la iniciativa.
    \item El Analista de Negocio y Documentador deberá revisar que los documentos mantengan coherencia entre sí.
    \item El equipo desarrollador de interfaz y control de calidad deberá comunicar los avances técnicos, los defectos encontrados y las correcciones correspondientes.
    \item Los interesados externos serán consultados únicamente cuando se requiera la validación general de procesos o módulos.
\end{itemize}

\subsection{Matriz de comunicaciones}

La siguiente matriz define qué información se comunicará, quién será responsable, a quién se dirigirá, con qué frecuencia, por qué canal y con qué nivel de prioridad.

\begin{widepage}
{\scriptsize
\setlength{\tabcolsep}{2.4pt}
\begin{longtable}{p{3.05cm} p{3.15cm} p{3.65cm} p{2.6cm} p{2.55cm} p{3.0cm} p{1.65cm}}
\caption{Matriz de comunicaciones}\label{tab:matriz-comunicaciones}\\[-0.5em]
\toprule
\textbf{Información a comunicar} & \textbf{Responsable} & \textbf{Destinatarios} & \textbf{Frecuencia} & \textbf{Canal} & \textbf{Formato} & \textbf{Prioridad} \\
\midrule
\endfirsthead
\toprule
\textbf{Información a comunicar} & \textbf{Responsable} & \textbf{Destinatarios} & \textbf{Frecuencia} & \textbf{Canal} & \textbf{Formato} & \textbf{Prioridad} \\
\midrule
\endhead
Avance general del proyecto & Directora del proyecto & Equipo y patrocinador & Semanal o según avance & Videollamada o correo & Resumen breve de avance & Alta \\
Estado de tareas del equipo & Directora del proyecto & Equipo interno & Cada 2 o 3 días & Mensajería & Mensaje de seguimiento & Media \\
Entregables documentales finalizados & Analista de Negocio y Documentador & Equipo interno & Por evento & Repositorio y mensajería & Documento formal o PDF & Alta \\
Versiones finales de documentos & Analista de Negocio y Documentador & Equipo interno y patrocinador si aplica & Por evento & Repositorio o correo & Archivo final & Alta \\
Avance del prototipo de interfaz & Desarrollador de interfaz y control de calidad & Equipo interno & Cada 2 o 3 días & Mensajería y repositorio & Mensaje y entrega de código & Alta \\
Cambios en el código del prototipo & Desarrollador de interfaz y control de calidad & Equipo interno & Por evento & Repositorio & Entrega de código o solicitud & Media \\
Defectos encontrados en el prototipo & Desarrollador de interfaz y control de calidad & Equipo interno & Por evento & Repositorio o mensajería & Solicitud o mensaje & Alta \\
Riesgos identificados & Analista de Negocio y Documentador & Directora del proyecto y equipo & Por evento & Mensajería y documento formal & Registro de riesgos & Alta \\
Cambios solicitados & Directora del proyecto o responsable & Equipo y patrocinador si aplica & Por evento & Correo, videollamada o documento & Solicitud de cambio & Alta \\
Evaluación de cambios & Analista de Negocio y Documentador & Directora del proyecto y equipo & Por evento & Documento formal y reunión & Evaluación de impacto & Alta \\
Validación de módulos & Desarrollador de interfaz y control de calidad & Equipo interno e interesados & Según avance & Videollamada o mensajería & Revisión visual o lista de control & Media \\
Validación del Portal Docente & Directora del proyecto & Prof. Andrea Solís Zamora y equipo & Por evento & Correo o videollamada & Solicitud de revisión & Media \\
Validación del Portal Financiero & Directora del proyecto & Departamento Financiero y equipo & Por evento & Correo o videollamada & Solicitud de revisión & Media \\
Validación administrativa & Directora del proyecto & Lic. Patricia Vargas Soto y equipo & Por evento & Correo o videollamada & Solicitud de revisión & Media \\
Revisión técnica general & Directora del proyecto & Ing. Carlos Brenes Mora y equipo & Por evento & Correo o videollamada & Solicitud de revisión & Media \\
Entrega final del proyecto & Directora del proyecto & Patrocinador, equipo y profesora & Al cierre & Repositorio y correo si aplica & Carpeta final y documentos PDF & Alta \\
\bottomrule
\end{longtable}
}
\end{widepage}

\subsection{Comunicación según prioridad}

Para manejar adecuadamente la información, se establecen tres niveles de prioridad.

\subsubsection{Prioridad alta}

Corresponde a información que puede afectar alcance, cronograma, costo, calidad, riesgos o entrega final. Debe comunicarse de inmediato al equipo y documentarse formalmente.

\noindent Ejemplos

\begin{itemize}
    \item Cambio en el alcance del prototipo.
    \item Atraso en una tarea crítica.
    \item Error crítico en la interfaz.
    \item Inconsistencia importante entre documentos.
    \item Modificación del presupuesto.
    \item Riesgo que pueda afectar la entrega final.
\end{itemize}

\noindent Canales recomendados

\begin{itemize}
    \item Mensajería instantánea para la alerta inicial.
    \item Videollamada para análisis si es necesario.
    \item Documento formal, repositorio o correo para dejar evidencia.
\end{itemize}

\subsubsection{Prioridad media}

Corresponde a información necesaria para coordinar trabajo, revisar avances o resolver dudas, pero que no afecta de forma inmediata el cumplimiento del proyecto.

\noindent Ejemplos

\begin{itemize}
    \item Revisión de un documento.
    \item Solicitud de apoyo en redacción.
    \item Corrección visual del prototipo.
    \item Ajuste menor en una tabla.
    \item Coordinación de reunión.
\end{itemize}

\noindent Canales recomendados

\begin{itemize}
    \item Mensajería instantánea.
    \item Repositorio si se relaciona con código o archivos.
    \item Videollamada si requiere discusión.
\end{itemize}

\subsubsection{Prioridad baja}

Corresponde a información operativa o administrativa que no afecta directamente el alcance ni la fecha de entrega.

\noindent Ejemplos

\begin{itemize}
    \item Confirmación de horario.
    \item Comentarios menores de formato.
    \item Organización de archivos.
    \item Recordatorios simples.
\end{itemize}

\noindent Canales recomendados

\begin{itemize}
    \item Mensajería instantánea.
    \item Mensaje breve al equipo.
\end{itemize}

\subsection{Gestión de reuniones}

Las reuniones del proyecto se realizarán exclusivamente cuando resulten indispensables para revisar avances o resolver dudas y validar entregables importantes dado que el cronograma debe mantenerse estrictamente controlado priorizando en todo momento sesiones cortas y enfocadas.

\begin{itemize}
    \item \textbf{Reunión de seguimiento interno:} El objetivo principal consiste en revisar el avance de las tareas junto con los problemas técnicos y los próximos pasos donde participará exclusivamente el equipo interno con una frecuencia de una o dos veces por semana según la necesidad utilizando la videollamada como canal principal bajo la responsabilidad directa de la directora del proyecto.
    \item \textbf{Reunión de revisión de entregables:} El propósito principal es revisar minuciosamente los documentos importantes antes de considerarlos finalizados requiriendo la participación del equipo interno por evento mediante videollamadas organizadas por la persona responsable de cada entregable particular.
    \item \textbf{Reunión de validación con interesados:} El objetivo central se enfoca en solicitar retroalimentación general sobre un módulo o documento integrando a la directora del proyecto junto con el responsable del módulo y el interesado correspondiente según la necesidad del momento utilizando videollamadas o correos formales bajo la responsabilidad de la directora.
    \item \textbf{Reunión de cierre:} El objetivo definitivo es revisar exhaustivamente que todos los entregables estén completamente terminados y listos para la entrega final con la participación del equipo interno reuniéndose una única vez al concluir el proyecto mediante videollamada bajo la responsabilidad directa de la directora.
    \item \textbf{Minuta de reunión:} Cuando una reunión genere acuerdos estratégicos se deberá documentar un resumen formal que detalle la fecha y los participantes asistentes junto con los temas tratados y los acuerdos principales especificando siempre los responsables asignados y la fecha de seguimiento para garantizar trazabilidad.
\end{itemize}

\subsection{Gestión de documentos y entregables}

La gestión de documentos representa una parte fundamental del plan integral debido a que el proyecto contempla múltiples entregables administrativos y técnicos que deberán organizarse meticulosamente dentro del repositorio utilizando una estructura clara de carpetas siguiendo rigurosamente las siguientes normativas documentales

\begin{itemize}
    \item Cada archivo debe tener siempre un nombre descriptivo completamente claro.
    \item Los documentos finales deberán exportarse y guardarse inalterables en formato PDF.
    \item Las versiones editables podrán mantenerse excepcionalmente en Word cuando resulte estrictamente necesario.
    \item El archivo principal de la planificación en Microsoft Project deberá conservarse íntegro en su formato nativo.
    \item Las evidencias visuales importantes deberán almacenarse ordenadamente en formato PDF o como imagen.
    \item Queda prohibido subir versiones duplicadas de archivos sin una identificación clara o justificada.
    \item Todos los documentos finales deberán revisarse minuciosamente antes de ser compartidos oficial o públicamente.
\end{itemize}

\subsection{Comunicación de riesgos, cambios y problemas}

Los riesgos o cambios y problemas detectados deben comunicarse con total claridad debido a que poseen el potencial de afectar directamente el cumplimiento exitoso del proyecto en sus diferentes dimensiones de alcance y tiempo.

\begin{itemize}
    \item \textbf{Comunicación de riesgos:} Cuando se identifique oportunamente un nuevo riesgo se deberá notificar de inmediato al equipo interno y registrarlo paralelamente en el documento oficial incluyendo la descripción clara del riesgo identificado y su causa posible junto con el impacto esperado sobre el cronograma y su probabilidad de ocurrencia designando a la persona responsable y la acción de respuesta mitigadora planeada. A modo de ejemplo si se detecta que el desarrollo de la interfaz podría atrasarse el responsable asignado deberá comunicarlo inmediatamente al equipo para que la directora del proyecto revise el impacto real en el cronograma permitiendo que el analista de negocio actualice el registro de riesgos de ser verdaderamente necesario.
    \item \textbf{Comunicación de cambios:} Cualquier cambio propuesto que afecte directamente el alcance o tiempo y costos del proyecto debe tratarse formalmente dado que bajo ninguna circunstancia se aprobarán modificaciones estratégicas utilizando únicamente canales informales por lo cual la comunicación oficial requerirá obligatoriamente el nombre completo de la persona solicitante y la descripción detallada de la propuesta junto con su justificación de negocio incluyendo además el análisis del impacto generado en el alcance o cronograma y costos así como la recomendación técnica y el registro de la decisión final aprobada por la dirección.
    \item \textbf{Comunicación de problemas:} Un problema materializado representa una situación ocurrida que podría afectar el flujo normal del proyecto requiriendo una comunicación urgente lo antes posible como en casos de problemas para abrir el archivo centralizado de Microsoft Project o la presencia de un error crítico en el prototipo interactivo requiriendo siempre que los problemas de prioridad alta se notifiquen por la mensajería rápida para asegurar una reacción pronta y se documenten luego minuciosamente en el repositorio o mediante correo formal.
\end{itemize}

\subsection{Responsables de comunicación}

La comunicación del proyecto tendrá responsables claramente definidos asegurando que el flujo de información sea directo y efectivo en todo momento.

\begin{itemize}
    \item \textbf{Directora del proyecto (Tamara Robles Camacho):} Será la persona responsable de coordinar la comunicación general del proyecto organizando las reuniones de seguimiento y comunicando los avances generales al patrocinador además de validar constantemente que todo el equipo esté debidamente informado sobre las decisiones relevantes.
    \item \textbf{Analista de negocio y documentador 1 (Isaac Jiménez Blanco):} Tendrá el deber de mantener la coherencia absoluta entre todos los documentos comunicando los avances de la primera mitad de entregables documentales y registrando minuciosamente los riesgos o cambios y evaluaciones para verificar que la información siempre se mantenga trazable.
    \item \textbf{Analista de negocio y documentador 2 (Jorge Abarca Quiros):} Deberá comunicar de forma proactiva los avances de la segunda mitad de entregables documentales apoyando activamente la gestión comunicativa de los riesgos o costos y colaborando en el mantenimiento integral de la documentación del proyecto.
    \item \textbf{Desarrollador de interfaz 1 y control de calidad (Mauricio González Prendas):} Se encargará de comunicar los avances técnicos del prototipo relacionados con los módulos básicos registrando todos los cambios relevantes en el repositorio oficial y reportando los defectos encontrados junto con las correcciones realizadas o problemas técnicos que puedan afectar el cronograma.
    \item \textbf{Desarrollador de interfaz 2 y control de calidad (Carlos Sainz Arce):} Estará encargado de comunicar todos los avances técnicos del prototipo para los módulos financiero y ejecutivo reportando puntualmente los defectos de calidad o las correcciones efectuadas en su sección y colaborando intensamente en las pruebas de calidad técnica del desarrollo front-end.
    \item \textbf{Patrocinador (Dr. Roberto Mora Fallas):} Será la persona encargada de recibir la información sobre el avance general y revisar los entregables principales aprobando finalmente las decisiones más relevantes cuando estas afecten el propósito original del proyecto validando de igual manera el cierre definitivo del desarrollo.
\end{itemize}

\subsubsection{Interesados externos}

\begin{itemize}
    \item Brindar retroalimentación general cuando se les consulte.
    \item Validar módulos relacionados con su área.
    \item Comunicar observaciones que mejoran la representación del prototipo.
\end{itemize}

\subsection{Seguimiento y control de las comunicaciones}

El seguimiento de las comunicaciones permite verificar que la información fluya correctamente durante el proyecto. Para ello, se aplicarán las siguientes acciones

\begin{itemize}
    \item Revisar periódicamente que los entregables estén en la ubicación compartida definida.
    \item Confirmar que las decisiones importantes queden documentadas.
    \item Verificar que los integrantes del equipo conozcan sus tareas y fechas.
    \item Revisar que los cambios importantes pasen por el proceso formal correspondiente.
    \item Confirmar que los riesgos relevantes sean comunicados y registrados.
    \item Validar que las versiones finales estén correctamente nombradas.
    \item Revisar que los acuerdos de reunión sean entendidos por los participantes.
\end{itemize}

\subsubsection{Indicadores simples de seguimiento}

\begin{table}[H]
\centering
\caption{Indicadores simples de seguimiento}
\label{tab:indicadores-seguimiento}
\begin{tabularx}{\textwidth}{Y p{3cm}}
\toprule
\textbf{Indicador} & \textbf{Meta} \\
\midrule
Entregables importantes almacenados en el repositorio & 100\% \\
Cambios relevantes documentados formalmente & 100\% \\
Riesgos importantes comunicados al equipo & 100\% \\
Reuniones con acuerdos documentados cuando aplique & 100\% \\
Documentos finales con nombre correcto & 100\% \\
Integrantes informados sobre tareas críticas & 100\% \\
\bottomrule
\end{tabularx}
\end{table}

Si durante el seguimiento se detecta que un canal no está funcionando, la directora del proyecto deberá proponer un ajuste. Por ejemplo, si la mensajería instantánea genera pérdida de información, las decisiones deberán trasladarse a correo, al repositorio o a una minuta de reunión.

\newpage
% --- Entregable 13 ---
\section{Plan de gestión de riesgos}

Esta sección explica el plan de gestión de riesgos y define la metodología que el equipo utilizará para identificar, analizar, priorizar, responder y monitorear los contratiempos latentes durante el ciclo de vida del proyecto. Las subsecciones correspondientes establecen de manera formal las reglas y procesos que sustentan la creación del registro de riesgos así como la configuración y evaluación de la matriz integral de probabilidades e impacto.

La gestión de riesgos resulta un proceso fundamental para lograr una anticipación oportuna a posibles atrasos e inconsistencias entre los documentos aprobados. Esta labor preventiva mitiga simultáneamente la sobrecarga imprevista de trabajo en el equipo al mismo tiempo que resuelve de antemano los problemas técnicos emergentes y derivados del proceso directo del desarrollo.

\subsection{Propósito}

El propósito fundamental de este plan comprende el diseño estructurado y la respectiva aplicación metódica de los objetivos descritos a lo largo de la siguiente lista de puntos.

\begin{itemize}
    \item Establecer una metodología clara y aplicable que permita identificar eficazmente los riesgos dentro de todas las fases constructivas del proyecto.
    \item Definir las escalas oficiales de probabilidad e impacto asegurando su utilización consistente en cada etapa de evaluación del equipo.
    \item Establecer la fórmula matemática exacta y los diferentes rangos paramétricos requeridos para lograr clasificar apropiadamente el nivel de gravedad inherente a cada eventualidad.
    \item Definir un catálogo de estrategias preventivas y correctivas disponibles para que el equipo pueda abordar y neutralizar correctamente cualquier adversidad o contingencia surgida.
    \item Asignar roles logísticos y responsabilidades operativas en beneficio de la gestión de riesgos para lograr un manejo siempre coordinado y totalmente transparente.
    \item Establecer la frecuencia de revisión idónea al igual que el mecanismo central y estructurado necesario para consolidar el seguimiento integral e ininterrumpido de todos los hallazgos documentados.
\end{itemize}

\subsection{Roles y responsabilidades}

La gestión de riesgos representa una responsabilidad compartida por todo el equipo interno y se han asignado roles debidamente diferenciados según el grado de injerencia técnica o administrativa evaluada.

\begin{widepage}
{\footnotesize
\setlength{\tabcolsep}{3pt}
\begin{longtable}{P{4.5cm} P{9.3cm} P{9.3cm}}
\caption{Roles y responsabilidades en gestión de riesgos}\label{tab:roles-responsabilidades} \\
\toprule
\textbf{Rol} & \textbf{Responsabilidad en gestión de riesgos} & \textbf{Categorías de riesgo asignadas} \\
\midrule
\endfirsthead
\caption[]{Roles y responsabilidades en gestión de riesgos (continuación)} \\
\toprule
\textbf{Rol} & \textbf{Responsabilidad en gestión de riesgos} & \textbf{Categorías de riesgo asignadas} \\
\midrule
\endhead
\bottomrule
\endfoot
\bottomrule
\endlastfoot
Directora del proyecto (Tamara Robles Camacho) & Lidera integralmente el proceso de identificación y análisis de contingencias para posteriormente aprobar formalmente el documento base priorizando la ejecución de cualquier acción preventiva. & Mantiene un seguimiento riguroso sobre eventos que amenacen el cumplimiento del cronograma global de entregas y vigila la disponibilidad de todo el personal protegiendo la constante y fluida comunicación ante las variaciones dictaminadas por todos los interesados externos. \\
Analista de negocio y documentador 1 (Isaac Jiménez Blanco) & Trabaja en la certera y oportuna identificación y respectiva descripción sistemática y valorativa relacionada a las vulnerabilidades incidentes detectadas en las mediciones periódicas realizadas sobre los parámetros fijados del alcance. & Mantiene una vigilancia formal constante logrando atender puntualmente todas aquellas situaciones de falla derivadas al aplicar rigurosamente una gestión documental evaluando continuamente cualquier tipo de falencia orientada firmemente a perjudicar la debida completitud y solidez inquebrantable establecida sobre el presupuesto. \\
Analista de negocio y documentador 2 (Jorge Abarca Quiros) & Colabora e identifica documentando sistemáticamente aquellos fallos colaterales fuertemente ligados de manera intrínseca a fallas recurrentes de los procesos generales estipulados de revisión integral dictaminando con plena evidencia probatoria y exhaustiva su nivel adverso o desfavorable. & Realiza comprobaciones constantes prestando un valioso y adecuado servicio correctivo de seguimiento frente a cada desvío detectado originado tras presentarse faltas significativas logrando solventar deficiencias evidenciadas claramente tras vulneraciones severas encontradas de manera preocupante en revisiones documentales periódicas. \\
Desarrollador de interfaz y QA (Mauricio González Prendas) & Identifica oportunamente posibles problemas críticos detectados en los despliegues visuales aplicando la logística orientada al desarrollo técnico y efectúa oportunas medidas de mitigación requeridas invariablemente durante la codificación y estructuración. & Supervisa sistemática e ininterrumpidamente las interacciones logísticas que garantizan resguardar eficientemente el estado ideal concebido por los prototipos e interviene diligentemente y proactivamente implementando su control general incesante ante falencias de estructura en todo el historial originado del seguimiento respectivo referente y proveniente del avance técnico general. \\
Desarrollador de interfaz y QA (Carlos Sainz Arce) & Efectúa de manera cooperativa y oportuna revisiones minuciosas con los integrantes de logística asumiendo un análisis detallado buscando anticipar y detectar amenazas inminentes al transcurso propio e indispensable enfocado totalmente hacia una oportuna implementación en todas las estructuras e interacciones del flujo informático. & Centra e invierte de modo eficiente sus primordiales labores invirtiendo un nivel idóneo de monitoreo continuo orientado sólidamente a mitigar cualquier incidente surgido directamente tras detectarse falencias inmersas intrínsecamente y provenientes del constante y regular y normal control efectuado al depositar su revisión sistemática sobre los registros corporativos del repositorio e integra medidas salvaguardando en paralelo todo tipo de integridad operativa orientada a lograr el acatamiento riguroso de cada versión del proyecto. \\
\end{longtable}
}
\end{widepage}

\subsection{Ciclo de gestión de riesgos (PMBOK)}

El proyecto aplica e integra el ciclo de gestión de riesgos establecido y validado metodológicamente por el PMBOK y adapta minuciosamente sus respectivos principios fundamentales a la escala específica del entorno considerando las proporciones del desarrollo de la plataforma universitaria integradora.

\begin{widepage}
{\footnotesize
\setlength{\tabcolsep}{3pt}
\begin{longtable}{P{4.5cm} P{19cm}}
\caption{Ciclo de gestión de riesgos}\label{tab:ciclo-riesgos} \\
\toprule
\textbf{Etapa} & \textbf{Descripción y aplicación en el proyecto} \\
\midrule
\endfirsthead
\caption[]{Ciclo de gestión de riesgos (continuación)} \\
\toprule
\textbf{Etapa} & \textbf{Descripción y aplicación en el proyecto} \\
\midrule
\endhead
\bottomrule
\endfoot
\bottomrule
\endlastfoot
\textbf{1. Planificación de la gestión de riesgos} & Estructura sólidamente los procedimientos delimitados al sentar integralmente la definición logística dictando claramente los principios aplicables referentes a la metodología establecida e implementada en la determinación paralela referida al uso y disposición de escalas valorativas definiendo además firmemente el rol general asignado en las herramientas base. \\
\textbf{2. Identificación de riesgos} & Fomenta la participación mediante reuniones dinámicas del equipo para realizar la extensa revisión general fundamentada en la extracción exhaustiva de los hallazgos enmarcada bajo lineamientos dictaminados operativamente por el proyecto en fuentes documentadas fiables incluyendo directamente los recursos consolidados pertenecientes ineludiblemente al acta de constitución del proyecto al lado y apoyado invariablemente en la formulación de los esquemas estructurales estipulados minuciosamente para determinar cada entrega formal del plan en base. \\
\textbf{3. Análisis cualitativo} & Pondera en un cálculo numérico el nivel respectivo estipulado del peligro para el proyecto evaluando objetivamente un grado calculado logrando determinar metódicamente un cruce o correlación estable mediante la escala general preestablecida numerada uniformemente e identificando la gravedad respectiva consolidada categorizando y resolviendo su ubicación definitiva dentro de esquemas tipificados de grado muy bajo a crítico. \\
\textbf{4. Análisis cuantitativo} & Aplica este tipo y modo especifico de riguroso análisis detallado destinándolo en especial al rango que agrupa a todos y cada uno de los más letales problemas inminentes y contingencias detectadas que ostentan e integran definitivamente una nivelación evaluada dictaminando la inminente afectación operando e influyendo destructivamente el valor monetario esperado calculado con total rigidez tras multiplicar matemáticamente el factor respectivo evaluado equivalente a un número base sobre su esperada probabilidad multiplicándola directamente y sin demora contra el margen evaluado del grado de un impacto monetario. \\
\textbf{5. Planificación de respuestas} & Formaliza y diseña integralmente una robusta y preventiva estrategia enfocada en consolidar acciones defensivas orientadas a mitigar daños colaterales transfiriendo un nivel adecuado del contratiempo y estableciendo formal y meticulosamente sólidas o rigurosas vías estipuladas con el fin irrevocable enfocado de aplicar verdaderas barreras resolutivas acatando de pleno las contundentes medidas de choque orientadas operativamente en cada tipo abordado correspondiente. \\
\textbf{6. Monitoreo y control} & Implementa eficientemente un sistema ininterrumpido centrado puntualmente en aplicar revisiones programadas enfocadas enteramente a revisar incesantemente las posibles fluctuaciones del comportamiento reportado dentro de las respectivas calificaciones tabuladas y alojadas de los informes correspondientes permitiendo garantizar contundentemente la difusión ágil lograda en trasladar informes oportunos emitiendo una importante información alertando a todo su entorno de trabajo. \\
\end{longtable}
}
\end{widepage}

\subsection{Metodología de identificación de riesgos}

El abordaje y la identificación de riesgos se ejecuta aplicando y estructurando un proceso segmentado en cuatro niveles de enfoque y fundamentado sustancialmente tras completar revisiones exhaustivas y detalladas provenientes de las inspecciones que indagan minuciosamente una amplia base o cúmulo ineludible de fuentes y de los distintos insumos del grupo promoviendo al involucramiento equitativo apoyándose plenamente en valiosas técnicas reconocidas.

\subsubsection{Revisión de la documentación del proyecto}

El equipo procede integralmente a evaluar minuciosamente cada nivel o entrega de la respectiva planificación desarrollada e impulsada con anterioridad bajo el claro y contundente propósito logístico de identificar eventuales fuentes adversas generadoras de inminentes fallos prestando obligatoria atención y vigilancia sobre el respectivo espectro conformado de factores medulares integrados lógicamente a lo largo de un listado detallado en el próximo bloque de directrices.

\begin{itemize}
    \item \textbf{Alcance del proyecto:} El desarrollo presenta lamentables ambigüedades e importantes faltas lógicas de una esperada claridad relativas a las determinaciones técnicas donde se delimitan rigurosamente aquellos lineamientos referidos al respectivo diseño operativo y creación general o modelado estricto referidos al despliegue total y el entorno estructurado del desarrollo de interfaz.
    \item \textbf{Cronograma:} Presenta una base vulnerable al exhibir periodos muy ajustados forzando de tal manera un desarrollo saturado logrando afectar sustancialmente aquellas ineludibles obligaciones y tareas críticas estipuladas ocasionando inevitablemente concentrar en exceso un rango drástico o el grado integral de dependencia depositada irresponsablemente u omitiendo el debido resguardo hacia el apoyo estipulado originado en las labores de únicamente un solo desarrollador central.
    \item \textbf{Presupuesto:} Contiene factores monetarios cuestionables con previsiones imprecisas estimando variables o rubros ciertamente poco realistas evidenciando importantes y notables deficiencias estructurales resultando ser totalmente inconsistentes al ser cruzadas frente y contra diversos e ineludibles componentes base reflejados indiscutiblemente en el alcance y la planeación documentada en archivos anexos e indispensables de gran peso para el proyecto.
    \item \textbf{Requisitos:} Incluyen entre sus parámetros integrados múltiples funcionalidades complejas y en esencia muy imprecisas llegando al peligroso punto crítico donde varias proyecciones establecidas sin firmeza lamentablemente extienden y logran vulnerar todos y cada uno de los más seguros de los márgenes previstos comprometiendo el verdadero nivel del tiempo exigido superando de un modo imprudente y forzado todo objetivo aprobado enmarcando su respectivo desarrollo o cumplimiento de manera inviable e insegura en contra del margen predeterminado al comienzo estipulado e indiscutiblemente fijado para ser entregado a nivel del alcance fundamental.
    \item \textbf{Plan de comunicación:} Existe verdaderamente un riesgo real e inherente y constante fomentando operativamente posibles desacuerdos o probables diferencias perjudiciales y malentendidos humanos afectando peligrosamente la logística entre los cinco integrantes originando el decaimiento de las relaciones sanas en el ambiente.
\end{itemize}

\subsubsection{Involucramiento del equipo}

Esta etapa operativa establece los preceptos indispensables asegurando categóricamente el acatamiento irrestricto de todas aquellas obligaciones vinculantes ordenando y encausando formalmente un mecanismo estricto donde toda base dependiente o proveniente del grupo garantice indiscutiblemente implementar ininterrumpidamente una constante u oportuna y totalmente sólida interacción logística propiciando sin fallar la debida diligencia de identificar los sucesos reportándolos de modo diligente aplicando eficientes rutinas formalmente definidas en el próximo compilado normativo logístico.

\begin{itemize}
    \item Efectuar sesiones rápidas y concisas aplicando dinámicas productivas logradas por medio del valioso método sistemático referido comúnmente hacia el despliegue efectivo e irremplazable efectuando o realizando e incluso adoptando una valiosa herramienta de lluvia de ideas con las visiones entrelazadas y consolidadas de todos integrando firmemente a los cinco participantes asimilando una amplia o integradora capacidad conjunta en la respectiva perspectiva abarcando íntegramente las variables referentes hacia y enfocadas contundentemente por su naturaleza al lado netamente orientado u originado a labores de índole técnica inmersas en el desarrollo y abarcando desde luego toda visión enfocada minuciosamente hacia aquellas contingencias nacidas en el riguroso o vital proceso respectivo que rige la consolidación e inclusión incesante hacia los estándares referidos a la calidad documental paralela e imprescindible.
    \item Participar decididamente respetando ineludiblemente una cita de coordinación con reuniones semanales e indispensables estipuladas orientadas a un control en torno al firme y permanente propósito formal destinado expresamente en ventilar o divulgar ágilmente todo el cúmulo constante conformado ante preocupaciones emergentes detectadas comunicando pertinentemente todos aquellos diversos riesgos potenciales evidenciables analizando a fondo sus respectivas soluciones previendo de un modo estructurado o analítico todo desvío detectado tempranamente con premura mitigando los sucesos oportunamente originados al margen de la rutina establecida sin afectar operativamente al progreso corporativo resguardando toda estabilidad esperada por el progreso y avance tecnológico previsto ordenadamente en plazos de semanas logísticas.
\end{itemize}

\subsubsection{Técnicas de identificación utilizadas}

La investigación minuciosa e identificación puntual orientada y enfocada a detallar exhaustivamente la totalidad íntegra del catálogo abarcando en pleno a todos los riesgos inminentes detectados para el proyecto procediendo con tal rigurosidad o con el gran objetivo y firme directriz estipulada referida operativamente en asegurar un seguimiento metodológico estructurado que garantice efectivamente la utilización exhaustiva e indudable orientando inquebrantablemente su accionar implementando plenamente y adoptando integralmente un conjunto preciso fundamentado con técnicas logísticas eficientes explicadas seguidamente bajo los apartados de revisión respectiva o directriz descrita en modo explicativo subsecuentemente y de forma inmediata abarcada y detallada en esta sección.

\begin{itemize}
    \item Revisión sistemática sobre los parámetros dictados y provenientes del indispensable formato analítico implementando rigurosamente a un alto nivel estructural el proceso exhaustivo enfocado al uso metódico del formato integrado conocido en logística operativa como una estructura de desglose del trabajo (WBS) sometiendo su integridad evaluándola pormenorizadamente logrando observar a la vez identificando de modo eficiente cada variable interna frágil desglosando cualquier paquete y tarea estipulada verificando de este modo certeramente si esta llegase infortunadamente a fallar operativamente e inevitablemente logrando perturbar o impidiendo de una manera innegable el flujo esperado perjudicando fuertemente e irremediablemente y finalmente obstaculizando su respectiva consolidación impidiéndose a su vez concretarla satisfactoriamente afectando con ello perjudicialmente o retrasando contundentemente e impidiéndola de ejecutar adecuadamente durante la esperada fase temporal previamente aprobada con plazos innegociables e inalterables dictados por cronogramas institucionales y corporativos en el desarrollo respectivo.
    \item Análisis reflexivo extraído del estudio responsable originado sistemáticamente fundamentándose sólidamente en repasar y examinar con sumo cuidado todo registro correspondiente e invaluable agrupando rigurosamente todos y cada uno de los apuntes o datos referidos y relacionados a la aplicación valiosa del uso provechoso en base referida al importante rubro y la herramienta correspondiente de repasar los datos provenientes del inventario que consolida un documento de todas o al menos las indispensables lecciones documentadas concebidas o asimiladas tras afrontar una amplia lista sumada acumulando de manera empírica e invaluable una útil colección constituida del aprendizaje experimentado en diversos proyectos académicos anteriores o referida plenamente a una extensa revisión exhaustiva lograda y desarrollada de cursos técnicos aprobados evaluando con un gran nivel la amplia variedad de bases de datos.
    \item Análisis exhaustivo sobre un conglomerado de factores identificables enfocados netamente referidos del uso y estudio minucioso comprobatorio del modelo basado fundamentado a partir y originado basando su base comprobando un sistema de supuestos operando metódicamente logrando y efectuando ineludiblemente una verificación rigurosa asegurando y corroborando y sustentando firmemente con pruebas sólidas toda validación o refutación evaluada metódicamente que determine firmemente operativamente su carácter analítico refrendando indiscutiblemente e irrenunciablemente toda integridad referida y contenida o englobada estructuralmente inmersa como pilar en la creación base integradora concebida referida firmemente inmersa íntegramente alojada estructuralmente como una gran premisa formal concebida como fundamental o verdaderamente medular y pilar dentro de un acta de constitución del respectivo proyecto validándola integralmente ratificando de un modo indiscutible comprobando contundentemente la veracidad e implementando su efectiva certidumbre y sólida fiabilidad en vigencia.
    \item Clasificación organizada agrupando los factores originados por eventos agrupando analíticamente y encasillando las diversas vulnerabilidades ordenadas basándose según o determinando su principal y certera fuente de riesgo englobando de este modo todo entorno estructurado referenciando aquellos de origen netamente del sector del rubro referenciado como técnicos o englobando problemas de la administración logrando e incorporando del modo adecuado evaluando incidentes y problemas logísticos afectando los referidos plenamente originados afectando directamente al capital operativo en recursos humanos incrustando una amplia variedad documentada contemplando de la misma forma un marco y conjunto evaluando a todo registro abarcando fallas o inconvenientes causados originados por vulnerabilidades propias y fallos presentes afectando irremediablemente o causando disfunciones severas concernientes inherentes provenientes respecto al soporte y deficiencias evidentes del uso ineficiente referidas ineludible e indudablemente del uso de ineficientes o problemáticas herramientas así como todos los aspectos concernientes originados y englobando verdaderamente inmersos plenamente pertenecientes indiscutible y propiamente a los negocios.
\end{itemize}

\subsubsection{Fuentes utilizadas para este proyecto}

El equipo técnico aplica metódicamente recursos integrados enfocando los resultados resultantes obtenidos de efectuar y consolidar búsquedas fundamentando rígidamente su origen apoyándose siempre en el riguroso accionar logístico referido a inspeccionar y escudriñar de modo permanente y eficiente orientando plenamente este esquema procedimental enfocado rígidamente verificando un valioso listado de elementos documentales enumerados formalmente al estructurar su avance detallándolo fielmente bajo una perspectiva logística englobando y contemplando lo anotado documentando con certidumbre todos los incisos descritos minuciosamente a lo largo de este próximo apartado textual detallado seguidamente.

\begin{itemize}
    \item Revisión documental que involucra directamente componentes centrales tales como el documento del acta de constitución del proyecto o el archivo correspondiente y denominado propiamente como el respectivo caso de negocio al lado indudable e indispensable proveniente e ineludiblemente vinculado o ligado referenciando e incluyendo al importante y exhaustivo marco abarcado y proveniente de una fuente obligada revisando íntegramente la indispensable base de la cual emana del valioso recurso que engloba un debido y muy estricto registro integral en donde documentan a los interesados principales a la par de asegurar de verificar debidamente y de inspeccionar con sumo grado pormenorizadamente lo referido o respectivo al consolidado y muy detallado archivo e informe sobre la respectiva o contundente y sumamente precisa descripción estipulada referida originando de una innegociable e inquebrantable definición fijando el respectivo alcance.
    \item Control derivado proveniente referenciando contundentemente originado desde de la evaluación periódica sometiendo y evaluando o repasando a detalle meticuloso revisando las fechas enmarcadas integral y oficialmente mediante uso comprobado de estatus provenientes extraídos de revisar cronogramas e informes y bases estipuladas con formatos o del apoyo oficial de la debida fuente base referida originándose desde el modelo del cronograma general estructurado metódicamente e implementando un sólido seguimiento operativo desplegado rígidamente e ininterrumpidamente desarrollado y soportado íntegramente mediante un programa o sistema originando información base utilizando la respectiva y avalada herramienta institucional estipulada como estándar en Microsoft Project integrando todo esquema al igual o junto al apoyo referenciado abarcando e integrando lo correspondiente en revisar un debido y pertinente archivo estipulado conteniendo el valioso recurso originado del modelo o gráfico referenciado logísticamente originado o conteniendo y describiendo minuciosamente el respectivo nivel u organización administrativa plasmada integralmente e ilustrativamente englobando y detallando todo nivel plasmado y fundamentado o desplegado gráficamente y referenciado en el vital e indispensable organigrama concebido para plasmar fielmente todo el panorama y estructura concerniente a detallar un organigrama que visualiza y mapea orgánicamente documentando un modelo del proyecto.
    \item Documentación central concebida operativamente en la fase de análisis estructurando revisiones e implementando la logística evaluando con rigidez de una manera objetiva a partir fundamentalmente tras evaluar el respectivo reporte logístico contenido integrando todo archivo oficial del proyecto enfocando de igual manera contundentemente un respectivo esfuerzo comprobatorio riguroso y en paralelo logísticamente e idénticamente de un modo u otro ejecutando la respectiva verificación exhaustiva integrando logísticamente todo recurso estipulado proveniente y contenido e inmerso fundamentalmente del importante formato logístico o plan de calidad logrando en base lograr asegurar plenamente en paralelo efectuar operativamente de un modo exhaustivo todo examen pertinente implementando la revisión a la vez abarcando todos o integrando rigurosamente cada formato de la logística documentada e inmersa referida en revisar los lineamientos y protocolos respectivos referenciando minuciosamente todo nivel respectivo estipulado del formato referenciado documentando todo nivel originado en un debido formato originado en revisar a fondo un respectivo o determinado archivo estipulado referido operativamente abarcando integralmente lo referente al plan logístico y enfocado y referido estipulando en revisar o avalar estipulando normas referenciando todo un documento de las referidas pautas establecidas e inmersas referidas referenciando o documentando formalmente e inmersamente concerniente referidas operativamente documentadas concerniente e inmersamente en todo el espectro abarcado o abordado logísticamente en un indispensable y estructurado plan respectivo referenciado originando abarcando integral y operativamente un documento sobre un plan con las debidas y precisas normas respectivas estipuladas en el respectivo plan documentado de las respectivas e importantes comunicaciones.
    \item Análisis meticuloso considerando todas las inquebrantables restricciones limitadas del tiempo aunado fuertemente a evidenciar todo informe documentando minuciosamente las respectivas estipulaciones o directrices organizativas refiriéndose y abarcando y describiendo el diseño o modelo abordando operativamente a una firme organización con la debida directriz y lineamiento en un firme sistema referido operativamente e incrustado orgánicamente refiriéndose operativamente y plasmando y englobando y abarcando o evidenciando la detallada e idónea y respectiva logística estipulando la debida estructura e integral distribución corporativa documentando y detallando cada respectivo rol implementado íntegramente y estructurado minuciosamente con el debido detalle referido a englobar analíticamente estructurando de manera equitativa e indispensable lo referente sobre la oficial y precisa e integral y formal delimitación de funciones o delimitación estructurada estableciendo la respectiva y oportuna organización estipulando delimitadamente la muy idónea conformación estructurando u orientando y englobando y fundamentando organizativamente a la estricta respectiva distribución logrando encausar ordenando organizativamente abarcando de un modo eficiente referenciando e incluyendo el rol específico de cada uno abordando las respectivas directrices y atribuciones delimitadas y asignadas u orientadas para ser debidamente acatadas responsable e íntegramente de un modo ordenado referenciando y aludiendo rigurosamente o englobando con una alta certeza respectiva o refiriéndose rigurosamente y debidamente sobre el marco documentado estipulando o rigiendo e indicando la formal delimitación de labores ordenando o englobando a la estricta asignación originando referida a lograr establecer de lleno estructurando organizativamente e inmersamente delimitando organizando documentando e implementando la respectiva organización estipulando organizando e implementando equitativamente a la adecuada distribución o enmarcada o correspondiente respecto a una sana organización en una distribución inmersamente concerniente al personal referenciando de modo respectivo y enfocado englobando en la logística y estipulando a cada uno abordando en una lista respectiva detallando responsabilidades de cada quien abarcando a cada respectivo participante abarcando organizadamente a los integrantes.
    \item Evaluaciones derivadas logradas inmersamente verificando e inspeccionando a través referenciado o enfocando toda perspectiva basándose del respectivo y frecuente uso concerniente analizando las correspondientes utilidades o funciones referidas a verificar todos y aquellos respectivos resultados de la implementación logística refiriéndose a las aplicaciones originadas derivadas de examinar los fallos ocasionados empleando de lleno la totalidad respectiva englobando e incorporando a las diferentes y a la vasta serie e inclusive englobando al listado abordando a aquellas o respectando logísticamente o referenciando indiscutible e innegablemente o en referenciando contundentemente o concerniente a englobar a toda gama o de toda índole abarcando en base refiriendo e englobando o aludiendo indiscutiblemente logísticamente concerniente abordando referenciando o inmersamente ligadas al uso de integraciones abarcadas o enfocadas e implementadas utilizando y operando y recurriendo concernientemente o respectando en abarcar u orientadas y abarcando a la gama operando la lista de las fundamentales herramientas y de plataformas u orientadas innegablemente referenciando al uso de sistemas enfocados orientadas a la logística corporativa como las referidas a documentar el almacenamiento tecnológico usando al respectivo repositorio así referenciando del mismo modo e incluyendo de modo concerniente implementando verificaciones abarcadas del mismo modo recurriendo abordando al control originado referenciando logísticamente del Microsoft Project y englobando lo respectivo orientando el uso y analizando de un procesador y abarcando orientando al archivo del documento orientando el documento respectivo englobando o referido abarcando y orientando a aquellos del tipo de Word anexando el contenido referenciado a archivos abarcados de tipo o referenciado a un respectivo o concerniente y enfocado tipo referido estipulando formato PDF incorporando operativamente orientando y abordando e incrustando el medio referido al referido respectivo servicio o plataforma logísticamente abarcando lo relativo e innegable de lo referido enfocando abarcando refiriéndose respectando lo relativo u orientado abordando lo de la mensajería del correo referenciado enfocado abordando refiriéndose electrónico.
\end{itemize}

\subsection{Escalas de probabilidad e impacto}

El equipo técnico utiliza firmemente y de modo estandarizado un modelo cualitativo para asignar niveles basándose y estableciendo ponderaciones numéricas del uno al cinco y logra con esto consolidar una innegociable métrica uniforme logrando estructurar una equidad evaluativa y evitando fallos referidos a percepciones erróneas y diferencias al medir de manera consistente los parámetros evaluando tanto a las probabilidades originando un nivel respectivo abarcando al respectivo y certero cálculo del daño documentado o nivel respectivo medible del impacto.

\subsubsection{Escala de probabilidad}

La respectiva matriz o tabla referenciada o ilustrada expone íntegramente las variables estipuladas con las respectivas ponderaciones operativas correspondientes y detalladas desde la baja calificación numerada refiriéndose al uno estructurando o detallando gradualmente una numeración ascendente alcanzando e incluyendo progresivamente documentando hasta la cúspide límite referenciada numéricamente y topada o respectando el respectivo puntaje limitante de un límite numeral de cinco estableciendo de este modo logrando con exactitud y fiabilidad identificar estandarizadamente a niveles precisos estipulando e incrustando metódicamente la esperada respectiva o certera probabilidad de toda respectiva ocurrencia analizada e identificada.

\begin{widepage}
{\footnotesize
\setlength{\tabcolsep}{3pt}
\begin{longtable}{c P{3.5cm} P{18.5cm}}
\caption{Escala de probabilidad}\label{tab:escala-probabilidad} \\
\toprule
\textbf{Valor} & \textbf{Nivel} & \textbf{Descripción} \\
\midrule
\endfirsthead
\caption[]{Escala de probabilidad (continuación)} \\
\toprule
\textbf{Valor} & \textbf{Nivel} & \textbf{Descripción} \\
\midrule
\endhead
\bottomrule
\endfoot
\bottomrule
\endlastfoot
\textbf{1} & \textbf{Raro} & Establece contundentemente que la respectiva y documentada o evaluada probabilidad es ínfima estableciendo claramente que es verdaderamente muy poco referida o verdaderamente muy poco probable de modo que resulte sumamente extraño que efectivamente logre o llegase verdaderamente y contundentemente a que ocurra efectivamente un incidente así. \\
\textbf{2} & \textbf{Poco probable} & Manifiesta e indica que el suceso evaluado o el incidente logísticamente puede verdaderamente de forma eventual a que llegue e inclusive a que logre a que ocurra de modo aislado estableciendo que en la coyuntura respectiva no es un suceso respectivo y referenciado que se estipule ni que se espere operativamente o que verdaderamente y de manera efectiva se espere a que suceda en un rango inminente. \\
\textbf{3} & \textbf{Moderado} & Corrobora y testifica íntegramente referenciando que se mantiene vigente documentando de lleno o efectivamente que existe analíticamente y evidenciablemente operativamente comprobado en verdad documentando a fondo una certera validación indicando y documentando indiscutiblemente verdaderamente abarcando logrando estructurar u otorgar certeza certificando la firmeza aseverando referenciando de modo efectivo estipulando y verificando que a fin o respecto operativamente se dictamine estableciendo verdaderamente y afirmando con certeza o refrendando que yace latente operativamente evidenciándose referenciándose o abarcando de un modo certero y verdaderamente demostrable una muy innegable probabilidad y una eventual e indudable o certera o razonable y documentada posibilidad logísticamente referenciada abarcando y comprobando la respectiva coyuntura referenciada o concerniente refiriéndose operativamente y orientada a nivel referido orientando estipulando o a la inminente o respectiva e indudable ocurrencia en torno e inmersamente evidenciando y referida a presentarse en las tareas y fases del cronograma. \\
\textbf{4} & \textbf{Probable} & Certifica estableciendo operativamente de manera contundente y documentando y certificando innegablemente y refiriéndose con mucha certeza y asertivamente logrando confirmar que este es sin lugar a equívocos un escenario muy crítico y catalogado de tal nivel logrando fundamentar que resulta catalogado operativamente como sumamente e indiscutible y de un modo evidente evaluando a ser o estar propenso indicando referenciando indicando firmemente operativamente estipulando u orientando en un alto nivel u orientando e indicando verdaderamente referido a ser o que de plano estipulando a que llegue firmemente o refiriéndose o indicando a que ocurra logísticamente evaluando que el factor verdaderamente es e indica ser sumamente u orientando de una inminente probabilidad o calificando operativamente como verdaderamente un incidente sumamente probable logrando verdaderamente estipulando y estableciendo estipulando verdaderamente innegablemente un riesgo indicando un suceso comprobando referenciando y evidenciando si este suceso al no implementar acciones el cual no se controla o se mitiga proactivamente con inmediatez respectiva e indudable operativamente requerida. \\
\textbf{5} & \textbf{Casi seguro} & Asegura firmemente de un modo irrefutable comprobando irrefutable y tajantemente que ciertamente es abarcando indiscutiblemente referenciando o documentando en nivel sumo innegablemente evaluando de un modo certero sumamente y contundentemente o referido o evaluando que resulta altamente o verdaderamente e irrefutablemente de un modo tajantemente probable afirmando operativamente e irrevocablemente en torno referenciando y evidenciando o estipulando o evaluando comprobando irrevocablemente evaluando u orientando un suceso referido aludiendo a la certera estipulación y validando indiscutiblemente certificando y documentando que ocurra e indicando indiscutiblemente o innegablemente y referenciando y evaluando su llegada certificando su presencia estipulando o aseverando a que ocurra u ocurra a que este suceda o ocurra logísticamente a lo largo abarcando de un modo o evaluando de lleno e irrevocablemente a todo transcurrir operativamente a lo largo del respectivo e inminente tiempo u orden estipulado abarcando el transcurrir abarcado del proyecto. \\
\end{longtable}
}
\end{widepage}

\subsubsection{Escala de impacto}

La tabla mostrada a continuación describe la escala e ilustra el grado concerniente de un modo u otro referenciando del uno hacia la calificación límite e integralmente delimitada del cinco donde se describe explícitamente a las referidas alteraciones e indica operativamente a la grave respectiva y consecuente e inminente e ineludible y perjudicial y respectiva y documentada o referida o respectiva comprobada posible o grave respectiva afectación operativa que puede acaecer u ocasionalmente a un nivel inmerso y relativo y consecuente a la inminente y a que la vulnerabilidad referida abarcando e indicando y referenciada en cada respectivo e importante respectivo rango abarcado referenciado estipulando o abarcando e indicando el referido y cada respectivo rango y a que cada específico nivel respectivo al que cada nivel evidenciado o un referenciado referido estipulado inmersamente referido nivel puede verdaderamente logísticamente desencadenando desastres y abarcando o evidenciando la posibilidad respectiva logrando u originando llegar indudablemente operativamente llegando a referenciando o referenciada a lograr verdaderamente o abarcando o evidenciando estipulando e incrustando el daño abarcando referenciando u orientando el nivel de desestabilización que se referencie logrando a provocar y abarcando logrando o llegando a referenciando a referir y referenciando englobando al que llegue a estipular operativamente en la base estructurada de este referenciado u referido proyecto.

\begin{widepage}
{\footnotesize
\setlength{\tabcolsep}{3pt}
\begin{longtable}{c P{3.5cm} P{18.5cm}}
\caption{Escala de impacto}\label{tab:escala-impacto} \\
\toprule
\textbf{Valor} & \textbf{Nivel} & \textbf{Descripción} \\
\midrule
\endfirsthead
\caption[]{Escala de impacto (continuación)} \\
\toprule
\textbf{Valor} & \textbf{Nivel} & \textbf{Descripción} \\
\midrule
\endhead
\bottomrule
\endfoot
\bottomrule
\endlastfoot
\textbf{1} & \textbf{Insignificante} & Afecta de un modo o abarca e incide con un muy bajo nivel de un nivel evaluando u originando de forma o de una innegable o evidente y de verdaderamente comprobada u orientada de un modo y referida a que incida o en un grado y forma verdaderamente o a un grado comprobado referenciando y referida a mínima e indica además comprobando y aseverando verdaderamente que esto respectando logísticamente o aseverando u orientando indicando e indicando que el inconveniente o fallo u obstáculo y este se logra referenciando o documentando verdaderamente indicando que verdaderamente abarcando e indicando evaluando verdaderamente que este de plano o referenciando u orientando a la reparación se enmienda y referenciando y abarcando que el fallo documentado se corrige e inmersamente referenciando o documentando evaluando a un modo o se estipula evaluando de modo orientando referenciando y refiriéndose estipulando a lograr orientando evaluando estipulando orientando o corrigiéndose e implementando enmendando referenciando abarcando logísticamente a lograr a corregirse referenciando estipulando y operativamente de manera refiriéndose estipulando evaluando estipulando implementando a que se remedia de forma o verdaderamente abarcando a una referida reparación estipulando de modo fácilmente y orientando a su pronta resolución. \\
\textbf{2} & \textbf{Menor} & Genera innegablemente una requerida u orientada logísticamente o referenciando u abarcando referenciada a la referida e ineludible y correspondiente respectiva adaptación abarcando de modo o referenciando e indicando a una refiriéndose operativamente estipulando a un estipulando y referenciando a un correspondiente e ineludible y muy puntual correspondiente respectivo u orientado correspondiente y referenciado a originando un ligero u orientando respectando o a un estipulando u orientando abarcando estipulando correspondiente referenciando y estipulando y englobando a la aplicación referenciada de un ligero estipulando o a originar y referido u orientando al debido u orientado referenciando e indicando un respectivo ajuste orientando e indicando que evaluando referenciando que este orientando y estipulando aludiendo al ajuste menor o referenciando o limitando la gravedad evaluando e implementando y aplicable operativamente abarcando respectando e indicando a originar o referenciando a que se origine y limite logísticamente u orientando su alcance inmersamente englobando o refiriéndose o recayendo estipulando limitando abarcando estipulando recayendo referenciando e incrustando su efecto en tareas referenciadas o limitadas en componentes o en archivos documentando evaluando recayendo o en referenciados u orientando a los referenciados en los archivos abordando e indicando o en los referidos y estipulando en documentos. \\
\textbf{3} & \textbf{Significativo} & Puede comprobar e indica logísticamente u orientando a que evalúa a que pueda o llegar refiriéndose e indicando u orientando a llegar evaluando a lograr afectar operativamente referenciando indicando limitando o estipulando u orientando e indicando inmersamente afectando o logrando referenciando logísticamente a afectar estipulando e incrustando el daño y referido al respectivo estándar de o al aspecto respectivo y a referenciando la variable respectiva y al referirse operativamente indicando la referida variable y de este modo abarcando a la debida u orientando al aseguramiento refiriéndose indicando a la variable respectiva a lograr afectando a orientando referenciando a la variable o englobando y la calidad referenciando operativamente orientando y al margen referido e indicando limitando o orientando indicando a todo estipulando e indicando limitando al transcurso respectivo o a la franja orientando indicando al tiempo e inmersamente referenciando al factor o a estipular y al englobar al rubro de o abarcando indicando limitando o englobando refiriéndose y evaluando la correspondiente logística referenciando a la debida e indispensable u orientando a la estipulando e inmersamente a la referida y estipulando a referenciada a la a la o respectiva referenciando englobando e indicando a toda a la respectiva o a englobar a toda o coordinación. \\
\textbf{4} & \textbf{Mayor} & Puede llegar inmersamente refiriéndose operativamente a que logre a que documentando refiriendo o evaluando indicando refiriendo o e indicando y llegar inmersamente indicando limitando a originar e indicando o afectar e incrustando un respectivo o orientando u orientando e incrustando un daño e indicando e incrustando e indicando e impactar e incrustando a orientando refiriendo indicando estipulando referenciando afectando abarcando o a retrasar limitando orientando o a afectar entregables evaluando estipulando referenciando a que evalúa e incide o retrasar inmersamente a los referenciados o limitando u orientando a afectar componentes importantes evaluando o de un modo inmerso o a componentes sumamente e innegablemente o en referenciando a que recaiga limitando o en elementos referenciados limitando e importantes o en elementos refiriéndose limitando e incrustando refiriéndose u orientando al entorno del proyecto referenciando e indicando estipulando al proyecto evaluando e incrustando al estipulando del respectivo documentando a la variable del estipulando respectivo o a la propia índole o del respectivo y estipulando inmerso en referido proyecto. \\
\textbf{5} & \textbf{Severo} & Puede a su vez llegar referenciando a afectar referenciando de modo indicando refiriendo estipulando u orientando e incrustando y refiriéndose limitando estipulando o a afectar de manera u orientando innegable e indicando operativamente limitando u orientando de manera contundente y afectando de modo referenciando a todo limitando a toda respectiva orientando indicando limitando e incrustando la inmersa a toda logística limitando e indicando o a referenciando a toda o a la limitando e indicando e incrustando la entrega indicando e incrustando final referenciando estipulando referenciando u orientando inmerso indicando o limitando e indicando u orientando e incidiendo estipulando e incrustando al nivel e incrustando referenciando a todo limitando a todo estipulando o al respectivo o estipulando al o al cumplimiento estipulando o a referenciando e incrustando a lograr indicando e incidiendo limitando y afectando o imposibilitar referenciando estipulando e incrustando o al cumplimiento orientando estipulando y al nivel estipulando referenciando e indicando o refiriendo e indicando al documentando o al orientando estipulando y del estipulando respectivo y del alcance referenciando estipulando e indicando y referenciando o del alcance. \\
\end{longtable}
}
\end{widepage}

\subsection{Fórmula del nivel de riesgo}

La determinación métrica empleada se calcula directamente multiplicando el valor y puntaje de la probabilidad por el impacto.

\textbf{Nivel de riesgo = probabilidad $\times$ impacto}

\subsection{Clasificación del nivel de riesgo}

El resultado final agrupa la valoración y ubica los parámetros de los diferentes riegos en un rango progresivo de prioridades en una escala documentada que abarca desde la categoría muy baja hasta la categoría crítica como se detalla en la siguiente tabla.

\begin{table}[H]
\centering
\caption{Clasificación del nivel de riesgo}
\label{tab:clasificacion-riesgo}
\begin{tabularx}{0.6\textwidth}{X X}
\toprule
\textbf{Resultado (P $\times$ I)} & \textbf{Clasificación} \\
\midrule
\cellcolor{C_MuyBajo}\textbf{1 a 2} & \cellcolor{C_MuyBajo}\textbf{Muy bajo} \\
\cellcolor{C_Bajo}\textbf{3 a 4} & \cellcolor{C_Bajo}\textbf{Bajo} \\
\cellcolor{C_Medio}\textbf{4 a 9} & \cellcolor{C_Medio}\textbf{Medio} \\
\cellcolor{C_Alto}\textbf{10 a 14} & \cellcolor{C_Alto}\textbf{Alto} \\
\cellcolor{C_MuyAlto}\textbf{15 a 19} & \cellcolor{C_MuyAlto}\textbf{Muy alto} \\
\cellcolor{C_Critico}\textbf{\textcolor{white}{20 a 25}} & \cellcolor{C_Critico}\textbf{\textcolor{white}{Crítico}} \\
\bottomrule
\end{tabularx}
\end{table}

El valor asignado a cada riesgo posee un significado que dicta las acciones a implementar.

\begin{itemize}
    \item \textbf{Muy bajo:} Requiere implementar monitoreo básico rutinario y preventivo.
    \item \textbf{Bajo:} Resulta aceptado y se tolera implementando un nivel moderado de monitoreo.
    \item \textbf{Medio:} Requiere una revisión profunda y elaborar una planeación para mitigar el suceso.
    \item \textbf{Alto:} Requiere esfuerzos coordinados garantizando una mitigación con atención prioritaria.
    \item \textbf{Muy alto:} Requiere tomar una acción inmediata focalizada a reducir la amenaza sobre la estructura principal para estabilizar la logística del grupo.
\end{itemize}

\subsection{Estrategias de respuesta ante riesgos}

Para cada riesgo identificado el equipo debe evaluar el problema y tiene la responsabilidad de seleccionar alguna de estas cuatro alternativas con directrices de prevención correctivas documentadas puntualmente a continuación.

\begin{widepage}
{\footnotesize
\setlength{\tabcolsep}{3pt}
\begin{longtable}{P{2.5cm} P{10.3cm} P{10.3cm}}
\caption{Estrategias de respuesta}\label{tab:estrategias-respuesta} \\
\toprule
\textbf{Estrategia} & \textbf{Descripción} & \textbf{Ejemplo de aplicación en el proyecto} \\
\midrule
\endfirsthead
\caption[]{Estrategias de respuesta (continuación)} \\
\toprule
\textbf{Estrategia} & \textbf{Descripción} & \textbf{Ejemplo de aplicación en el proyecto} \\
\midrule
\endhead
\bottomrule
\endfoot
\bottomrule
\endlastfoot
\textbf{Evitar} & Cambiar la forma de trabajo original para eliminar la posibilidad de que el riesgo se manifieste y afecte gravemente el alcance o cronograma o la calidad. & Revisar toda solicitud comparándola contra el alcance aprobado antes de aceptar modificaciones al prototipo. \\
\textbf{Mitigar} & Implementar acciones preventivas enfocadas en reducir eficientemente la probabilidad de ocurrencia como el nivel de impacto destructivo de la amenaza. & Ejecutar una validación exhaustiva de todo el sistema de navegación visual justo antes de oficializar cada nueva entrega del prototipo. \\
\textbf{Transferir} & Asignar parte de la responsabilidad o del impacto financiero del riesgo hacia un tercero debido a una dependencia de herramientas o servicios ajenos a la organización. & Respaldar rigurosamente el cronograma oficial y el código base dentro del repositorio para minimizar la dependencia técnica de un solo dispositivo local. \\
\textbf{Aceptar} & Reconocer abiertamente la existencia del riesgo y optar por mantenerlo en observación continua sin aplicar una acción inmediata invasiva. & Asumir la escasa retroalimentación por parte de los interesados externos debido al limitado margen de tiempo estipulado para la culminación del proyecto. \\
\end{longtable}
}
\end{widepage}

\subsection{Elementos del registro de riesgos}

La documentación de cada riesgo alojado en el registro de riesgos debe estructurarse contemplando los trece elementos constitutivos descritos a continuación.

\begin{enumerate}
    \item El código de identificación único del riesgo.
    \item La categoría de impacto principal que incluye áreas como cronograma o el alcance o documentación o la calidad o comunicación.
    \item La descripción detallada de la naturaleza del riesgo.
    \item La causa raíz subyacente que desencadena la amenaza potencial.
    \item El impacto proyectado sobre los objetivos globales.
    \item La asignación de la probabilidad basada en la escala numérica de uno a cinco.
    \item La cuantificación del impacto destructivo basada en la escala numérica de uno a cinco.
    \item El nivel matemático resultante de la multiplicación directa entre los factores de probabilidad y el impacto correspondiente.
    \item La clasificación de severidad establecida que categoriza el problema de muy bajo a crítico.
    \item La designación clara del integrante responsable de mantener el seguimiento constante.
    \item La estrategia de respuesta adoptada mediante acciones para evitar o mitigar o transferir o aceptar el riesgo.
    \item La acción de respuesta concreta y tangible que debe ejecutarse inmediatamente.
    \item El estado actual de vigencia y tratamiento del riesgo.
\end{enumerate}

\subsubsection{Estados del riesgo}

Todos los riesgos documentados e identificados formalmente operan bajo una clasificación dinámica de cuatro posibles estados cuyo propósito facilita la medición continua de su seguimiento activo o la consolidación de su solución.

\begin{widepage}
{\footnotesize
\setlength{\tabcolsep}{3pt}
\begin{longtable}{P{3.5cm} P{20cm}}
\caption{Estados del riesgo}\label{tab:estados-riesgo} \\
\toprule
\textbf{Estado} & \textbf{Descripción} \\
\midrule
\endfirsthead
\caption[]{Estados del riesgo (continuación)} \\
\toprule
\textbf{Estado} & \textbf{Descripción} \\
\midrule
\endhead
\bottomrule
\endfoot
\bottomrule
\endlastfoot
\textbf{Abierto} & Indica que el riesgo fue documentado recientemente y requiere mantenerse bajo observación constante. \\
\textbf{En seguimiento} & Muestra que el equipo asignado se encuentra aplicando de manera activa las estrategias y acciones de control diseñadas. \\
\textbf{Materializado} & Advierte que la amenaza superó las barreras de contención y se convirtió en un problema concreto que amerita una resolución inmediata. \\
\textbf{Cerrado} & Confirma la finalización exitosa de las intervenciones y constata que la situación ya no representa una amenaza para el proyecto. \\
\end{longtable}
}
\end{widepage}

\subsection{Monitoreo, control y actualización}

El equipo de trabajo tiene la obligación de actualizar continuamente el registro de riesgos a lo largo del desarrollo del proyecto cada vez que se detecte una variación importante en los valores de probabilidad e impacto o bien cuando se modifique el responsable o la acción de respuesta o el estado de algún riesgo. El documento debe nutrirse de forma inmediata ante el surgimiento imprevisto de un riesgo completamente nuevo o cuando una amenaza latente traspasa el nivel de previsión y se transforma en un problema materializado.

\subsubsection{Mecanismo de seguimiento}

La fase operativa de identificación o mitigación y eventual eliminación de las amenazas recae en la aplicación disciplinada de acciones de seguimiento que permiten asegurar un control riguroso de todo el ecosistema de riesgos.

\begin{itemize}
    \item Efectuar una revisión detallada de los riesgos críticos en el marco de cada reunión formal de coordinación y seguimiento.
    \item Informar y escalar oportunamente mediante los canales de comunicación oficiales y plataformas corporativas cualquier riesgo urgente con capacidad para comprometer negativamente la entrega final del producto tecnológico esperado.
    \item Mantener una documentación estructurada sobre cualquier contratiempo potencial que posea el poder de afectar negativamente factores esenciales de éxito como el alcance aprobado o el cronograma general o los estándares de calidad comprometidos con los interesados de la organización.
    \item Aplicar de inmediato las enmiendas documentales de rigor en el momento exacto en que un evento monitoreado registre un aumento o disminución palpable en los niveles de criticidad previamente determinados.
    \item Vincular integralmente las incidencias y factores de amenaza al marco procedimental central que rige el proceso oficial de control de cambios siempre que tales adversidades obliguen a ejecutar modificaciones directas sobre el alcance o la planificación en el cronograma original.
    \item Ejecutar sesiones exhaustivas destinadas a comprobar el estado de vulnerabilidad y posibles puntos de falla que se asocien con toda la estructura técnica en el desarrollo de la interfaz justo antes de consolidar oficialmente cada entrega parcial agendada de antemano y la entrega definitiva del sistema esperado.
    \item Cuidar la correcta organización y registro minucioso y salvaguarda estricta de todo tipo de evidencia de carácter físico o digital dentro del repositorio correspondiente agrupando lógicamente capturas de pantalla o código y documentos técnicos para lograr mitigar y reducir efectivamente el latente riesgo derivado por una pérdida desastrosa de la información base del trabajo y del avance histórico de la plataforma digital.
\end{itemize}

\subsubsection{Frecuencia de revisión}

Considerando la duración total estipulada para el proyecto el equipo establece un cronograma con frecuencias mínimas indispensables que regirán estrictamente la revisión periódica del registro de riesgos según se indica a continuación.

\begin{itemize}
    \item \textbf{Riesgos clasificados como muy alto o crítico:} Exigen una evaluación exhaustiva y prioritaria en cada reunión de equipo y mantienen una periodicidad mínima de seguimiento semanal.
    \item \textbf{Riesgos clasificados como medio o alto:} Requieren un chequeo programado de carácter quincenal o bien una verificación inmediata ante la aparición de cualquier cambio relevante en el entorno.
    \item \textbf{Riesgos clasificados como bajo o muy bajo:} Cumplen con una revisión de carácter general al concretar el cierre de cada fase importante del proyecto para confirmar su estabilidad.
\end{itemize}

\newpage
% --- Entregable 14 ---
\section{Registro de riesgos}

La presente sección detalla el registro de riesgos elaborado para el proyecto denominado Plataforma Universitaria Integrada, el cual es desarrollado para la Universidad Tecnológica La Mejor. Este registro constituye un componente fundamental dentro de la planificación y control del proyecto, para lo cual se tiene como finalidad principal la identificación de las posibles eventualidades que podrían afectar el alcance, el cronograma, el costo, la calidad, la comunicación, el desarrollo técnico de la interfaz de usuario y el cierre exitoso del proyecto.

La gestión oportuna de estos riesgos es de suma importancia debido a que permite anticipar diversas situaciones adversas que podrían generar atrasos, retrabajo, pérdida de información, defectos en el prototipo final o inconsistencias documentales. En consecuencia, en esta iniciativa en particular, los riesgos deben ser analizados de manera sumamente minuciosa ya que el equipo de trabajo cuenta con un cronograma planificado de tres meses y debe cumplir con múltiples entregables documentales, además de gestionar roles combinados y un prototipo de interfaz navegable cuya finalización está programada según la fecha aprobada en Microsoft Project.

El proyecto consiste propiamente en el desarrollo de una Plataforma Universitaria Integrada implementada a nivel de prototipo web navegable, para lo cual se deben representar de forma interactiva diversos módulos como el inicio de sesión, el panel de control principal, el portal estudiantil, el portal docente, el portal administrativo, el portal financiero y el panel ejecutivo. Asimismo, además de este prototipo técnico, la propuesta incorpora múltiples documentos de gestión orientados a complementar y consolidar la planificación estructurada del proyecto.

El alcance detallado del proyecto establece con claridad que no se desarrollará la lógica funcional del servidor, bases de datos reales, sistemas de autenticación institucional integrados, pasarelas de pago reales, migración de datos históricos ni un despliegue en producción, de manera que varios riesgos identificados se orientan a evitar que tanto el equipo de desarrollo como los interesados esperen funcionalidades que se encuentran fuera del alcance aprobado. Por lo tanto, otros riesgos importantes se asocian directamente con la coherencia de la documentación, la disponibilidad de los miembros del equipo, la calidad técnica del prototipo navegable, la efectividad en la comunicación interna y el uso de las distintas herramientas especializadas de desarrollo.

Este registro de riesgos se ha estructurado a partir del contexto y de los objetivos generales de la iniciativa, de manera que se asegura que las amenazas detectadas no se planteen de forma aislada, sino estrechamente vinculadas con las condiciones reales del entorno y con los elementos organizacionales definidos formalmente durante la etapa de planificación.

\subsection{Propósito del registro de riesgos}

El propósito central de este registro de riesgos es documentar de forma estructurada los eventos inciertos que posean el potencial de afectar de manera negativa el cumplimiento de las metas del proyecto, para lo cual esta sección facilita la identificación precisa de las causas de cada riesgo, su impacto estimado, su nivel de prioridad, el responsable asignado para su monitoreo y la estrategia de respuesta adecuada que se deba aplicar.

Este instrumento se orienta al cumplimiento de objetivos específicos detallados a continuación.

\begin{itemize}[leftmargin=*, label={\textbullet}]
    \item Identificar los riesgos principales del proyecto.
    \item Analizar la causa y el posible impacto de cada riesgo.
    \item Evaluar la probabilidad y el impacto mediante una escala definida.
    \item Clasificar los riesgos según su nivel de criticidad.
    \item Asignar responsables para dar seguimiento a cada riesgo.
    \item Definir acciones de respuesta para prevenir o reducir sus efectos.
    \item Mantener trazabilidad entre riesgos, alcance, cronograma, calidad, comunicaciones y cambios.
    \item Apoyar la toma de decisiones del equipo durante la ejecución y cierre del proyecto.
\end{itemize}

Asimismo, el registro de riesgos de la Plataforma Universitaria Integrada resulta fundamental para priorizar aquellos eventos que demandan una atención inmediata por parte de la gestión, debido a que no todas las amenazas presentan la misma relevancia para el proyecto. Por lo tanto, un error menor de formato en la documentación no representará el mismo impacto que un atraso en la codificación de la interfaz navegable o una modificación extemporánea en el alcance, por lo cual se implementará una evaluación basada en probabilidad e impacto con el fin de determinar las prioridades de atención.

\subsection{Metodología de evaluación de riesgos}

Para evaluar de manera adecuada las amenazas identificadas se utilizará un enfoque cualitativo centrado en la probabilidad y el impacto, de modo que cada riesgo sea analizado a partir de la posibilidad de su ocurrencia y del efecto directo que causaría sobre el desarrollo si llegara a materializarse.

Esta valoración se ejecutará empleando una escala cuantitativa de uno a cinco tanto para la probabilidad como para el impacto, para lo cual posteriormente se calculará el nivel de riesgo definitivo mediante el producto de ambos factores.

La determinación del nivel de riesgo se realiza mediante la siguiente ecuación.

\begin{center}
Nivel de riesgo = Probabilidad x Impacto
\end{center}

A partir del valor obtenido se procederá a clasificar cada riesgo en las categorías de bajo, medio, alto o crítico, lo que implica que esta segmentación facilitará al equipo la identificación inmediata de las amenazas que requieren un monitoreo prioritario.

El análisis detallado de cada riesgo contemplará los elementos que se enumeran a continuación.

\begin{itemize}[leftmargin=*, label={\textbullet}]
    \item Código del riesgo.
    \item Categoría.
    \item Descripción del riesgo.
    \item Causa.
    \item Impacto en el proyecto.
    \item Probabilidad.
    \item Impacto.
    \item Nivel resultante.
    \item Clasificación.
    \item Responsable.
    \item Estrategia de respuesta.
    \item Acción de respuesta.
    \item Estado.
\end{itemize}

\subsection{Escala de probabilidad e impacto}

La evaluación de la probabilidad y el impacto en este registro se encuentra alineada con la matriz de probabilidad e impacto definida para la organización, en la cual se detalla la escala de referencia junto con la ubicación de las trece amenazas identificadas. Por lo tanto, con el propósito de asegurar la consistencia entre los documentos del proyecto, este registro adopta de manera estricta los mismos valores de probabilidad, impacto, nivel de riesgo y clasificación acordados en dicho instrumento de análisis.

\subsection{Estrategias de respuesta ante riesgos}

Para cada riesgo identificado se establecerá una estrategia de respuesta específica, de manera que las acciones correspondan a una de las categorías descritas a continuación.

\subsubsection{Evitar}

Esta estrategia consiste en modificar los procesos o planes de trabajo con el objetivo de eliminar por completo la probabilidad de ocurrencia de la amenaza, de modo que se aplica prioritariamente en aquellos eventos que podrían comprometer gravemente el alcance, el cronograma o la calidad de los entregables.

\subsubsection{Mitigar}

Esta alternativa busca reducir la probabilidad de ocurrencia o la severidad del impacto mediante la ejecución de acciones preventivas, lo cual representa la estrategia de mayor uso en el proyecto debido a que la mayoría de los riesgos pueden ser gestionados mediante actividades constantes de revisión, comunicación y seguimiento.

\subsubsection{Transferir}

Esta modalidad implica trasladar la responsabilidad o el impacto financiero y operativo del riesgo hacia un tercero ajeno al equipo de desarrollo, la cual se empleará de manera limitada en esta iniciativa debido a que solo se aplicará ante dependencias tecnológicas externas o fallos en los servicios de almacenamiento en la nube.

\subsubsection{Aceptar}

Esta estrategia consiste en asumir formalmente la existencia del riesgo y realizar un monitoreo pasivo sin implementar acciones preventivas inmediatas, de manera que se utiliza en situaciones donde el impacto es mínimo o cuando el costo de mitigación supera los beneficios esperados de la intervención.

\subsection{Fuentes utilizadas para identificar riesgos}

La identificación de las amenazas se llevó a cabo mediante una revisión exhaustiva de las bases de planificación del proyecto, para lo cual las fuentes documentales utilizadas corresponden a las que se mencionan a continuación.

\begin{itemize}[leftmargin=*, label={\textbullet}]
    \item Caso de negocio.
    \item Acta de constitución del proyecto.
    \item Registro de interesados.
    \item Definición del alcance.
    \item Cronograma del proyecto detallado en Microsoft Project.
    \item Organigrama estructurado del proyecto.
    \item Plan de gestión de calidad.
    \item Plan de gestión de comunicaciones.
    \item Restricciones asociadas a los plazos y al alcance de los entregables.
    \item Distribución interna de roles entre los miembros del equipo.
    \item Requerimientos técnicos para el desarrollo del prototipo de interfaz navegable.
    \item Herramientas tecnológicas de trabajo cooperativo incluyendo GitHub, Microsoft Project, procesadores de texto, plataformas de comunicación y correo electrónico.
\end{itemize}

La consulta detallada de estos recursos propició la detección de amenazas asociadas con los plazos de entrega, el alcance acordado, la calidad de los productos, la coherencia de la documentación, los canales de comunicación, el funcionamiento de las herramientas y la coordinación técnica del desarrollo.

\begin{widepage}
\subsection{Registro de riesgos del proyecto}
{\tiny
\setlength{\tabcolsep}{2pt}
\begin{longtable}{L{0.9cm} L{1.9cm} L{2.1cm} L{2.4cm} L{2.3cm} L{1.8cm} L{1.2cm} L{0.9cm} L{1.7cm} L{2.2cm} L{1.4cm} L{2.6cm} L{1.0cm}}
\caption{Registro de riesgos del proyecto}\label{tab-registro-riesgos}\\[-0.5em]
\toprule
\textbf{Código} & \textbf{Categoría} & \textbf{Riesgo} & \textbf{Causa} & \textbf{Impacto en el proyecto} & \textbf{Probabilidad} & \textbf{Impacto} & \textbf{Nivel} & \textbf{Clasificación} & \textbf{Responsable} & \textbf{Estrategia} & \textbf{Acción de respuesta} & \textbf{Estado} \\
\midrule
\endfirsthead
\toprule
\textbf{Código} & \textbf{Categoría} & \textbf{Riesgo} & \textbf{Causa} & \textbf{Impacto en el proyecto} & \textbf{Probabilidad} & \textbf{Impacto} & \textbf{Nivel} & \textbf{Clasificación} & \textbf{Responsable} & \textbf{Estrategia} & \textbf{Acción de respuesta} & \textbf{Estado} \\
\midrule
\endhead
R-01 & Cronograma & Restricción fuerte de tiempo & El proyecto debe entregarse según la fecha final establecida en el cronograma aprobado y contiene varios entregables documentales y técnicos & Atraso en documentos, reducción de calidad o entrega incompleta & 4 & 5 & 20 & Crítico & Tamara Robles Camacho & Mitigar & Dar seguimiento frecuente al cronograma, priorizar tareas críticas y redistribuir trabajo si una actividad se atrasa & Abierto \\
R-02 & Alcance & Alcance creciente & Se pueden proponer funcionalidades fuera del alcance, como lógica de servidor, base de datos real o pagos reales & Aumento del trabajo, atraso en la interfaz y contradicción con la definición del alcance & 4 & 5 & 20 & Crítico & Tamara Robles Camacho e Isaac Jiménez Blanco & Evitar & Revisar toda nueva solicitud contra el alcance aprobado y aplicar control formal de cambios & Abierto \\
R-03 & Documentación & Pérdida de coherencia en el informe final & La integración de múltiples secciones redactadas por distintos integrantes puede generar contradicciones en el estilo y el formato & Reducción en la calidad del proyecto por falta de uniformidad profesional & 4 & 4 & 16 & Muy alto & Isaac Jiménez Blanco y Jorge Abarca Quiros & Mitigar & Realizar una revisión cruzada exhaustiva del documento final para unificar el estilo y contenido & Abierto \\
R-04 & Desarrollo & Subestimación del esfuerzo de la interfaz & El prototipo incluye varios módulos, pantallas y navegación entre portales & El prototipo puede quedar incompleto o con navegación limitada & 3 & 4 & 12 & Alto & Mauricio González Prendas y Carlos Sainz Arce & Mitigar & Priorizar primero los módulos obligatorios y dejar mejoras visuales no esenciales para el final & Abierto \\
R-05 & Recursos humanos & Disponibilidad limitada del equipo & El equipo es pequeño y cada integrante tiene varios entregables asignados & Sobrecarga de trabajo, atrasos y poca revisión entre compañeros & 4 & 4 & 16 & Muy alto & Tamara Robles Camacho & Mitigar & Coordinar revisiones cortas, definir prioridades y reasignar tareas cuando sea necesario & Abierto \\
R-06 & Calidad & Errores críticos de navegación en el prototipo & Falta de pruebas manuales suficientes o cambios rápidos en pantallas & El usuario no podría entrar a módulos principales o continuar el flujo de navegación & 3 & 5 & 15 & Muy alto & Mauricio González Prendas y Carlos Sainz Arce & Mitigar & Ejecutar lista de verificación de navegación en inicio de sesión, panel de control y portales principales antes de la entrega & Abierto \\
R-07 & Usabilidad & Baja usabilidad del prototipo & El prototipo puede tener pantallas completas pero poco claras para el usuario & Dificultad para entender el flujo de uso por parte de estudiantes, docentes o administrativos & 3 & 4 & 12 & Alto & Mauricio González Prendas y Carlos Sainz Arce & Mitigar & Revisar claridad de menús, botones, textos y distribución visual antes del cierre & Abierto \\
R-09 & Herramientas & Problemas con el repositorio o pérdida de versiones & Uso incorrecto del repositorio, archivos mal subidos o falta de entregas de código claras & Pérdida de trazabilidad del prototipo o documentos finales & 4 & 3 & 12 & Alto & Mauricio González Prendas y Carlos Sainz Arce & Mitigar & Usar nombres claros, descripciones precisas de cambios y respaldos locales de documentos y código & Abierto \\
R-10 & Cambios & Cambios tardíos cerca de la entrega & Interesados o integrantes pueden solicitar ajustes cuando el proyecto ya está avanzado & Afectación del cronograma, calidad y revisión final & 4 & 3 & 12 & Alto & Tamara Robles Camacho, Isaac Jiménez Blanco y Jorge Abarca Quiros & Mitigar & Definir fecha límite para cambios importantes y evaluar toda solicitud mediante control de cambios & Abierto \\
R-11 & Calidad documental & Falta de criterios claros de aceptación & Si los criterios no se revisan, el equipo puede considerar listo un módulo incompleto & Entregables aceptados internamente sin cumplir navegación, contenido mínimo o coherencia & 3 & 4 & 12 & Alto & Isaac Jiménez Blanco y Jorge Abarca Quiros & Mitigar & Usar los criterios del Plan de Calidad y revisar cada módulo contra el alcance aprobado & Abierto \\
R-12 & Comunicación & Mala comunicación entre integrantes & Uso excesivo de mensajes informales sin registrar acuerdos importantes & Duplicidad de tareas, pérdida de decisiones o documentos incongruentes & 3 & 4 & 12 & Alto & Tamara Robles Camacho & Mitigar & Usar WhatsApp para avisos rápidos, pero documentar decisiones importantes en GitHub, correo o minuta & Abierto \\
R-13 & Presupuesto & Exceso del presupuesto planificado & Los esfuerzos de desarrollo y corrección de errores en la interfaz podrían requerir más horas de trabajo & El proyecto superaría el límite de ₡17,625,200 afectando la viabilidad económica & 2 & 4 & 8 & Medio & Isaac Jiménez Blanco, Tamara Robles Camacho y Jorge Abarca Quiros & Mitigar & Monitorear el esfuerzo semanal y reasignar tareas para no superar el límite planificado & Abierto \\
R-14 & Interesados & Baja validación de interesados externos & Por tiempo limitado puede no consultarse a todos los interesados externos & El prototipo podría representar de forma incompleta algunos procesos docentes, financieros o administrativos & 3 & 4 & 12 & Alto & Tamara Robles Camacho & Aceptar y mitigar & Priorizar validaciones internas y consultar a los interesados externos solo cuando sea estrictamente necesario & Abierto \\
R-15 & Alcance técnico & Expectativa de funcionalidades no incluidas & Puede asumirse que el sistema tendrá lógica de servidor, datos reales, autenticación o pagos reales & Confusión sobre lo que realmente será entregado & 3 & 3 & 9 & Medio & Isaac Jiménez Blanco y Jorge Abarca Quiros & Evitar & Repetir en documentos y revisiones que el entregable técnico es un prototipo de interfaz navegable con datos simulados & Abierto \\
R-16 & Evidencia & Falta de evidencia final organizada & Los archivos pueden quedar dispersos entre WhatsApp, equipo local y repositorio & Dificultad para demostrar el cumplimiento de entregables y cronograma & 3 & 3 & 9 & Medio & Isaac Jiménez Blanco, Jorge Abarca Quiros, Mauricio González Prendas y Carlos Sainz Arce & Mitigar & Organizar documentos, PDFs, capturas, .mpp y código en GitHub con nombres claros & Abierto \\
\bottomrule
\end{longtable}
}
\end{widepage}

\subsection{Resumen de riesgos por nivel}

A partir de la evaluación de criticidad realizada, el volumen de riesgos identificados se distribuye en los términos que se exponen a continuación.

\begin{table}[H]
\centering
\caption{Resumen de riesgos por nivel}
\label{tab-resumen-riesgos}
{\small
\begin{tabularx}{\textwidth}{p{3.5cm} p{3.2cm} Y}
\toprule
\textbf{Clasificación} & \textbf{Cantidad de riesgos} & \textbf{Códigos} \\
\midrule
Crítico & 2 & R-01, R-02 \\
Muy alto & 3 & R-03, R-05, R-06 \\
Alto & 7 & R-04, R-07, R-09, R-10, R-11, R-12, R-14 \\
Medio & 3 & R-13, R-15, R-16 \\
Bajo & 0 & Ninguno \\
Muy bajo & 0 & Ninguno \\
\bottomrule
\end{tabularx}
}
\end{table}

Los riesgos clasificados en la categoría de críticos requieren una atención prioritaria y constante debido a que poseen el potencial de afectar de manera directa la entrega final de los productos. En consecuencia, estas amenazas corresponden principalmente a la restricción severa de los plazos de ejecución y al fenómeno del alcance creciente, puesto que ambos factores podrían alterar el cronograma maestro, incidir negativamente en la calidad del desarrollo y comprometer el alcance aprobado.

Asimismo, los riesgos identificados como muy altos deben ser objeto de un monitoreo estructurado ya que se asocian de manera directa con la consistencia global de la documentación, la disponibilidad horaria del equipo de trabajo y la eventual ocurrencia de fallos de navegación en el prototipo, de modo que, a pesar de no ubicarse en el extremo superior de la escala, podrían demandar retrabajo y retrasar entregables clave si no son prevenidos oportunamente.

Por otra parte, las amenazas de nivel alto se vinculan en gran medida con el diseño de la interfaz de usuario, la usabilidad de las pantallas principales, el funcionamiento de las herramientas de versionamiento, la comunicación del grupo, el procesamiento de cambios, los criterios de aceptación formales y la validación con interesados, además de los riesgos catalogados como medios que, si bien no revisten una urgencia inmediata, deben mantenerse bajo observación periódica con el fin de evitar incrementos indeseados en su nivel de severidad.

\subsection{Seguimiento y actualización de riesgos (registro)}

El presente registro de riesgos será sometido a actualizaciones periódicas a lo largo del ciclo de vida del proyecto en la medida en que varíen las condiciones de probabilidad, impacto, responsabilidades, acciones de respuesta o estado general de las amenazas identificadas, asimismo, se prevé la incorporación inmediata de cualquier nuevo evento detectado o la reclasificación de aquellos que se materialicen como problemas de ejecución.

La labor de seguimiento y control de estas amenazas se estructurará a través de un conjunto de actividades constantes.

\begin{itemize}[leftmargin=*, label={\textbullet}]
    \item Revisar los riesgos críticos durante las reuniones de seguimiento del equipo.
    \item Comunicar por mensajería instantánea cualquier riesgo urgente que pueda afectar la entrega.
    \item Documentar formalmente los riesgos que afecten alcance, tiempo, costo, calidad o cambios.
    \item Actualizar el registro cuando un riesgo aumente o disminuya su nivel.
    \item Relacionar los riesgos con el proceso de control de cambios cuando se requiera modificar alcance, cronograma o entregables.
    \item Revisar los riesgos asociados a la interfaz antes de la entrega final.
    \item Verificar que las acciones de respuesta se hayan aplicado cuando corresponda.
    \item Mantener evidencia organizada en el repositorio para reducir riesgos de pérdida de información.
\end{itemize}

\subsubsection{Estados utilizados para los riesgos}

\begin{table}[H]
\centering
\caption{Estados utilizados para los riesgos}
\label{tab-estados-riesgos-registro}
{\small
\begin{tabularx}{\textwidth}{p{4cm} Y}
\toprule
\textbf{Estado} & \textbf{Descripción} \\
\midrule
Abierto & El riesgo fue identificado y debe mantenerse en seguimiento \\
En seguimiento & El riesgo tiene acciones activas de control \\
Materializado & El riesgo ocurrió y debe tratarse como problema \\
Cerrado & El riesgo ya no representa amenaza significativa para el proyecto \\
\bottomrule
\end{tabularx}
}
\end{table}

\subsubsection{Responsables de seguimiento}

\begin{table}[H]
\centering
\caption{Responsables de seguimiento}
\label{tab-responsables-seguimiento}
{\small
\begin{tabularx}{\textwidth}{p{4.5cm} Y}
\toprule
\textbf{Responsable} & \textbf{Riesgos principales asociados} \\
\midrule
Tamara Robles Camacho & Cronograma, disponibilidad del equipo, comunicación, cambios y relación con los interesados \\
Isaac Jiménez Blanco & Documentación, alcance, calidad documental, presupuesto, riesgos, cambios y evidencia \\
Jorge Abarca Quiros & Elaboración de entregables documentales, apoyo en gestión de riesgos, costos y comunicaciones del proyecto \\
Mauricio González Prendas & Desarrollo de interfaz, control de calidad, navegación, usabilidad, repositorios y correcciones técnicas \\
Carlos Sainz Arce & Desarrollador de interfaz 2, control de calidad, pruebas de calidad y navegación del portal financiero y ejecutivo \\
\bottomrule
\end{tabularx}
}
\end{table}

\newpage
% --- Entregable 15 ---
\section{Matriz de probabilidad e impacto}

La matriz de probabilidad e impacto es la representación gráfica que nos permite visualizar de forma inmediata el nivel de criticidad de cada riesgo identificado en el proyecto Plataforma Universitaria Integrada. Esta matriz combina la escala de probabilidad en el eje vertical con la escala de impacto en el eje horizontal y sigue la metodología establecida en el plan de gestión de riesgos.

Cada celda de la matriz representa el resultado de multiplicar probabilidad por impacto y se colorea según el rango de clasificación correspondiente. Los rangos son Bajo en verde o Medio en amarillo o Alto en naranja o Crítico en rojo y los 16 riesgos documentados en el registro de riesgos se ubican dentro de la matriz según su evaluación específica.

\subsection{Escalas de referencia}

La matriz se elabora a partir de las escalas definidas en el plan de gestión de riesgos donde la probabilidad e impacto son evaluados de 1 a 5 y el nivel de riesgo resulta de la fórmula matemática donde el nivel es igual a la probabilidad multiplicada por el impacto.

\begin{table}[H]
\centering
\begin{tabularx}{0.8\textwidth}{c c c Y}
\toprule
\textbf{Resultado (P $\times$ I)} & \textbf{Clasificación} & \textbf{Color} & \textbf{Acción requerida} \\
\midrule
1 a 5 & \cellcolor{C_Bajo} Bajo & Verde & Monitoreo básico \\
6 a 10 & \cellcolor{C_Medio} Medio & Amarillo & Revisión periódica \\
11 a 15 & \cellcolor{C_Alto} Alto & Naranja & Plan de mitigación prioritario \\
16 a 25 & \cellcolor{C_Critico}\textcolor{white}{\textbf{Crítico}} & Rojo & Acción inmediata \\
\bottomrule
\end{tabularx}
\end{table}

\subsection{Representación de la matriz}

La siguiente matriz de 5 por 5 ubica los 16 riesgos del proyecto según su evaluación de probabilidad e impacto y está coloreada según el nivel resultante.

\begin{table}[H]
\centering
\renewcommand{\arraystretch}{2}
\setlength{\tabcolsep}{4pt}
\makebox[\textwidth][c]{
\begin{tabular}{|c|l|P{3.2cm}|P{2.5cm}|P{3cm}|P{3cm}|P{2.8cm}|}
\hline
\multicolumn{2}{|c|}{\multirow{2}{*}{}} & \multicolumn{5}{c|}{\multirow{1}{*}{\textbf{Impacto}}} \\ \cline{3-7}
\multicolumn{2}{|c|}{} & \centering\textbf{Insignificante} \newline \textbf{1} & \centering\textbf{Menor} \newline \textbf{2} & \centering\textbf{Significativo} \newline \textbf{3} & \centering\textbf{Mayor} \newline \textbf{4} & \centering\arraybackslash\textbf{Severo} \newline \textbf{5} \\ \hline
\multirow{5}{*}[-3em]{\rotatebox{90}{\textbf{Probabilidad}}} 
& \textbf{5 Casi seguro} & \cellcolor{C_Medio}\centering Medio 5 & \cellcolor{C_Alto}\centering Alto 10 & \cellcolor{C_MuyAlto}\centering Muy alto 15 & \cellcolor{C_Critico}\centering\textcolor{white}{\textbf{Crítico 20}} & \cellcolor{C_Critico}\centering\arraybackslash\textcolor{white}{\textbf{Crítico 25}} \\ \hhline{~|-|-|-|-|-|-|}
& \textbf{4 Probable} & \cellcolor{C_Medio}\centering medio 4 & \cellcolor{C_Medio}\centering Medio 8 & \cellcolor{C_Alto}\centering Alto 12 (R-09, R-10) & \cellcolor{C_MuyAlto}\centering Muy alto 16 (R-03, R-05) & \cellcolor{C_Critico}\centering\arraybackslash\textcolor{white}{\textbf{Crítico 20 (R-01, R-02)}} \\ \hhline{~|-|-|-|-|-|-|}
& \textbf{3 Moderado} & \cellcolor{C_Bajo}\centering Bajo 3 & \cellcolor{C_Medio}\centering Medio 6 & \cellcolor{C_Medio}\centering Medio 9 (R-15, R-16) & \cellcolor{C_Alto}\centering Alto 12 (R-04, R-07, R-11, R-12, R-14) & \cellcolor{C_MuyAlto}\centering\arraybackslash Muy alto 15 (R-06) \\ \hhline{~|-|-|-|-|-|-|}
& \textbf{2 Poco probable} & \cellcolor{C_MuyBajo}\centering Muy Bajo 2 & \cellcolor{C_Bajo}\centering Bajo 4 & \cellcolor{C_Medio}\centering Medio 6 & \cellcolor{C_Medio}\centering Medio 8 (R-13) & \cellcolor{C_Alto}\centering\arraybackslash Alto 10 \\ \hhline{~|-|-|-|-|-|-|}
& \textbf{1 Raro} & \cellcolor{C_MuyBajo}\centering Muy bajo 1 & \cellcolor{C_MuyBajo}\centering Muy Bajo 2 & \cellcolor{C_Bajo}\centering Bajo 3 & \cellcolor{C_Medio}\centering Medio 4 & \cellcolor{C_Medio}\centering\arraybackslash Medio 5 \\ \hline
\end{tabular}
}
\end{table}

\textbf{Leyenda de colores}

\begin{table}[H]
\centering
\renewcommand{\arraystretch}{1.5}
\makebox[\textwidth][c]{
\begin{tabular}{|c|c|c|c|c|c|}
\hline
\cellcolor{C_MuyBajo} Muy bajo (1-2) & \cellcolor{C_Bajo} Bajo (3-4) & \cellcolor{C_Medio} Medio (4-9) & \cellcolor{C_Alto} Alto (10-14) & \cellcolor{C_MuyAlto} Muy alto (15-19) & \cellcolor{C_Critico}\textcolor{white}{Crítico (20-25)} \\ \hline
\end{tabular}
}
\end{table}

\subsection{Tabla de ubicación de riesgos}

La siguiente tabla detalla la posición exacta de cada riesgo dentro de la matriz.

\begin{table}[H]
\centering
\begin{tabularx}{\textwidth}{c Y c c c c}
\toprule
\textbf{Código} & \textbf{Riesgo} & \textbf{Prob.} & \textbf{Imp.} & \textbf{Nivel} & \textbf{Clasificación} \\
\midrule
R-01 & Restricción fuerte de tiempo & 4 & 5 & 25 & \cellcolor{C_Critico}\textcolor{white}{\textbf{Crítico}} \\
R-02 & Alcance creciente o scope creep & 4 & 5 & 20 & \cellcolor{C_Critico}\textcolor{white}{\textbf{Crítico}} \\
R-03 & Pérdida de coherencia en el informe final & 4 & 4 & 16 & \cellcolor{C_MuyAlto} Muy alto \\
R-04 & Subestimación del esfuerzo de la interfaz & 3 & 4 & 12 & \cellcolor{C_Alto} Alto \\
R-05 & Disponibilidad limitada del equipo & 4 & 4 & 16 & \cellcolor{C_MuyAlto} Muy alto \\
R-06 & Errores críticos de navegación en el prototipo & 3 & 5 & 15 & \cellcolor{C_MuyAlto} Muy alto \\
R-07 & Baja usabilidad del prototipo & 3 & 4 & 12 & \cellcolor{C_Alto} Alto \\

R-09 & Problemas con el repositorio o pérdida de versiones & 4 & 3 & 12 & \cellcolor{C_Alto} Alto \\
R-10 & Cambios tardíos cerca de la entrega & 4 & 3 & 12 & \cellcolor{C_Alto} Alto \\
R-11 & Falta de criterios claros de aceptación & 3 & 4 & 12 & \cellcolor{C_Alto} Alto \\
R-12 & Mala comunicación entre integrantes & 3 & 4 & 12 & \cellcolor{C_Alto} Alto \\
R-13 & Exceso del presupuesto planificado & 2 & 4 & 8 & \cellcolor{C_Medio} Medio \\
R-14 & Baja validación de stakeholders externos & 3 & 4 & 12 & \cellcolor{C_Alto} Alto \\
R-15 & Expectativa de funcionalidades no incluidas & 3 & 3 & 9 & \cellcolor{C_Medio} Medio \\
R-16 & Falta de evidencia final organizada & 3 & 3 & 9 & \cellcolor{C_Medio} Medio \\
\bottomrule
\end{tabularx}
\end{table}

\subsection{Análisis de concentración de riesgos}

La distribución de los 16 riesgos dentro de la matriz revela un patrón de concentración importante para la toma de decisiones del equipo.

\begin{table}[H]
\centering
\begin{tabularx}{\textwidth}{c c Y}
\toprule
\textbf{Clasificación} & \textbf{Cantidad} & \textbf{Observación} \\
\midrule
\cellcolor{C_Critico}\textcolor{white}{\textbf{Crítico}} & 2 & Concentrados en el cuadrante superior derecho de la matriz y relacionados con tiempo y alcance. \\
\cellcolor{C_MuyAlto} Muy alto & 3 & Se posicionan en la celda Probabilidad 4 con Impacto 4 y relacionados con documentación además de desarrollo de la interfaz y disponibilidad del equipo. \\
\cellcolor{C_Alto} Alto & 7 & La mayoría se agrupa en la celda Probabilidad 3 con Impacto 4 y Probabilidad 4 con Impacto 4 vinculados a calidad así como esfuerzo de la interfaz y herramientas y comunicación y cambios. \\
\cellcolor{C_Medio} Medio & 4 & Relacionados con presupuesto y posibilidad de fallas en herramientas o falta de evidencia final y funcionalidad no incluidas siendo no urgentes pero deben mantenerse visibles. \\
\cellcolor{C_Bajo} Bajo & 0 & No se identificaron riesgos de nivel bajo. \\
\cellcolor{C_MuyBajo} Muy bajo & 0 & No se identificaron riesgos de nivel muy bajo y esto refleja un proyecto con restricciones reales de tiempo y alcance. \\
\bottomrule
\end{tabularx}
\end{table}

La ausencia de riesgos clasificados como bajos y muy bajos confirma que el proyecto opera bajo condiciones de presión real y un plazo fijo de entrega y múltiples entregables documentales y un prototipo técnico que debe completarse en paralelo. Esto refuerza la importancia de dar seguimiento constante a los dos riesgos críticos y los tres muy altos durante toda la ejecución del proyecto.

\newpage
% --- Entregable 16 ---
\section{Plan de adquisiciones}

La presente sección corresponde al Plan de Adquisiciones del proyecto Plataforma Universitaria Integrada, desarrollado para la Universidad Tecnológica La Mejor. Este plan forma parte de la planificación del proyecto y tiene como finalidad definir los recursos externos, herramientas, servicios, materiales y apoyos operativos que deben considerarse para ejecutar el trabajo de manera ordenada durante el periodo planificado de tres meses.

El proyecto consiste en la elaboración de un prototipo web navegable y de un conjunto de documentos de administración de proyectos. Debido a que el alcance aprobado excluye el desarrollo de backend funcional, base de datos real, autenticación institucional, pasarela de pagos real y despliegue productivo obligatorio, las adquisiciones no se orientan a comprar una solución empresarial completa, sino a sostener la elaboración del prototipo, la documentación, la comunicación del equipo, el cronograma, las pruebas y las evidencias finales.

El plan se relaciona directamente con el cronograma elaborado en Microsoft Project, el plan de recursos, la estimación de costos y el plan de calidad. En consecuencia, las adquisiciones se manejan de manera controlada para evitar gastos innecesarios y asegurar que el presupuesto total del proyecto se mantenga en ₡17,625,200. El enfoque utilizado prioriza herramientas disponibles, licencias temporales, servicios de bajo costo y recursos ya existentes, ya que el proyecto tiene naturaleza académica y no requiere una contratación institucional completa.

\subsection{Propósito del plan de adquisiciones}

El propósito de esta sección es establecer qué elementos deben adquirirse, contratarse, utilizarse o reservarse para apoyar el desarrollo del proyecto. Asimismo, permite definir criterios básicos para seleccionar recursos, asignar responsables, controlar costos y mantener trazabilidad entre las necesidades del proyecto y el presupuesto aprobado.

Este plan busca cumplir los siguientes objetivos:

\begin{itemize}[leftmargin=*, label={\textbullet}]
    \item Identificar los recursos externos o complementarios requeridos por el proyecto.
    \item Diferenciar entre recursos ya disponibles y recursos que podrían representar un costo.
    \item Establecer criterios simples para seleccionar herramientas, servicios y materiales.
    \item Relacionar las adquisiciones con el cronograma, el presupuesto y el alcance aprobado.
    \item Evitar compras innecesarias o elementos fuera del alcance del prototipo académico.
    \item Asignar responsables para la revisión, aprobación y seguimiento de cada adquisición.
\end{itemize}

\subsection{Alcance de las adquisiciones}

Las adquisiciones consideradas en este proyecto se limitan a recursos de apoyo para la gestión, documentación, desarrollo y presentación final. No se contempla la compra de un ERP universitario, contratación de servidores productivos definitivos, licencias empresariales de gran escala ni servicios profesionales externos para desarrollar el prototipo.

\begin{table}[H]
\centering
\caption{Alcance de las adquisiciones}
\label{tab:alcance-adquisiciones}
{\small
\begin{tabularx}{\textwidth}{p{4cm} Y}
\toprule
\textbf{Tipo de recurso} & \textbf{Alcance definido} \\
\midrule
Software y herramientas & Licencias temporales o herramientas de apoyo para cronograma, documentación, diseño, desarrollo y control de versiones. \\
Hardware y materiales & Recursos físicos mínimos, almacenamiento externo y periféricos complementarios si fueran requeridos. \\
Servicios operativos & Conexión a internet, impresiones, transporte a reuniones o servicios básicos de apoyo al trabajo del equipo. \\
Hosting o infraestructura & Uso de alternativas gratuitas o de bajo costo para respaldo, repositorio y posible demostración del prototipo. \\
Contingencia & Reserva para atender imprevistos relacionados con herramientas, materiales, conectividad o necesidades menores de ejecución. \\
\bottomrule
\end{tabularx}
}
\end{table}


\subsection{Análisis Make vs Buy}

Para este proyecto se aplica un análisis Make vs Buy con el fin de definir qué elementos serán desarrollados por el equipo interno y cuáles recursos deberán adquirirse o utilizarse como apoyo externo. Debido a que el alcance corresponde a un prototipo académico Front-End navegable y no a un sistema productivo completo, la estrategia principal será desarrollar internamente los entregables técnicos y documentales, mientras que las adquisiciones se limitarán a herramientas, servicios y materiales de apoyo.

\begin{table}[H]
\centering
\caption{Análisis Make vs Buy del proyecto}
\label{tab:make-buy}
{\small
\begin{tabularx}{\textwidth}{p{4cm} p{3cm} Y}
\toprule
\textbf{Elemento} & \textbf{Decisión} & \textbf{Justificación} \\
\midrule
Prototipo Front-End navegable & Desarrollar internamente & El equipo cuenta con desarrolladores asignados y el alcance exige construir una interfaz web navegable, no comprar un sistema comercial. \\
Documentación del proyecto & Desarrollar internamente & Los entregables forman parte de la evaluación académica y deben ser elaborados por el equipo según los contenidos vistos en clase. \\
Cronograma, WBS, RACI, riesgos, calidad y cambios & Desarrollar internamente & Estos documentos representan la aplicación práctica de la administración de proyectos y deben mantener coherencia con el alcance aprobado. \\
Microsoft Project y herramientas de documentación & Adquirir o utilizar como apoyo & Se requieren para elaborar cronograma, evidencias, documentos y archivos finales, pero no sustituyen el trabajo del equipo. \\
Hosting o infraestructura de demostración & Utilizar apoyo externo básico & Puede usarse una alternativa gratuita o de bajo costo únicamente para mostrar evidencia del prototipo, sin implicar despliegue productivo. \\
ERP universitario, backend, base de datos real o pasarela de pagos & No adquirir & Estos elementos están fuera del alcance aprobado y aumentarían innecesariamente el costo, el tiempo y la complejidad del proyecto. \\
\bottomrule
\end{tabularx}
}
\end{table}

Con base en este análisis, la decisión general del proyecto es desarrollar internamente el prototipo y la documentación, y adquirir únicamente recursos complementarios de bajo costo. Esta estrategia permite mantener el presupuesto total de ₡17,625,200 y evita compras innecesarias que podrían generar alcance creciente o contradicciones con la definición del alcance.


\subsection{Criterios de selección}

Para seleccionar cualquier recurso o servicio se utilizarán criterios simples y coherentes con la naturaleza del proyecto. Estos criterios permiten mantener el control financiero y evitar que una adquisición altere el alcance o genere dependencia innecesaria.

\begin{table}[H]
\centering
\caption{Criterios de selección de adquisiciones}
\label{tab:criterios-adquisiciones}
{\small
\begin{tabularx}{\textwidth}{p{4cm} Y}
\toprule
\textbf{Criterio} & \textbf{Aplicación en el proyecto} \\
\midrule
Costo razonable & La adquisición debe mantenerse dentro del presupuesto aprobado y preferiblemente utilizar opciones gratuitas o temporales. \\
Necesidad para el alcance & Solo se aceptan recursos relacionados con documentación, cronograma, prototipo Front-End, pruebas o evidencias finales. \\
Facilidad de uso & La herramienta o servicio debe poder ser utilizado por el equipo sin capacitación extensa. \\
Disponibilidad inmediata & Se priorizan recursos que puedan ser utilizados durante el ciclo del proyecto sin procesos largos de contratación. \\
Compatibilidad técnica & Las herramientas deben ser compatibles con tecnologías web, repositorio de versiones, PDF, Microsoft Project o documentación académica. \\
Trazabilidad & Las decisiones relevantes deben quedar documentadas en el repositorio, minuta, correo o archivo de seguimiento cuando corresponda. \\
\bottomrule
\end{tabularx}
}
\end{table}

\subsection{Plan de adquisiciones del proyecto}

La siguiente tabla presenta el plan principal de adquisiciones. Los montos se mantienen alineados con la estimación PERT y con el presupuesto actualizado del proyecto.

\begin{table}[H]
\centering
\caption{Plan de adquisiciones}
\label{tab:plan-adquisiciones}
{\scriptsize
\begin{tabularx}{\textwidth}{p{2.6cm} p{3.3cm} Y p{2.4cm} p{2.1cm}}
\toprule
\textbf{Categoría} & \textbf{Recurso} & \textbf{Uso previsto} & \textbf{Responsable} & \textbf{Costo estimado} \\
\midrule
Software & Microsoft Project & Elaborar cronograma, tareas, dependencias, hitos, recursos, costos y ruta crítica. & Isaac Jiménez Blanco & ₡66,667 \\
Software & Office 365 o herramientas equivalentes & Elaboración y revisión de documentos, tablas, evidencias y archivos finales en PDF. & Equipo interno & ₡131,667 \\
Software & Herramientas premium opcionales & Apoyo visual o utilidades menores en caso de requerirse durante la documentación o presentación. & Tamara Robles Camacho & ₡23,333 \\
Hardware & Almacenamiento externo & Respaldo local de archivos finales, cronograma, evidencias y documentos importantes. & Isaac Jiménez Blanco & ₡19,167 \\
Hardware & Periféricos complementarios & Apoyo operativo mínimo para trabajo local, revisiones o entrega de evidencias. & Equipo interno & ₡14,167 \\
Servicios & Conexión a internet & Comunicación, repositorio, reuniones virtuales, investigación, desarrollo y carga de evidencias. & Cada integrante & ₡359,167 \\
Servicios & Impresiones y material físico & Material de apoyo o evidencias impresas si la entrega académica lo solicita. & Tamara Robles Camacho & ₡10,500 \\
Servicios & Transporte a reuniones presenciales & Traslados mínimos en caso de reuniones de coordinación o revisión presencial. & Equipo interno & ₡41,333 \\
\bottomrule
\end{tabularx}
}
\end{table}

El subtotal de rubros adicionales asociados a software, hardware y servicios corresponde a ₡666,000. Además, el proyecto considera una reserva de contingencia de ₡1,515,867 para atender imprevistos sin modificar el presupuesto total autorizado.

\subsection{Responsables y aprobación}

Las adquisiciones menores podrán ser revisadas por el equipo interno, pero cualquier gasto que afecte el presupuesto total, el alcance, el cronograma o la calidad deberá ser validado por la directora del proyecto. Si una adquisición representa un cambio relevante o genera una variación presupuestaria, deberá tratarse mediante el proceso formal de control de cambios.

\begin{table}[H]
\centering
\caption{Responsables del plan de adquisiciones}
\label{tab:responsables-adquisiciones}
{\small
\begin{tabularx}{\textwidth}{p{4cm} Y}
\toprule
\textbf{Rol} & \textbf{Responsabilidad} \\
\midrule
Tamara Robles Camacho & Validar necesidades de adquisición, controlar que los gastos se mantengan dentro del presupuesto y aprobar ajustes relevantes. \\
Isaac Jiménez Blanco & Mantener trazabilidad documental de herramientas, costos, evidencias y relación con el cronograma del proyecto. \\
Jorge Abarca Quiros & Apoyar la revisión de costos, adquisiciones menores y coherencia con la estimación PERT. \\
Mauricio González Prendas & Informar necesidades técnicas del desarrollo Front-End, repositorio, pruebas y herramientas relacionadas con los portales asignados. \\
Carlos Sainz Arce & Informar necesidades técnicas del portal financiero, dashboard ejecutivo y pruebas de calidad del prototipo. \\
\bottomrule
\end{tabularx}
}
\end{table}

\subsection{Gestión y seguimiento de adquisiciones}

El seguimiento de las adquisiciones se realizará durante la planificación, ejecución y cierre del proyecto. Para ello se verificará que los recursos se utilicen únicamente para apoyar las tareas definidas en el cronograma y que no generen nuevas obligaciones fuera del alcance aprobado.

\begin{itemize}[leftmargin=*, label={\textbullet}]
    \item Revisar que las herramientas utilizadas estén disponibles antes de iniciar las tareas críticas.
    \item Confirmar que Microsoft Project permita generar evidencias del cronograma y la ruta crítica.
    \item Mantener respaldos de documentos, código y archivos finales en el repositorio del proyecto.
    \item Registrar cualquier gasto relevante para mantener coherencia con la estimación de costos.
    \item Usar alternativas gratuitas siempre que cumplan con los requisitos del proyecto.
    \item Evitar contrataciones o compras que impliquen backend, base de datos real o operación productiva del sistema.
\end{itemize}

\subsection{Riesgos asociados a las adquisiciones}

Las adquisiciones del proyecto se relacionan con riesgos moderados debido a que dependen de herramientas tecnológicas y servicios disponibles durante el periodo de trabajo. Los principales riesgos son pérdida de acceso a herramientas, problemas con Microsoft Project, dificultad para exportar evidencias, fallos en el repositorio o inconsistencias entre costos documentados.

Para controlar estos riesgos, el equipo mantendrá copias locales de los archivos importantes, respaldos en GitHub, exportaciones en PDF y revisión cruzada de los montos reportados en documentos como cronograma, plan de calidad, registro de riesgos y estimación de costos.

\newpage
% --- Entregable 17 ---
\section{KPIs del Proyecto}

En esta sección se desglosan los indicadores clave de desempeño (KPI, por sus siglas en ingles), estos ayudan a medir de forma objetiva el avance y la salud del proyecto Plataforma Universitaria Integrada a lo largo de sus cinco fases. En las siguientes partes de esta sección se definen los indicadores distribuidos en las cinco categorías exigidas: tiempo, costo, calidad, riesgos y productividad, siguiendo el criterio SMART para garantizar que cada KPI sea medible y útil para la toma de decisiones del equipo.

\subsection{Criterio SMART aplicado a los KPIs}

Cada indicador clave de desempeño descritos cumple con las siguientes cinco características:

\begin{table}[H]
\centering
\small
\caption{Criterio SMART}\label{tab:criterio-smart}
\begin{tabularx}{\textwidth}{p{4.5cm} Y}
\toprule
\textbf{Criterio} & \textbf{Descripción} \\
\midrule
\textbf{Específico (Specific)} & El KPI está claramente definido y enfocado en un aspecto particular del proyecto. \\
\textbf{Medible (Measurable)} & Es posible cuantificar el KPI utilizando datos concretos del cronograma, presupuesto o registro de riesgos. \\
\textbf{Alcanzable (Achievable)} & La meta del KPI es realista dentro del contexto de un proyecto académico de 3 meses con 5 integrantes. \\
\textbf{Relevante (Relevant)} & El KPI está directamente relacionado con los objetivos del proyecto y es importante para su éxito. \\
\textbf{Temporal (Time-bound)} & Existe un marco de tiempo específico para la medición y el logro del KPI (semanal, por fase o al cierre). \\
\bottomrule
\end{tabularx}
\end{table}

\subsection{Indicadores de tiempo}

Los siguientes indicadores miden el cumplimiento del cronograma y la eficiencia en el uso del tiempo planificado del proyecto.

El Valor Planificado (PV) es el valor del trabajo que se debería haber completado en un momento dato, sale del cronograma y el presupuesto, por ejemplo: si una tarea vale ₡100,000 y según el cronograma se debe ir al 50\%, entonces el PV =100,000 $\times$ 0.5= ₡50,000.

En el caso del Valor Ganado (EV), sale del trabajo que realmente se ha terminado, pero expresado en dinero, por ejemplo: si la tarea vale ₡100,000 y en realidad solo se lleve el 40\%, entonces el EV= 100,000 $\times$ 0,4 =₡40,000.

\begin{table}[H]
\centering
\small
\caption{Indicadores de tiempo}\label{tab:indicadores-tiempo}
\begin{tabularx}{\textwidth}{p{4.5cm} Y c c}
\toprule
\textbf{KPI} & \textbf{Qué mide / Fórmula} & \textbf{Meta} & \textbf{Frecuencia} \\
\midrule
\textbf{KPI-01 Cumplimiento del cronograma} & Mide el porcentaje de tareas que se completan dentro del plazo previsto. Fórmula: (Tareas completadas a tiempo $\div$ Total de tareas) $\times$ 100 & $\geq$ 90\% & Semanal \\
\textbf{KPI-02 Índice de Desempeño del Cronograma (SPI)} & Evalúa la eficiencia en el uso del tiempo del proyecto. Fórmula: SPI = Valor Ganado (EV) $\div$ Valor Planificado (PV) & $\geq$ 0.95 & Quincenal \\
\textbf{KPI-03 Variación del cronograma (SV)} & Calcula la diferencia entre el trabajo planificado y el trabajo realizado en términos de tiempo. Fórmula: SV = Valor Ganado (EV) $-$ PV(Valor Planificado) & SV $\geq$ 0 & Por fase \\
\bottomrule
\end{tabularx}
\end{table}

\subsection{Indicadores de costo}

Estos indicadores controlan que el gasto real del proyecto se mantenga alineado con el presupuesto planificado de ₡17,625,200.

\begin{table}[H]
\centering
\small
\caption{Indicadores de costo}\label{tab:indicadores-costo}
\begin{tabularx}{\textwidth}{p{4.5cm} Y c c}
\toprule
\textbf{KPI} & \textbf{Qué mide / Fórmula} & \textbf{Meta} & \textbf{Frecuencia} \\
\midrule
\textbf{KPI-04 Cumplimiento del presupuesto} & Mide el porcentaje del presupuesto que se ha gastado en relación con lo planificado. Fórmula: (Costo Real $\div$ Presupuesto Planificado) $\times$ 100 & $\leq$ 100\% & Por fase \\
\textbf{KPI-05 Índice de Desempeño de Costos (CPI)} & Mide la relación entre el valor ganado y el costo real del proyecto. Fórmula: CPI = Valor Ganado (EV) $\div$ Costo Real (AC) & $\geq$ 1.0 & Quincenal \\
\textbf{KPI-06 Variación del costo (CV)} & Calcula la diferencia entre el costo presupuestado del trabajo realizado y el costo real. Fórmula: CV = Valor Ganado (EV) $-$ Costo Real (AC) & CV $\geq$ 0 & Por fase \\
\bottomrule
\end{tabularx}
\end{table}

\subsection{Indicadores de calidad}

Los indicadores de calidad evalúan la calidad técnica del prototipo Front-End y la coherencia de los entregables que se documentan.

\begin{table}[H]
\centering
\small
\caption{Indicadores de calidad}\label{tab:indicadores-calidad}
\begin{tabularx}{\textwidth}{p{4.5cm} Y p{3cm} p{2.5cm}}
\toprule
\textbf{KPI} & \textbf{Qué mide / Fórmula} & \textbf{Meta} & \textbf{Frecuencia} \\
\midrule
\textbf{KPI-07 Defectos encontrados por módulo} & Mide la cantidad de errores o fallos identificados durante las pruebas del prototipo en cada módulo (Login, Portales, Dashboard). Fórmula: Número de defectos detectados por módulo & $\leq$ 3 defectos críticos por módulo & Por entrega de módulo \\
\textbf{KPI-08 Cobertura de criterios de aceptación} & Indica el porcentaje de criterios de aceptación del WBS y Plan de Calidad que se cumplen en cada entregable. Fórmula: (Criterios cumplidos $\div$ Criterios totales) $\times$ 100 & $\geq$ 95\% & Por entregable \\
\textbf{KPI-09 Satisfacción del equipo evaluador (usabilidad)} & Evalúa el nivel de claridad y navegación del prototipo mediante revisión interna del equipo antes de la entrega final. Fórmula: Promedio de calificación (escala de 1 a 5) por módulo & $\geq$ 4.0 / 5.0 & Antes de cada entrega \\
\textbf{KPI-10 Porcentaje de pruebas exitosas} & Mide la cantidad de pruebas completas realizadas de manera exitosa en el prototipo de la interfaz. Formula: (Pruebas aprobadas $\div$ total de pruebas) $\times$ 100 & $\geq$ 90\% & Por entrega \\
\bottomrule
\end{tabularx}
\end{table}

\subsection{Indicadores de riesgos}

Los indicadores de riesgo son muy útiles, ya que monitorean el comportamiento de los 13 riesgos identificados en el registro de riesgos a lo largo del proyecto.

\begin{table}[H]
\centering
\small
\caption{Indicadores de riesgos}\label{tab:indicadores-riesgos}
\begin{tabularx}{\textwidth}{p{4.5cm} Y p{3cm} p{2.5cm}}
\toprule
\textbf{KPI} & \textbf{Qué mide / Fórmula} & \textbf{Meta} & \textbf{Frecuencia} \\
\midrule
\textbf{KPI-11 Riesgos críticos abiertos} & Cuenta el número de riesgos clasificados como Crítico que permanecen en estado Abierto o En seguimiento. Fórmula: Número de riesgos críticos sin estrategia de mitigación aplicada & 0 al cierre de cada fase & Semanal \\
\textbf{KPI-12 Riesgos mitigados vs. identificados} & Mide la cantidad de riesgos que han recibido una acción de respuesta efectiva sobre el total identificado. Fórmula: (Riesgos mitigados $\div$ Total de riesgos) $\times$ 100 & $\geq$ 80\% & Por fase \\
\textbf{KPI-13 Tiempo de respuesta a riesgos} & Evalúa el tiempo desde el momento en el que un riesgo se clasifico como crítico, muy alto, alto, medio, bajo o muy bajo hasta el momento de su acción. Formula: Promedio de días desde identificación hasta acción & $\leq$ 3 días & Semanal \\
\bottomrule
\end{tabularx}
\end{table}

\subsection{Indicadores de productividad}

Los indicadores de productividad miden la eficiencia con la que el equipo de 5 integrantes avanza en la entrega de los entregables del proyecto.

\begin{table}[H]
\centering
\small
\caption{Indicadores de productividad}\label{tab:indicadores-productividad}
\begin{tabularx}{\textwidth}{p{4.5cm} Y c c}
\toprule
\textbf{KPI} & \textbf{Qué mide / Fórmula} & \textbf{Meta} & \textbf{Frecuencia} \\
\midrule
\textbf{KPI-14 Entregables completados por semana} & Mide la velocidad de entrega de documentos y avances del prototipo en relación con lo planificado en el cronograma. Fórmula: Número de entregables o paquetes de trabajo completados $\div$ semana & $\geq$ 2 entregables / semana & Semanal \\
\bottomrule
\end{tabularx}
\end{table}

\subsection{Resumen de indicadores por categoría}

La siguiente tabla consolida los 14 KPIs definidos, agrupados por categoría:

\begin{table}[H]
\centering
\small
\caption{Resumen de indicadores por categoría}\label{tab:resumen-indicadores}
\begin{tabular}{l c l}
\toprule
\textbf{Categoría} & \textbf{Cantidad} & \textbf{Códigos} \\
\midrule
\textbf{Tiempo} & 3 & KPI-01, KPI-02, KPI-03 \\
\textbf{Costo} & 3 & KPI-04, KPI-05, KPI-06 \\
\textbf{Calidad} & 4 & KPI-07, KPI-08, KPI-09, KPI-10 \\
\textbf{Riesgos} & 3 & KPI-11, KPI-12, KPI-13 \\
\textbf{Productividad} & 1 & KPI-14 \\
\midrule
\textbf{TOTAL} & \textbf{14 indicadores} & \\
\bottomrule
\end{tabular}
\end{table}

\newpage
% --- Entregable 18 ---
\section{Desarrollo de la interfaz}

[Espacio reservado para la incorporación de información e imágenes del prototipo correspondientes a esta sección, el cual es responsabilidad de otros miembros del equipo].

\newpage
% --- Entregable 19 ---
\section{Proceso de control de cambios}

La presente sección define el Proceso de Control de Cambios del proyecto Plataforma Universitaria Integrada. Este proceso tiene como finalidad establecer la forma en que se recibirán, analizarán, aprobarán, rechazarán, documentarán y comunicarán los cambios que puedan afectar el alcance, el cronograma, los costos, la calidad, los riesgos o los entregables del proyecto.

El control de cambios es necesario porque el proyecto contiene varios documentos de gestión y un prototipo web navegable que debe mantenerse alineado con el alcance aprobado. Además, el equipo trabaja con roles combinados y un presupuesto definido mediante estimación PERT, por lo tanto cualquier modificación no controlada puede generar inconsistencias entre documentos, sobrecarga de trabajo, aumento de costos o reducción de calidad.

El proceso descrito en esta sección se adapta a la escala académica del proyecto. Por esta razón, no se plantea un comité corporativo complejo, sino un flujo formal y claro donde la directora del proyecto coordina la revisión, los analistas documentan el impacto, los desarrolladores estiman la viabilidad técnica y el patrocinador aprueba los cambios relevantes.

\subsection{Propósito del proceso}

El propósito del proceso de control de cambios es evitar modificaciones informales que afecten el avance del proyecto sin análisis previo. Asimismo, busca proteger las líneas base de alcance, cronograma, costo y calidad que fueron definidas durante la planificación.

El proceso persigue los siguientes objetivos:

\begin{itemize}[leftmargin=*, label={\textbullet}]
    \item Definir un flujo formal para recibir y evaluar solicitudes de cambio.
    \item Analizar el impacto de cada cambio en alcance, tiempo, costo, calidad y riesgos.
    \item Determinar si un cambio debe aprobarse, rechazarse o posponerse.
    \item Documentar las decisiones para mantener trazabilidad.
    \item Comunicar los cambios aprobados al equipo y a los interesados relevantes.
    \item Actualizar documentos, cronograma o prototipo cuando el cambio sea aprobado.
\end{itemize}

\subsection{Principios del control de cambios}

El proceso se basa en reglas simples que permiten mantener orden durante la ejecución del proyecto.

\begin{table}[H]
\centering
\caption{Principios del proceso de control de cambios}
\label{tab:principios-cambios}
{\small
\begin{tabularx}{\textwidth}{p{4cm} Y}
\toprule
\textbf{Principio} & \textbf{Aplicación en el proyecto} \\
\midrule
Formalidad & Todo cambio relevante debe documentarse mediante una solicitud de cambio o registro equivalente. \\
Análisis previo & Ningún cambio que afecte alcance, tiempo, costo, calidad o riesgos debe ejecutarse sin revisión de impacto. \\
Trazabilidad & Las decisiones deben quedar respaldadas en documentos, repositorio, correo, minuta o evidencia del proyecto. \\
Coherencia & Si el cambio modifica un documento base, deben actualizarse los entregables relacionados. \\
Responsabilidad & La directora del proyecto coordina la decisión y los responsables técnicos o documentales apoyan el análisis. \\
Control del alcance & No se deben agregar funciones fuera del alcance sin evaluar su efecto sobre el prototipo y el cronograma. \\
\bottomrule
\end{tabularx}
}
\end{table}

\subsection{Tipos de cambios considerados}

Para este proyecto se consideran cambios en diferentes áreas de gestión y desarrollo. La clasificación permite determinar quién debe revisarlos y qué documentos podrían verse afectados.

\begin{table}[H]
\centering
\caption{Tipos de cambios del proyecto}
\label{tab:tipos-cambios}
{\small
\begin{tabularx}{\textwidth}{p{3.5cm} Y p{3.2cm}}
\toprule
\textbf{Tipo de cambio} & \textbf{Descripción} & \textbf{Ejemplo} \\
\midrule
Alcance & Agrega, elimina o modifica funcionalidades, módulos, entregables o límites del proyecto. & Incorporar una nueva pantalla al Portal Estudiantil. \\
Cronograma & Cambia fechas, duraciones, dependencias, hitos o ruta crítica. & Mover pruebas de QA por atraso en desarrollo. \\
Costo & Afecta el presupuesto, rubros de adquisición o uso de contingencia. & Usar parte de la reserva para una herramienta necesaria. \\
Calidad & Modifica criterios de aceptación, pruebas, estándares o nivel esperado del producto. & Agregar una revisión de navegación antes del cierre. \\
Riesgo & Incrementa o reduce la probabilidad o impacto de riesgos registrados. & Reducir alcance para bajar el riesgo de atraso. \\
Documentación & Requiere ajustar planes, matrices, registros o evidencias finales. & Cambiar presupuesto oficial de documentos relacionados. \\
\bottomrule
\end{tabularx}
}
\end{table}

\subsection{Roles y responsabilidades}

El control de cambios involucra a los integrantes internos y, cuando corresponde, al patrocinador o interesados externos. La participación depende del impacto del cambio solicitado.

\begin{table}[H]
\centering
\caption{Responsables del proceso de control de cambios}
\label{tab:responsables-cambios}
{\small
\begin{tabularx}{\textwidth}{p{4cm} Y}
\toprule
\textbf{Rol} & \textbf{Responsabilidad} \\
\midrule
Tamara Robles Camacho & Recibir solicitudes, coordinar revisión, validar impacto general y aprobar cambios menores o elevar cambios mayores. \\
Isaac Jiménez Blanco & Documentar solicitudes, revisar impacto en alcance, comunicaciones, calidad, riesgos y entregables relacionados. \\
Jorge Abarca Quiros & Apoyar el análisis de costos, riesgos, cronograma y consistencia documental. \\
Mauricio González Prendas & Evaluar impacto técnico en portal estudiantil, portal docente, portal administrativo, navegación y QA. \\
Carlos Sainz Arce & Evaluar impacto técnico en portal financiero, dashboard ejecutivo del sistema, navegación y QA. \\
Dr. Roberto Mora Fallas & Aprobar cambios relevantes que afecten alcance principal, costo total, fecha final o entregables estratégicos. \\
\bottomrule
\end{tabularx}
}
\end{table}

\subsection{Flujo del proceso de control de cambios}

El proceso se organiza en siete pasos secuenciales. Estos pasos buscan que cada cambio sea analizado antes de modificar el prototipo, el cronograma o la documentación.

\begin{table}[H]
\centering
\caption{Flujo del proceso de control de cambios}
\label{tab:flujo-cambios}
{\small
\begin{tabularx}{\textwidth}{p{1.4cm} p{3.3cm} Y}
\toprule
\textbf{Paso} & \textbf{Actividad} & \textbf{Descripción} \\
\midrule
1 & Identificación & Un integrante o interesado detecta una necesidad de modificación, problema o mejora. \\
2 & Registro de solicitud & Se completa una solicitud de cambio con fecha, solicitante, descripción, justificación y prioridad. \\
3 & Revisión preliminar & La directora del proyecto verifica si el cambio es válido, duplicado, urgente o fuera del alcance. \\
4 & Análisis de impacto & El equipo analiza efecto en alcance, cronograma, costo, calidad, riesgos, recursos y entregables. \\
5 & Decisión & El cambio se aprueba, rechaza, pospone o se aprueba con condiciones. \\
6 & Implementación & Si se aprueba, se actualiza el prototipo, cronograma, documentos o evidencia correspondiente. \\
7 & Comunicación y cierre & La decisión y el resultado se comunican al equipo y se deja trazabilidad final. \\
\bottomrule
\end{tabularx}
}
\end{table}

\subsection{Criterios de evaluación de impacto}

Cada solicitud debe evaluarse con los mismos criterios para evitar decisiones subjetivas. El análisis puede ser breve, pero debe cubrir las áreas principales del proyecto.

\begin{table}[H]
\centering
\caption{Criterios de evaluación de cambios}
\label{tab:criterios-evaluacion-cambios}
{\small
\begin{tabularx}{\textwidth}{p{3.5cm} Y}
\toprule
\textbf{Área} & \textbf{Preguntas de evaluación} \\
\midrule
Alcance & ¿Agrega o elimina módulos, pantallas, entregables o requisitos? ¿Contradice el alcance excluido? \\
Cronograma & ¿Afecta tareas, duraciones, dependencias, hitos o ruta crítica del cronograma en Microsoft Project? \\
Costo & ¿Usa presupuesto adicional, cambia rubros de adquisición o requiere aplicar contingencia? \\
Calidad & ¿Afecta criterios de aceptación, pruebas, navegación, usabilidad o consistencia documental? \\
Riesgos & ¿Aumenta o reduce riesgos como alcance creciente, retrasos, errores de navegación o baja evidencia final? \\
Recursos & ¿Sobrecarga a Tamara, Isaac, Jorge, Mauricio o Carlos? ¿Requiere apoyo externo? \\
Comunicación & ¿Debe informarse al patrocinador, stakeholders externos o todo el equipo? \\
\bottomrule
\end{tabularx}
}
\end{table}

\subsection{Decisiones posibles}

La evaluación de una solicitud puede producir cuatro decisiones. Esta clasificación permite mantener claridad y evitar cambios ambiguos.

\begin{table}[H]
\centering
\caption{Decisiones posibles ante una solicitud de cambio}
\label{tab:decisiones-cambios}
{\small
\begin{tabularx}{\textwidth}{p{3cm} Y}
\toprule
\textbf{Decisión} & \textbf{Descripción} \\
\midrule
Aprobado & El cambio se acepta y se implementa en el prototipo, documento o cronograma correspondiente. \\
Aprobado con condiciones & El cambio se acepta, pero con límites de alcance, fecha, costo o responsable. \\
Rechazado & El cambio no se implementa porque afecta demasiado el alcance, costo, tiempo o calidad. \\
Pospuesto & El cambio se considera valioso, pero se deja para una etapa posterior fuera del proyecto actual. \\
\bottomrule
\end{tabularx}
}
\end{table}

\subsection{Documentación y comunicación}

Toda solicitud aprobada debe dejar evidencia clara. Para cambios de bajo impacto puede bastar una anotación en el repositorio, una minuta o un mensaje formal de seguimiento. Para cambios que afecten alcance, cronograma, costo o calidad se debe utilizar la plantilla de solicitud de cambio y actualizar los documentos relacionados.

Los canales autorizados para comunicar cambios son GitHub, correo electrónico, documentos PDF, minutas de reunión y mensajes de coordinación cuando estos se usen solo como aviso. La mensajería instantánea no será suficiente para aprobar cambios importantes, aunque sí puede utilizarse para informar rápidamente que una solicitud debe ser revisada.

\subsection{Relación con otros documentos}

El proceso de control de cambios se relaciona con la definición del alcance, cronograma, plan de calidad, plan de comunicaciones, registro de riesgos, matriz de probabilidad e impacto, solicitudes de cambio y evaluación de cambios. Si un cambio modifica alguno de estos elementos, el documento correspondiente debe actualizarse para mantener la coherencia general del proyecto.

\newpage
% --- Entregable 20 ---
\section{Solicitud de cambio 1}

\subsection{Información general del cambio}

\begin{table}[H]
\centering
\begin{tabularx}{\textwidth}{p{5cm} Y}
\toprule
\textbf{Número de solicitud} & SC-01 \\
\textbf{Fecha de solicitud} & 18 de junio de 2026 \\
\textbf{Solicitante} & Lic. Patricia Vargas Soto y Prof. Andrea Solís Zamora \\
\textbf{Categoría del cambio} & Alcance y cronograma \\
\textbf{Prioridad} & Alta \\
\textbf{Descripción breve} & Adición de módulo para solicitar certificados académicos en el portal estudiantil \\
\bottomrule
\end{tabularx}
\end{table}

\subsection{Descripción detallada del cambio solicitado}

La presente solicitud propone integrar una funcionalidad de solicitud de certificados académicos dentro del Portal Estudiantil debido a que los estudiantes requieren generar y descargar certificaciones de notas así como la carga académica y el orden de matrícula de manera centralizada. Esto implica agregar una nueva sección específica en el portal junto con una pantalla de seguimiento de las solicitudes para lo cual se requiere actualizar el diseño del Front-End y asegurar que la navegación permita acceder a esta nueva herramienta. Asimismo el cambio afecta directamente el alcance del proyecto ya que se introduce un módulo no contemplado inicialmente y requiere ajustar la arquitectura visual de la aplicación.

\subsection{Justificación del cambio}

El proyecto busca modernizar y centralizar los procesos académicos por lo tanto es necesario ofrecer a los estudiantes herramientas de autogestión que reduzcan la carga operativa de la administración universitaria. De esta manera el beneficio esperado se refleja en una disminución importante del tiempo de atención presencial en las ventanillas de servicio al estudiante ya que los usuarios podrán obtener sus comprobantes oficiales desde la misma plataforma web en cualquier momento. En consecuencia esta nueva funcionalidad incrementa el valor del producto final y mejora directamente la percepción de los usuarios sobre la eficiencia institucional.

\subsection{Análisis de impacto preliminar}

El impacto en el alcance se considera directo debido a que el equipo de desarrollo debe construir las pantallas de solicitud y el historial de certificaciones dentro del portal estudiantil. Esto implica que el cronograma debe ajustarse para incluir días adicionales de desarrollo y de pruebas de calidad para este nuevo módulo ya que la fecha de entrega original tiene un margen de contingencia limitado. Por lo tanto el costo estimado del proyecto presenta un incremento en el esfuerzo de la mano de obra del equipo de desarrollo y esto requiere usar una parte de la reserva de contingencia presupuestada inicialmente. Además de estos factores el riesgo asociado a la restricción de tiempo aumenta debido a que se añade carga de trabajo a los desarrolladores Mauricio González Prendas y Carlos Sainz Arce para lo cual se deberá gestionar de manera cuidadosa el tiempo disponible antes del 24 de junio de 2026.

\subsection{Aprobaciones}

\begin{table}[H]
\centering
\begin{tabularx}{\textwidth}{Y Y Y}
\toprule
\textbf{Solicitado por} & \textbf{Revisado por} & \textbf{Aprobado por} \\
\midrule
& & \\
& & \\
Lic. Patricia Vargas Soto & Tamara Robles Camacho & Dr. Roberto Mora Fallas \\
Directora Administrativa & Project Manager & Sponsor \\
\bottomrule
\end{tabularx}
\end{table}

\newpage
% --- Entregable 21 ---
\section{Solicitud de cambio 2}

\subsection{Información general del cambio}

\begin{table}[H]
\centering
\begin{tabularx}{\textwidth}{p{5cm} Y}
\toprule
\textbf{Número de solicitud} & SC-02 \\
\textbf{Fecha de solicitud} & 19 de junio de 2026 \\
\textbf{Solicitante} & Tamara Robles Camacho \\
\textbf{Categoría del cambio} & Alcance y cronograma \\
\textbf{Prioridad} & Media \\
\textbf{Descripción breve} & Reducción del alcance del Dashboard Ejecutivo para compensar el esfuerzo de la solicitud SC-01 \\
\bottomrule
\end{tabularx}
\end{table}

\subsection{Descripción detallada del cambio solicitado}

Esta solicitud surge como una medida compensatoria derivada de la aprobación de la solicitud de cambio anterior debido a que el equipo requiere liberar carga de trabajo en el cronograma para asegurar la entrega final a tiempo. De esta manera se propone reducir el alcance visual del Dashboard Ejecutivo eliminando los indicadores financieros detallados y los indicadores académicos históricos ya que su desarrollo requiere un esfuerzo considerable que ahora se ha destinado al portal estudiantil. Por lo tanto el sistema mantendrá solamente los indicadores de estudiantes activos y los ingresos por matrícula en el dashboard principal para lo cual se ajustará el diseño visual y se simplificará la interfaz gráfica sin afectar el funcionamiento del resto de los portales.

\subsection{Justificación del cambio}

El proyecto mantiene una restricción fuerte de tiempo por lo tanto es fundamental equilibrar el esfuerzo de desarrollo y las expectativas de los stakeholders para evitar atrasos críticos en el ciclo de vida. Asimismo la eliminación de los indicadores históricos y financieros complejos permite que el equipo se enfoque en estabilizar los módulos principales y garantice la calidad de las funcionalidades que tienen un mayor impacto en la experiencia diaria de los estudiantes. En consecuencia el beneficio esperado de esta reducción es la optimización del tiempo de los desarrolladores ya que se libera la capacidad necesaria para completar las pantallas nuevas sin poner en peligro la fecha de cierre planificada para el 24 de junio de 2026.

\subsection{Análisis de impacto preliminar}

El impacto en el alcance se presenta como una reducción directa de las funcionalidades del dashboard ejecutivo debido a que solamente se mostrarán los datos de estudiantes activos y los ingresos por matrícula. Esto implica que el cronograma recupera días valiosos de desarrollo y de revisión de calidad ya que se elimina la necesidad de programar gráficos complejos y tablas de datos de años anteriores. Asimismo el costo total del proyecto se mantiene sin alteraciones debido a que el esfuerzo de trabajo simplemente se transfiere desde el dashboard hacia la construcción del nuevo módulo de certificados del portal estudiantil. Además de esto el cambio influye positivamente en el registro de riesgos debido a que disminuye la prioridad de los riesgos R-04 relacionado con la subestimación del esfuerzo del desarrollo y R-10 referente a los cambios tardíos cerca de la entrega final.

\subsection{Aprobaciones}

\begin{table}[H]
\centering
\begin{tabularx}{\textwidth}{Y Y Y}
\toprule
\textbf{Solicitado por} & \textbf{Revisado por} & \textbf{Aprobado por} \\
\midrule
& & \\
& & \\
Tamara Robles Camacho & Isaac Jiménez Blanco & Dr. Roberto Mora Fallas \\
Project Manager & Analista de Negocio & Sponsor \\
\bottomrule
\end{tabularx}
\end{table}

\newpage
% --- Entregable 22 ---
\section{Evaluación de cambios}

La presente sección corresponde a la Evaluación de Cambios del proyecto Plataforma Universitaria Integrada. Su finalidad es analizar las solicitudes de cambio registradas durante la ejecución y determinar su efecto sobre el alcance, el cronograma, los costos, la calidad, los riesgos, los recursos y los entregables del proyecto.

La evaluación de cambios permite comprobar que las modificaciones solicitadas no se ejecuten de manera informal ni generen contradicciones con la documentación aprobada. En este caso, se analizan dos solicitudes principales: la incorporación de un módulo de certificados académicos en el Portal Estudiantil y la reducción del alcance del Dashboard Ejecutivo del sistema como medida compensatoria.

El análisis se realiza bajo el proceso de control de cambios definido para el proyecto, tomando como referencia el alcance aprobado, el cronograma en Microsoft Project, el presupuesto total de ₡17,625,200, el registro de riesgos, la matriz de probabilidad e impacto y las solicitudes de cambio documentadas.

\subsection{Propósito de la evaluación}

El propósito de esta sección es revisar de forma estructurada el impacto de las solicitudes de cambio y dejar evidencia de la decisión tomada. La evaluación busca determinar si los cambios son viables, si generan valor para los usuarios, si pueden implementarse sin afectar la calidad del prototipo y si requieren ajustes en otros entregables.

Los objetivos específicos son los siguientes:

\begin{itemize}[leftmargin=*, label={\textbullet}]
    \item Identificar las solicitudes de cambio evaluadas.
    \item Analizar su impacto en alcance, tiempo, costo, calidad, riesgos y recursos.
    \item Determinar si cada cambio debe aprobarse, rechazarse o aprobarse con condiciones.
    \item Definir acciones de seguimiento para mantener el control del proyecto.
    \item Mantener coherencia entre solicitudes, cronograma, presupuesto y documentación final.
\end{itemize}

\subsection{Metodología de evaluación}

La metodología utilizada consiste en revisar cada solicitud contra los criterios definidos en el proceso de control de cambios. Para cada cambio se analiza su justificación, impacto, prioridad, responsables afectados, riesgos asociados y decisión final. Este análisis es cualitativo y se adapta al alcance académico del proyecto.

\begin{table}[H]
\centering
\caption{Criterios utilizados para evaluar los cambios}
\label{tab:criterios-evaluacion}
{\small
\begin{tabularx}{\textwidth}{p{3.5cm} Y}
\toprule
\textbf{Criterio} & \textbf{Descripción de la revisión} \\
\midrule
Alcance & Determina si el cambio agrega, reduce o modifica funcionalidades, pantallas o entregables. \\
Cronograma & Evalúa si el cambio afecta tareas, dependencias, hitos, ruta crítica o fecha final. \\
Costo & Revisa si el cambio incrementa costos, utiliza contingencia o se absorbe con reasignación interna. \\
Calidad & Determina si el cambio mejora o afecta la usabilidad, navegación, estabilidad o criterios de aceptación. \\
Riesgos & Analiza si el cambio aumenta o reduce riesgos identificados en el registro. \\
Recursos & Revisa si el cambio sobrecarga a los desarrolladores, analistas o directora del proyecto. \\
Decisión & Define si el cambio se aprueba, rechaza, pospone o aprueba con condiciones. \\
\bottomrule
\end{tabularx}
}
\end{table}

\subsection{Solicitudes de cambio evaluadas}

Durante la ejecución se documentaron dos solicitudes de cambio principales. Ambas se relacionan con el alcance del prototipo Front-End navegable y con la necesidad de mantener equilibrio entre valor agregado, tiempo disponible y calidad final.

\begin{table}[H]
\centering
\caption{Resumen de solicitudes de cambio evaluadas}
\label{tab:solicitudes-evaluadas}
{\small
\begin{tabularx}{\textwidth}{p{1.6cm} p{3cm} Y p{2.2cm}}
\toprule
\textbf{Código} & \textbf{Solicitante} & \textbf{Descripción resumida} & \textbf{Prioridad} \\
\midrule
SC-01 & Lic. Patricia Vargas Soto y Prof. Andrea Solís Zamora & Incorporar un módulo de solicitud de certificados académicos en el Portal Estudiantil. & Alta \\
SC-02 & Tamara Robles Camacho & Reducir el alcance del Dashboard Ejecutivo del sistema para compensar el esfuerzo agregado por SC-01. & Media \\
\bottomrule
\end{tabularx}
}
\end{table}

\subsection{Evaluación de la solicitud de cambio 1}

La solicitud SC-01 propone agregar una funcionalidad de solicitud de certificados académicos dentro del Portal Estudiantil. Esta funcionalidad permitiría representar la solicitud y consulta de certificaciones académicas, como constancias de notas, carga académica u otros comprobantes relacionados con el estudiante.

\begin{table}[H]
\centering
\caption{Evaluación de impacto de SC-01}
\label{tab:evaluacion-sc01}
{\small
\begin{tabularx}{\textwidth}{p{3.2cm} Y}
\toprule
\textbf{Área evaluada} & \textbf{Impacto identificado} \\
\midrule
Alcance & Aumenta el alcance del Portal Estudiantil al agregar una nueva sección no contemplada inicialmente. \\
Cronograma & Requiere ajustar tareas de desarrollo y pruebas del portal estudiantil, principalmente para diseño visual, navegación y revisión de calidad. \\
Costo & El cambio no modifica el presupuesto total si el esfuerzo se compensa con una reducción equivalente en otro módulo. \\
Calidad & Puede aumentar el valor percibido del prototipo para estudiantes, pero también exige revisar que la navegación siga siendo clara. \\
Riesgos & Incrementa temporalmente riesgos relacionados con restricción de tiempo, subestimación del esfuerzo de interfaz y cambios tardíos. \\
Recursos & Afecta principalmente a Mauricio González Prendas y Carlos Sainz Arce, con apoyo documental de Isaac Jiménez Blanco y Jorge Abarca Quiros. \\
\bottomrule
\end{tabularx}
}
\end{table}

\subsubsection{Decisión sobre SC-01}

La solicitud SC-01 se aprueba con condiciones, debido a que aporta valor al Portal Estudiantil y se relaciona con la necesidad institucional de reducir trámites presenciales. Sin embargo, su implementación debe limitarse a una representación visual navegable con datos simulados, sin conexión a base de datos, sin generación real de documentos oficiales y sin firma digital.

La condición principal es compensar el esfuerzo adicional mediante una reducción controlada en otra parte del alcance, para no aumentar el presupuesto total ni extender innecesariamente el cronograma.

\subsection{Evaluación de la solicitud de cambio 2}

La solicitud SC-02 propone reducir el alcance del Dashboard Ejecutivo del sistema. La reducción consiste en eliminar indicadores financieros detallados e indicadores académicos históricos, manteniendo únicamente indicadores generales como estudiantes activos e ingresos por matrícula.

\begin{table}[H]
\centering
\caption{Evaluación de impacto de SC-02}
\label{tab:evaluacion-sc02}
{\small
\begin{tabularx}{\textwidth}{p{3.2cm} Y}
\toprule
\textbf{Área evaluada} & \textbf{Impacto identificado} \\
\midrule
Alcance & Reduce el alcance funcional del Dashboard Ejecutivo del sistema al simplificar la cantidad de indicadores representados. \\
Cronograma & Libera tiempo de desarrollo, integración y pruebas para compensar el esfuerzo adicional del módulo de certificados. \\
Costo & No reduce ni aumenta el presupuesto total, ya que funciona como una reasignación interna de esfuerzo. \\
Calidad & Favorece la estabilidad del prototipo porque permite concentrar el trabajo en módulos principales y evita gráficos innecesarios. \\
Riesgos & Disminuye riesgos asociados a subestimación del esfuerzo técnico y cambios tardíos cerca del cierre. \\
Recursos & Reduce carga sobre Carlos Sainz Arce y Mauricio González Prendas en el dashboard ejecutivo, permitiendo priorizar integración y QA. \\
\bottomrule
\end{tabularx}
}
\end{table}

\subsubsection{Decisión sobre SC-02}

La solicitud SC-02 se aprueba debido a que compensa de forma razonable el incremento de alcance producido por SC-01. Esta decisión permite mantener el cronograma dentro del periodo planificado de tres meses y protege el presupuesto total de ₡17,625,200.

El cambio no elimina por completo el Dashboard Ejecutivo del sistema, sino que reduce su profundidad. Por lo tanto, el módulo sigue cumpliendo su objetivo visual básico, pero con menor detalle histórico y financiero.

\subsection{Impacto neto de los cambios}

El impacto conjunto de SC-01 y SC-02 se considera controlado. La primera solicitud aumenta el valor del Portal Estudiantil, mientras que la segunda reduce el esfuerzo en el Dashboard Ejecutivo. En consecuencia, el balance permite conservar el alcance general del prototipo y mantener la coherencia con el cronograma, el presupuesto y los criterios de calidad.

\begin{table}[H]
\centering
\caption{Impacto neto de los cambios evaluados}
\label{tab:impacto-neto}
{\small
\begin{tabularx}{\textwidth}{p{3.5cm} Y}
\toprule
\textbf{Área} & \textbf{Resultado neto} \\
\midrule
Alcance & Se agrega una sección de certificados académicos y se simplifica el Dashboard Ejecutivo del sistema. \\
Cronograma & El esfuerzo adicional se compensa, por lo que la fecha final del proyecto se mantiene dentro de los tres meses planificados. \\
Costo & El presupuesto total se mantiene en ₡17,625,200 sin incremento adicional. \\
Calidad & Se priorizan módulos de mayor valor para usuarios finales y se reducen elementos visuales complejos no esenciales. \\
Riesgos & Se controla el riesgo de alcance creciente mediante una reducción compensatoria y seguimiento formal. \\
Recursos & Se mantiene la asignación general del equipo sin agregar nuevos integrantes ni contratar recursos externos. \\
\bottomrule
\end{tabularx}
}
\end{table}

\subsection{Actualizaciones requeridas}

Después de aprobar las solicitudes se deben actualizar los documentos y evidencias afectadas. En particular, el alcance, el cronograma, el plan de calidad, el registro de riesgos y los documentos de cierre deben reflejar que el módulo de certificados se incorpora de manera representativa y que el Dashboard Ejecutivo del sistema queda simplificado.

\begin{itemize}[leftmargin=*, label={\textbullet}]
    \item Ajustar el alcance del Portal Estudiantil para incluir la sección de certificados académicos.
    \item Ajustar el alcance del Dashboard Ejecutivo para mantener solo indicadores generales.
    \item Verificar que el cronograma mantenga la fecha final dentro del periodo planificado.
    \item Revisar riesgos R-01, R-04 y R-10 para confirmar que el cambio quedó controlado.
    \item Mantener el presupuesto total en ₡17,625,200.
    \item Documentar evidencia de aprobación y comunicar el resultado al equipo.
\end{itemize}

\subsection{Seguimiento posterior}

El seguimiento de los cambios aprobados será responsabilidad de la directora del proyecto, con apoyo de los analistas documentadores y los desarrolladores Front-End. Durante las revisiones finales se deberá verificar que el módulo agregado no genere errores críticos de navegación y que la reducción del Dashboard Ejecutivo no contradiga los criterios de aceptación del prototipo.

La comunicación de la decisión se realizará al equipo interno mediante el canal de coordinación definido y se dejará evidencia en los documentos finales. Si durante las pruebas se detecta que el cambio genera nuevos problemas, estos deberán registrarse como defectos o riesgos activos según corresponda.


\newpage
% --- Entregable 24 ---
\section{Informe Ejecutivo}

\subsection{Información del proyecto}

\begin{table}[H]
\centering
\begin{tabularx}{\textwidth}{p{5cm} Y}
\toprule
\textbf{Nombre del proyecto} & Plataforma Universitaria Integrada \\
\textbf{Organización} & Universidad Tecnológica La Mejor \\
\textbf{Dirigido a} & Dr. Roberto Mora Fallas,Sponsor / Rector \\
\textbf{Directora del proyecto} & Tamara Robles Camacho \\
\textbf{Período del proyecto} & 8 de junio – 12 de agosto, 2026 (33 días hábiles / 3 meses) \\
\textbf{Presupuesto total aprobado} & ₡17,625,200 \\
\textbf{Versión} & 1.0  |  Junio 2026 \\
\bottomrule
\end{tabularx}
\end{table}

\subsection{Resumen Ejecutivo}
El proyecto Plataforma Universitaria Integrada fue desarrollado para la Universidad Tecnológica La Mejor con el objetivo de modernizar y centralizar sus procesos académicos, administrativos y financieros mediante un prototipo Front-End navegable. El proyecto se ejecutó bajo la metodología PMI a lo largo de cinco procesos (Inicio, Planificación, Ejecución, Monitoreo y Control, y Cierre), con una duración total planificada de 3 meses hábiles distribuidos entre el 8 de junio y el 9 de septiembre de 2026.\\
El equipo del proyecto estuvo conformado por cinco integrantes: una Project Manager, dos analistas de negocio/documentadores y dos desarrolladores Front-End, quienes produjeron 25 entregables documentales bajo el marco del PMBOK y un prototipo funcional que cubre los módulos de Portal Estudiantil, Portal Docente, Portal Administrativo, Portal Financiero y Dashboard Ejecutivo.\\
Durante la ejecución se gestionaron dos solicitudes de cambio formales, las cuales fueron evaluadas, aprobadas y ejecutadas internamente sin afectar el presupuesto total ni la fecha de cierre planificada al inicio del proyecto. El proyecto se considera exitoso en términos de alcance, calidad y control, cumpliendo con los criterios de éxito definidos en el Project Charter.

\subsection{Objetivos del Proyecto vs. Resultados Obtenidos}
En la siguiente tabla se compara los objetivos del proyecto contra los resultados obtenidos al cierre del proyecto:
\begin{table}[H]
\centering
\caption{Objetivos del proyecto y resultados obtenidos}
\label{tab:objetivos-resultados}
{\small
\begin{tabularx}{\textwidth}{p{6cm}Y}
\toprule
\textbf{Objetivo planteado} &
\textbf{Resultado obtenido} \\
\midrule

Desarrollar un prototipo Front-End web navegable que represente los módulos principales de la Universidad Tecnológica La Mejor.
&
El prototipo es completamente navegable y se implementaron los cinco módulos principales.
\\

Centralizar visualmente los procesos académicos, administrativos y financieros dentro de una misma plataforma.
&
Los procesos fueron centralizados mediante una navegación integrada entre todos los módulos.
\\

Diseñar un Portal Estudiantil con inicio de sesión, perfil, matrícula, horarios, calificaciones, estado de cuenta e historial académico.
&
El Portal Estudiantil fue completado al 100\% y aprobado mediante pruebas de calidad y usabilidad.
\\

Diseñar un Portal Docente con gestión de cursos, registro de asistencia y calificaciones.
&
El Portal Docente fue completado al 100\% y validado mediante pruebas de calidad.
\\

Diseñar un Portal Administrativo con gestión de estudiantes, docentes, cursos y períodos académicos.
&
El Portal Administrativo fue completado al 100\% y superó las pruebas definidas.
\\

Diseñar un Portal Financiero con pagos de matrícula, pagos de cursos e historial de pagos.
&
El Portal Financiero fue desarrollado completamente y aprobado durante las pruebas.
\\

Diseñar un Dashboard Ejecutivo con indicadores académicos y financieros.
&
El Dashboard fue implementado con los indicadores definidos y aprobados.
\\

Aplicar los grupos de procesos del PMI durante el desarrollo del proyecto.
&
Los grupos de Inicio, Planificación, Ejecución, Monitoreo y Control y Cierre fueron completados exitosamente.
\\

Elaborar los entregables documentales requeridos por el proyecto.
&
Se completaron los 26 entregables documentales establecidos.
\\

Mantener el proyecto dentro del presupuesto aprobado de ₡17,625,200.
&
El presupuesto se mantuvo sin incremento mediante una reasignación interna de esfuerzo.
\\

Completar el proyecto dentro del plazo establecido.
&
El cronograma se cumplió dentro del periodo planificado.
\\

Gestionar formalmente las solicitudes de cambio.
&
Las solicitudes SC-01 y SC-02 fueron evaluadas, aprobadas y documentadas mediante el proceso de control de cambios.
\\

\bottomrule
\end{tabularx}
}
\end{table}

\subsection{Desempeño en Tiempo y Costo}

En la siguiente sección se evalúa el desarrollo del proyecto en términos de cronograma y costo, considerando hitos, fechas de cumplimiento y el estado de avance. Esto nos permite identificar posibles desviaciones y valorar la efectividad de la gestión implementada durante las diferentes fases del proyecto.


\subsubsection{Cronograma}
El proyecto se planificó con una duración total de 3 meses, distribuidos en las fases de Inicio (33 días), Planificación (20 días), Supervisión y control, Ejecución/Desarrollo Front-End (38 días en paralelo con Planificación) y Cierre (27 días), con fecha de inicio el 8 de junio y fecha de cierre el 9 de septiembre de 2026.\\
El módulo de Project Charter aprobado se cumplió el 12 de junio, el de Alcance definido el 19 de junio, el de Cronograma aprobado el 3 de julio, el de Registro de Riesgos completado el 1 de julio y el de Prototipo Front-End completado el 29 de julio. El hito final de Entrega del proyecto se programó para el 9 de septiembre de 2026.


\begin{table}[H]
\centering
\caption{Cumplimiento de hitos del proyecto}
\label{tab:hitos}
{\small
\begin{tabularx}{\textwidth}{p{6cm}p{3cm}Y}
\toprule
\textbf{Hito} &
\textbf{Fecha planificada} &
\textbf{Estado} \\
\midrule

Project Charter aprobado &
12 jun 2026 &
Cumplido \\

Alcance definido &
19 jun 2026 &
Cumplido \\

Cronograma aprobado &
3 jul 2026 &
Cumplido \\

Registro de riesgos completado &
1 jul 2026 &
Cumplido \\

Prototipo Front-End completado &
29 jul 2026 &
Cumplido \\

Cambios evaluados &
8 jun 2026 &
Cumplido \\

Entrega final del proyecto &
9 sep 2026 &
Cumplido dentro del plazo \\

\bottomrule
\end{tabularx}
}
\end{table}

\subsubsection{Costo}

El presupuesto total aprobado para el proyecto fue de ₡17,625,200, financiado mediante la estimación de mano de obra del equipo de cinco integrantes, software, hardware, servicios operativos y reserva de contingencia. Las dos solicitudes de cambio gestionadas durante la ejecución no generaron incremento en el presupuesto total, ya que el esfuerzo adicional del módulo de certificados académicos se compensó mediante el ajuste del alcance del Dashboard Ejecutivo del sistema.

\begin{table}[H]
\centering
\caption{Desempeño del costo del proyecto}
\label{tab:costo-proyecto}
\begin{tabularx}{\textwidth}{p{8cm}Y}
\toprule
\textbf{Concepto} & \textbf{Valor} \\
\midrule
Presupuesto total aprobado (PV) & ₡17,625,200 \\
Valor ganado al cierre (EV) & ₡17,625,200 \\
Costo real al cierre (AC) & ₡17,279,608 \\
Variación del costo (CV = EV - AC) & +₡345,592 \\
\bottomrule
\end{tabularx}
\end{table}

El equipo logró completar el 100\% del alcance planificado con un gasto real (AC) de ₡17,279,608, lo cual representa un ahorro de ₡345,592 frente al presupuesto aprobado. Este resultado se atribuye al uso de herramientas gratuitas o de licencia comercial y a que las solicitudes de cambio SC-01 y SC-02 se gestionaron mediante reasignación interna de esfuerzo en lugar de recurrir a la reserva de contingencia.

\subsection{Gestión de Solicitudes de Cambio}

Durante la ejecución del proyecto se documentaron y gestionaron dos solicitudes de cambio formales, siguiendo el proceso de control de cambios establecido. Ambas solicitudes fueron evaluadas en conjunto, dado que la segunda funcionó como medida compensatoria de la primera.

\begin{table}[H]
\centering
\caption{Solicitudes de cambio del proyecto}
\label{tab:cambios}
{\small
\begin{tabularx}{\textwidth}{p{1.5cm}p{3cm}Yp{1.7cm}p{2cm}}
\toprule
\textbf{Código} &
\textbf{Solicitante} &
\textbf{Descripción} &
\textbf{Prioridad} &
\textbf{Decisión} \\
\midrule

SC-01 &
Lic. Patricia Vargas Soto y Prof. Andrea Solís Zamora &
Incorporar un módulo de solicitud de certificados académicos dentro del Portal Estudiantil.
&
Alta &
Aprobada con condiciones
\\

SC-02 &
Tamara Robles Camacho &
Reducir el alcance del Dashboard Ejecutivo como medida compensatoria.
&
Media &
Aprobada
\\

\bottomrule
\end{tabularx}
}
\end{table}

La solicitud SC-01 se aprobó con la condición de limitar su implementación a una representación visual navegable con datos simulados, sin conexión a base de datos ni generación real de documentos oficiales. La solicitud SC-02 permitió liberar el esfuerzo necesario para absorber el cambio anterior sin afectar el cronograma ni el presupuesto, simplificando el Dashboard Ejecutivo a únicamente los indicadores de estudiantes activos e ingresos por matrícula.\\
El impacto neto de ambos cambios se considera controlado: se incrementó el valor del Portal Estudiantil para los usuarios finales, se mantuvo el presupuesto total sin alteraciones, se conservó la fecha de cierre planificada y se redujeron los riesgos R-04 (subestimación del esfuerzo de desarrollo) y R-10 (cambios tardíos cerca del cierre).


\subsection{Riesgos Relevantes Durante la Ejecución}

El proyecto identificó 16 riesgos en su fase de planificación, de los cuales 5 fueron clasificados como Críticos: restricción fuerte de tiempo (R-01), alcance creciente o scope creep (R-02), inconsistencias entre documentos (R-03), subestimación del esfuerzo del Front-End (R-04) y disponibilidad limitada del equipo (R-05).\\
Durante la ejecución, los riesgos R-04 y R-10 se activaron parcialmente como consecuencia de la solicitud de cambio SC-01, pero fueron gestionados de forma efectiva mediante la aprobación compensatoria de SC-02, lo cual redujo su nivel de criticidad antes de que afectaran la entrega final del proyecto.


\begin{table}[H]
\centering
\caption{Riesgos relevantes al cierre del proyecto}
\label{tab:riesgos}
\begin{tabularx}{\textwidth}{p{1.5cm}Yp{4cm}}
\toprule
\textbf{Código} &
\textbf{Riesgo} &
\textbf{Estado al cierre} \\
\midrule

R-01 &
Restricción fuerte de tiempo del proyecto &
Gestionado; cronograma mantenido
\\

R-02 &
Alcance creciente (scope creep) &
Controlado mediante SC-02
\\

R-04 &
Subestimación del esfuerzo del Front-End &
Mitigado mediante reasignación interna
\\

R-10 &
Cambios tardíos cerca de la entrega &
Reducido tras aprobación de SC-02
\\

\bottomrule
\end{tabularx}
\end{table}

\subsection{Indicadores de Desempeño (KPIs)}

Al cierre del proyecto se midieron los 14 indicadores definidos en el Plan de KPIs, distribuidos en las categorías de tiempo, costo, calidad, riesgos y productividad. Los resultados finales confirman que el proyecto cumplió o superó la totalidad de las metas establecidas.

\subsubsection{Indicadores de Tiempo}

La siguiente tabla resume el comportamiento de los indicadores asociados al
cronograma.

\begin{table}[H]
\centering
\caption{Indicadores de tiempo}
\label{tab:kpi-tiempo}
\begin{tabularx}{\textwidth}{Yp{3cm}p{2.5cm}p{2.5cm}}
\toprule
\textbf{KPI} &
\textbf{Valor final} &
\textbf{Meta} &
\textbf{Estado} \\
\midrule

Cumplimiento del cronograma &
95\% &
$\geq$ 90\% &
En meta
\\

Eficiencia de cronograma (SPI) &
1.00 &
$\geq$ 0.95 &
En meta
\\

Variación del cronograma (SV) &
0 días &
$\geq$ 0 &
En meta
\\

\bottomrule
\end{tabularx}
\end{table}

\subsubsection{Indicadores de Costo}

La siguiente tabla resume el comportamiento de los indicadores asociados al
desempeño financiero del proyecto.

\begin{table}[H]
\centering
\caption{Indicadores de costo}
\label{tab:kpi-costo}
\begin{tabularx}{\textwidth}{Yp{3cm}p{2.5cm}p{2.5cm}}
\toprule
\textbf{KPI} &
\textbf{Valor final} &
\textbf{Meta} &
\textbf{Estado} \\
\midrule

Cumplimiento del presupuesto &
92\% ejecutado &
$\leq$ 100\% &
En meta
\\

Eficiencia de costos (CPI) &
1.02 &
$\geq$ 1.0 &
En meta
\\

Variación del costo (CV) &
+₡345,592 &
$\geq$ 0 &
En meta
\\

\bottomrule
\end{tabularx}
\end{table}

\subsubsection{Indicadores de Calidad}

La calidad del proyecto fue evaluada mediante indicadores relacionados con
errores, pruebas y satisfacción del equipo.

\begin{table}[H]
\centering
\caption{Indicadores de calidad}
\label{tab:kpi-calidad}
\begin{tabularx}{\textwidth}{Yp{3cm}p{2.5cm}p{2.5cm}}
\toprule
\textbf{KPI} &
\textbf{Valor final} &
\textbf{Meta} &
\textbf{Estado} \\
\midrule

Defectos por módulo &
0 errores reportados &
$\leq$ 3 &
En meta
\\

Cobertura de criterios de aceptación &
100\% cumplidos &
$\geq$ 95\% &
En meta
\\

Pruebas exitosas (aprobación) &
100\% &
$\geq$ 90\% &
En meta
\\

Satisfacción del equipo (escala 1-5) &
4.5 &
$\geq$ 4.0 &
En meta
\\

\bottomrule
\end{tabularx}
\end{table}

\subsubsection{Indicadores de Riesgos}

La siguiente tabla presenta el comportamiento de los principales indicadores
asociados a la gestión de riesgos.

\begin{table}[H]
\centering
\caption{Indicadores de riesgos}
\label{tab:kpi-riesgos}
\begin{tabularx}{\textwidth}{Yp{3cm}p{2.5cm}p{2.5cm}}
\toprule
\textbf{KPI} &
\textbf{Valor final} &
\textbf{Meta} &
\textbf{Estado} \\
\midrule

Riesgos críticos abiertos al cierre &
0 &
0 al cierre &
En meta
\\

Tiempo promedio de respuesta ante riesgos &
2 días &
$\leq$ 3 días &
En meta
\\

Riesgos mitigados sobre riesgos identificados &
85\% &
$\geq$ 80\% &
En meta
\\

\bottomrule
\end{tabularx}
\end{table}

\subsubsection{Indicadores de Productividad}

El desempeño del equipo también fue medido mediante indicadores de
productividad asociados al ritmo de entrega de productos y documentos.

\begin{table}[H]
\centering
\caption{Indicadores de productividad}
\label{tab:kpi-productividad}
\begin{tabularx}{\textwidth}{Yp{3cm}p{2.5cm}p{2.5cm}}
\toprule
\textbf{KPI} &
\textbf{Valor final} &
\textbf{Meta} &
\textbf{Estado} \\
\midrule

Entregables completados por semana &
9.6 / semana &
$\geq$ 2 / semana &
En meta
\\

\bottomrule
\end{tabularx}
\end{table}

Los catorce indicadores definidos en el Plan de KPIs finalizaron dentro o por
encima de las metas establecidas. Destacan los valores del SPI y del CPI,
los cuales demuestran que el proyecto se ejecutó eficientemente tanto en
tiempo como en costo.

\subsection{Conclusiones y Recomendaciones}

El proyecto \textit{Plataforma Universitaria Integrada} se considera exitoso al cierre de su ejecución. Se cumplieron los objetivos de alcance establecidos en el Project Charter, el presupuesto se mantuvo sin incremento en ₡17,625,200, el cronograma se completó dentro del periodo de tres meses planificado y las dos solicitudes de cambio surgidas durante la ejecución se implementaron de forma controlada sin comprometer la calidad del prototipo ni la fecha de entrega.\\
El prototipo Front-End navegable entregado cubre los cinco módulos requeridos por la Universidad Tecnológica La Mejor (Portal Estudiantil, Portal Docente, Portal Administrativo, Portal Financiero y Dashboard Ejecutivo), incluyendo la funcionalidad adicional de solicitud de certificados académicos aprobada durante la ejecución, lo cual aumenta el valor percibido por los estudiantes sin afectar el proyecto.








\newpage
% --- Entregable 25 ---
\section{Acta de cierre}

\subsection{Información del proyecto}

\begin{table}[H]
\centering
\begin{tabularx}{\textwidth}{p{5cm} Y}
\toprule
\textbf{Nombre del proyecto} & Plataforma Universitaria Integrada \\
\textbf{Patrocinador} & Dr. Roberto Mora Fallas \\
\textbf{Directora del proyecto} & Tamara Robles Camacho \\
\textbf{Fecha de inicio} & 8 de junio de 2026 \\
\textbf{Fecha de cierre} & 23 de junio de 2026 \\
\textbf{Presupuesto asignado} & ₡17,625,200 \\
\bottomrule
\end{tabularx}
\end{table}

\subsection{Resumen de ejecución del proyecto}

La ejecución del proyecto logró cumplir con los objetivos principales establecidos en el inicio debido a que el equipo completó el desarrollo del prototipo Front-End navegable y generó toda la documentación de gestión requerida por el estándar evaluado. De esta manera se entregó una solución visual que incluye el inicio de sesión y los portales estudiantil docente administrativo y financiero con los ajustes realizados durante el control de cambios. Por lo tanto el desempeño general del equipo permitió abarcar las expectativas de los stakeholders ya que se aplicaron adecuadamente los procesos de planificación y ejecución en un margen de tiempo sumamente reducido.

\begin{longtable}{L{2cm} L{9cm} L{3cm}}
\caption{Estado de los entregables del proyecto}\label{tab:estado-entregables}\\[-0.5em]
\toprule
\textbf{Código} & \textbf{Nombre del entregable} & \textbf{Estado} \\
\midrule
\endfirsthead
\toprule
\textbf{Código} & \textbf{Nombre del entregable} & \textbf{Estado} \\
\midrule
\endhead
E-01 & Business Case & Completado \\
E-02 & Project Charter & Completado \\
E-03 & Registro de Stakeholders & Completado \\
E-04 & Definición del Alcance & Completado \\
E-05 & WBS del proyecto & Completado \\
E-06 & Cronograma en Microsoft Project & Completado \\
E-07 & Matriz RACI & Completado \\
E-08 & Organigrama del Proyecto & Completado \\
E-09 & Plan de Recursos & Completado \\
E-10 & Estimación de Costos & Completado \\
E-11 & Plan de Calidad & Completado \\
E-12 & Plan de Comunicaciones & Completado \\
E-13 & Matriz de probabilidad e impacto & Completado \\
E-14 & Registro de Riesgos & Completado \\
E-15 & Plan de Adquisiciones & Completado \\
E-16 & KPIs del proyecto & Completado \\
E-17 & Proceso de control de cambios & Completado \\
E-18 & Prototipo Front-End navegable & Completado \\
E-19 & Plantilla de Solicitud de Cambio & Completado \\
E-20 & Solicitud de Cambio 1 & Completado \\
E-21 & Solicitud de Cambio 2 & Completado \\
E-22 & Evaluación de Cambios & Completado \\
E-23 & Dashboard ejecutivo del proyecto & Completado \\
E-24 & Informe ejecutivo & Completado \\
E-25 & Acta de Cierre & Completado \\
E-26 & Reflexiones y Recomendaciones & Completado \\
\bottomrule
\end{longtable}

\subsection{Verificación del alcance}

El equipo de proyecto confirma que el alcance aprobado y documentado fue entregado en su totalidad mediante el prototipo web funcional y los artefactos de administración de proyectos. Asimismo los elementos que se definieron como fuera del alcance tales como el backend funcional y la base de datos real se mantuvieron efectivamente excluidos durante toda la ejecución técnica. Esto implica que la navegación y las pantallas cumplen con los criterios de aceptación acordados para lo cual se realizaron las pruebas de calidad correspondientes antes de formalizar la entrega de los productos finales al patrocinador.

\subsection{Desviaciones y cambios aprobados durante el proyecto}

Durante el transcurso del trabajo se presentaron desviaciones controladas que fueron procesadas y aprobadas mediante el sistema formal de gestión de cambios institucionales. Por lo tanto se aprobó la Solicitud de Cambio 1 para incluir el módulo de solicitud de certificados en el portal estudiantil ya que esto representaba un valor agregado importante para los usuarios y la administración. En consecuencia para balancear el cronograma y mantener el control del esfuerzo se aprobó la Solicitud de Cambio 2 que redujo el alcance de los indicadores históricos en el dashboard ejecutivo. De esta manera el impacto neto sobre el proyecto resultó equilibrado debido a que el incremento en el trabajo de desarrollo del portal estudiantil se compensó con la disminución de funciones en el área ejecutiva manteniendo intacto el presupuesto global y cumpliendo la fecha final establecida.

\subsection{Presupuesto ejecutado versus planificado}

\begin{table}[H]
\centering
\begin{tabularx}{\textwidth}{p{5cm} Y Y Y}
\toprule
\textbf{Categoría} & \textbf{Planificado caso de negocio} & \textbf{Planificado PERT} & \textbf{Ejecutado real estimado} \\
\midrule
Mano de obra y desarrollo & ₡14,000,000 & ₡8,500,000 & ₡9,200,000 \\
Licencias e infraestructura & ₡3,600,000 & ₡2,100,000 & ₡2,100,000 \\
Pruebas y calidad & ₡3,000,000 & ₡2,500,000 & ₡2,800,000 \\
Gestión y documentación & ₡5,000,500 & ₡3,525,200 & ₡3,525,200 \\
Contingencias & ₡1,600,000 & ₡1,000,000 & ₡0 \\
\midrule
\textbf{Total general} & \textbf{₡17,625,200} & \textbf{₡17,625,200} & \textbf{₡17,625,200} \\
\bottomrule
\end{tabularx}
\end{table}

El presupuesto ejecutado real estimado refleja los ajustes realizados tras la aprobación de los cambios en el proyecto ya que se reasignaron los esfuerzos de desarrollo sin incrementar el costo total y se absorbió el margen operativo planificado de manera exitosa. 

\subsection{Lecciones aprendidas}

La principal lección del proyecto indica que definir un alcance claro desde el primer día resulta fundamental para manejar las expectativas de todos los involucrados debido a que un límite bien documentado evita que se soliciten funciones complejas que no pueden ser atendidas en el corto plazo. Asimismo la gestión de riesgos demostró que mantener un monitoreo constante sobre la disponibilidad del equipo y sobre las subestimaciones del esfuerzo de desarrollo permite aplicar medidas correctivas a tiempo ya que las acciones de mitigación salvaron la entrega final al reducir la carga del dashboard ejecutivo. Además de esto la documentación compartida y estandarizada es vital para conservar la coherencia entre los múltiples archivos generados por lo tanto el uso de herramientas centralizadas y revisiones cruzadas disminuye significativamente las contradicciones en fechas y presupuestos.

Por otro lado el trabajo técnico requiere de revisiones de usabilidad continuas para garantizar que el producto final sea valioso para el estudiante ya que un buen diseño de interfaz compensa la ausencia de funcionalidades avanzadas como la conexión a bases de datos en los ambientes controlados. Finalmente la comunicación abierta y constante mediante plataformas ágiles como WhatsApp y GitHub se consolidó como un factor determinante para el éxito del proyecto debido a que la rápida resolución de dudas evitó bloqueos técnicos y permitió avanzar con los entregables documentales sin perder el ritmo de trabajo exigido por el cronograma inicial.

\subsection{Formalización del cierre}

\begin{table}[H]
\centering
\begin{tabularx}{\textwidth}{Y Y Y}
\toprule
\textbf{Directora del proyecto} & \textbf{Analista de negocio} & \textbf{Desarrollador de la interfaz} \\
\midrule
& & \\
& & \\
Tamara Robles Camacho & Isaac Jiménez Blanco & Mauricio González Prendas \\
\midrule
& & \\
& & \\
Dr. Roberto Mora Fallas & Jorge Abarca Quiros & Carlos Sainz Arce \\
\textbf{Patrocinador} & \textbf{Analista de negocio} & \textbf{Desarrollador de la interfaz} \\
\bottomrule
\end{tabularx}
\end{table}

\newpage
% --- Entregable 26 ---
\section{Reflexiones y recomendaciones}

\subsection{Reflexión sobre la gestión del proyecto}

La gestión del proyecto se caracterizó por una aplicación rigurosa de los grupos de procesos propuestos por el estándar en medio de un marco de tiempo sumamente desafiante debido a que el trabajo inició el 8 de junio y debía entregarse el 24 de junio de 2026. Durante la fase de inicio y planificación se logró establecer un alcance realista y una estructura de trabajo bien delimitada ya que los documentos como el acta de constitución y el cronograma brindaron el horizonte necesario para ordenar las tareas de todo el equipo de desarrollo. Asimismo la ejecución y el monitoreo permitieron identificar de manera oportuna las desviaciones en el esfuerzo de desarrollo del Front-End por lo tanto la implementación del control de cambios funcionó como una herramienta estratégica para proteger la calidad del producto y asegurar que el proyecto finalizara exitosamente en la fecha estipulada por la institución.

\subsection{Reflexión sobre el trabajo en equipo}

El desempeño colaborativo del grupo compuesto por cinco miembros resultó ser uno de los pilares más fuertes de la iniciativa ya que la división clara de los roles entre documentadores y desarrolladores permitió avanzar en paralelo sin generar cuellos de botella importantes. De esta manera el uso de herramientas de comunicación inmediata como WhatsApp agilizó las consultas cotidianas y los acuerdos rápidos mientras que el uso de GitHub garantizó un control efectivo de las versiones del código del prototipo web en cada fase. En consecuencia las reuniones de seguimiento regulares facilitaron la sincronización entre el diseño técnico y los artefactos de gestión para lo cual el liderazgo de la administración del proyecto logró integrar las diferentes perspectivas y mantener la motivación a pesar de la presión fuerte del cronograma semanal.

\subsection{Reflexión sobre la gestión de riesgos}

El análisis de los 16 riesgos identificados en la planificación demostró que el equipo anticipó correctamente los principales retos que afectaban el proyecto ya que la restricción de tiempo y la complejidad técnica figuraron como las amenazas de mayor impacto desde el primer momento de trabajo. Por lo tanto algunos riesgos críticos como el incremento desmedido del alcance y la subestimación del esfuerzo de desarrollo se materializaron parcialmente cuando se solicitaron funcionalidades adicionales en el portal estudiantil durante el ciclo iterativo. Esto implica que las estrategias de mitigación establecidas fueron completamente efectivas debido a que la reducción del alcance en el dashboard ejecutivo permitió controlar el riesgo R-04 y el riesgo R-10 garantizando así que el producto conservara su estabilidad general y se entregara sin defectos mayores a los usuarios finales.

\subsection{Reflexión sobre la calidad del proyecto}

El cumplimiento de los estándares de calidad se evidenció a través de las diversas revisiones aplicadas tanto a la documentación como al prototipo final debido a que el grupo se apegó a los criterios de aceptación definidos para cada módulo de la plataforma universitaria. Asimismo las pruebas manuales ejecutadas sobre el Front-End aseguraron que la navegación y el diseño visual ofrecieran una experiencia intuitiva para los estudiantes docentes y administrativos ya que se priorizaron los elementos funcionales sobre los adornos innecesarios en la interfaz gráfica. Además de esto la revisión cruzada de los archivos escritos garantizó una alta coherencia entre los presupuestos los cronogramas y los planes de gestión por lo tanto el resultado refleja un esfuerzo sistemático por entregar un trabajo académico profesional confiable y sumamente ordenado.

\subsection{Recomendaciones para proyectos futuros}

A partir de las lecciones asimiladas durante la ejecución se recomienda definir el alcance del producto con el mayor nivel de detalle posible en las primeras reuniones de inicio ya que esto ayuda a prevenir discusiones complejas cuando el tiempo de trabajo se encuentra avanzado. Asimismo resulta fundamental asignar un margen de contingencia amplio en el cronograma para todas las tareas de programación y revisión técnica debido a que el desarrollo de interfaces suele enfrentar imprevistos que consumen más horas de las planificadas originalmente en las estimaciones. Además de estas medidas se sugiere mantener una plantilla unificada para todos los documentos del proyecto desde el día uno por lo tanto el equipo evita retrabajos de formato continuos y asegura que la imagen del trabajo sea consistente ante los evaluadores académicos.

Por otro lado es recomendable establecer canales formales de aprobación de cambios que funcionen de manera ágil y sin excesiva burocracia ya que en proyectos cortos las decisiones deben tomarse rápidamente para no paralizar el avance de los programadores front-end. Finalmente se aconseja fomentar una cultura de revisión por pares en todos los entregables técnicos y documentales para lo cual cada miembro debe comprometerse a leer el trabajo de sus compañeros antes de consolidar la versión definitiva mitigando así los errores por omisión y elevando sustancialmente el estándar general del proyecto universitario.

\subsection{Conclusión general del proyecto}

El desarrollo de la plataforma universitaria integrada representa un logro valioso que combina la disciplina de la gestión de proyectos con las habilidades técnicas de la construcción de sistemas web debido a que el equipo consiguió transformar una necesidad institucional abstracta en un prototipo funcional. De esta manera el esfuerzo invertido en cada fase demuestra que la planificación adecuada y el control de los recursos son herramientas indispensables para superar las barreras del tiempo estricto y las limitaciones tecnológicas en cualquier entorno de trabajo colaborativo. En consecuencia la experiencia adquirida durante este corto período fortalece las competencias profesionales de todos los participantes dedicados y establece una base sólida para enfrentar futuros desafíos en el desarrollo de software y la administración experta de requerimientos complejos.

\end{document}
